\clearpage
\thispagestyle{plain}
\oldpage[IV-3]{5}
\phantomsection
\label{chapter:IV3-fr}

\begin{center}
{\large CHAPITRE IV (suite)\par}
\vspace{1.5em}
{\Large\bfseries ÉTUDE LOCALE DES SCHÉMAS\\
ET DES MORPHISMES DE SCHÉMAS\par}
\end{center}

\setcounter{section}{7}
\section{Limites projectives de préschémas}
\label{section:IV.8-fr}

\setcounter{subsection}{0}
\subsection{Introduction}
\label{subsection:IV.8.1-fr}

\begin{env}[8.1.1]
\label{IV.8.1.1-fr}
Dans ce paragraphe, nous allons étudier systématiquement la situation
suivante. Soient $I$ un ensemble préordonné filtrant croissant,
$(A_\lambda,\varphi_{\mu\lambda})$ un système inductif d'anneaux ayant $I$
pour ensemble d'indices, $A=\varinjlim A_\lambda$ sa limite inductive. Pour
tout $\alpha\in I$ et tout $A_\alpha$-préschéma $X_\alpha$, considérons les
$A_\lambda$-préschémas
$X_\lambda=X_\alpha\otimes_{A_\alpha}A_\lambda$ pour $\lambda\geq\alpha$,
et le $A$-préschéma $X=X_\alpha\otimes_{A_\alpha}A$~; il est clair que les
préschémas $X_\lambda$ (pour $\lambda\geq\alpha$) forment un \emph{système
projectif}, et on verra (8.2.5) que $X$ est une \emph{limite projective} de ce
système dans la catégorie des préschémas. Nous nous proposons de trouver des
conditions sur $X_\alpha$ ou sur les $A_\lambda$ permettant de démontrer des
propriétés du type suivant~: pour que $X$ possède une propriété $P$ (par
exemple, la propriété d'être propre sur $S=\operatorname{Spec}(A)$, ou
irréductible, ou connexe, etc.), il faut et il suffit qu'il existe un indice
$\mu\geq\alpha$ tel que, pour tout $\lambda\geq\mu$, $X_\lambda$ ait (par
rapport à $S_\lambda=\operatorname{Spec}(A_\lambda)$, le cas échéant), la
même propriété $P$. Nous obtiendrons des énoncés analogues pour des propriétés
de $\mathcal O_X$-Modules, de $A$-morphismes de $A$-préschémas, etc. Nous
montrerons également (8.9.1) que la donnée d'un $A$-préschéma de présentation
finie (1.6.1) équivaut essentiellement à la donnée d'un $A_\lambda$-préschéma
de présentation finie $X_\lambda$ pour un $\lambda$ assez grand, $X$ étant
alors \emph{isomorphe à} $X_\lambda\otimes_{A_\lambda}A$. On a aussi des
énoncés analogues pour les $\mathcal O_X$-Modules de présentation finie,
leurs homomorphismes, les $A$-morphismes de $A$-préschémas de présentation
finie, etc.
\end{env}

\begin{env}[8.1.2]
\label{IV.8.1.2-fr}
L'utilité de tels résultats apparaîtra par exemple dans les questions
suivantes~:
\begin{enumerate}
\item[\emph{a)}] Soient $Y$ un préschéma, $y$ un point de $Y$, $(U_\lambda)$
le système projectif filtrant décroissant des voisinages ouverts affines de
$y$ dans $Y$~; si $A_\lambda$ est l'anneau de $U_\lambda$, les $A_\lambda$
forment un système inductif filtrant croissant, dont la limite inductive $A$
est l'anneau local $\mathcal O_y$~; d'ailleurs, si on désigne par
$\mathfrak p_\lambda$ l'idéal premier $\mathfrak j_y$ dans l'anneau
$A_\lambda$, le système inductif $(A_\lambda)$ est cofinal à tout système
inductif $(A_\alpha)_f$, où $f$ parcourt
$A_\alpha-\mathfrak p_\alpha$ (pour un $\alpha$ \emph{fixé}), car les $D(f)$
forment un système fondamental de voisinages de $y$ dans $U_\alpha$,
\oldpage[IV-3]{6}
donc dans $Y$. Les résultats du présent paragraphe impliqueront que la
géométrie algébrique des $\mathcal O_y$-préschémas de présentation finie (et
la théorie des Modules de présentation finie sur ces préschémas) est
essentiellement équivalente à la géométrie algébrique des préschémas de
présentation finie sur des voisinages ouverts «~assez petits~» de $y$. Ainsi,
l'énoncé (8.10.5, (xii)) implique que si un morphisme $f:X\to Y$ est de
présentation finie, alors, pour que
$X\times_Y\operatorname{Spec}(\mathcal O_y)$ soit propre sur
$\operatorname{Spec}(\mathcal O_y)$, il faut et il suffit qu'il existe un
voisinage ouvert $U$ de $y$ dans $Y$ tel que $f^{-1}(U)$ soit propre sur $U$.

Un cas particulièrement important, et dans une certaine mesure classique,
est celui où $Y$ est intègre et où $y$ est son point générique, de sorte que
$\mathcal O_y$ n'est autre que \emph{le corps $R(Y)=K$ des fonctions
rationnelles sur $Y$}. Les résultats du présent paragraphe reviennent alors à
interpréter la géométrie algébrique sur $K$ en termes de la géométrie
algébrique au-dessus d'ouverts non vides «~assez petits~» de $Y$, c'est-à-dire,
intuitivement, en termes de «~familles~» d'objets géométriques indexées par les
points d'un tel ouvert. Ce point de vue a d'ailleurs été couramment utilisé
depuis longtemps, non seulement en Géométrie algébrique sur des corps
algébriquement clos, mais aussi dans l'étude arithmétique des variétés définies
sur un \emph{corps de nombres} $K$ (extension finie de $\mathbf Q$), en
considérant ce dernier comme le corps des fractions de son anneau des entiers
$A$ («~théorie de la \emph{réduction modulo} $\mathfrak p$~» ($\mathfrak p$
idéal premier de $A$)~; cf. (I, 3.7)). Les résultats des \S\S~8 et 9
fournissent donc entre autres des fondements du langage de la «~réduction
modulo $\mathfrak p$~» en arithmétique.

On notera que dans l'exemple envisagé ici, les morphismes
$S_\mu\to S_\lambda$ (pour $\lambda\leq\mu$) sont les \emph{immersions
ouvertes} canoniques $U_\mu\to U_\lambda$, et \emph{a fortiori} sont des
morphismes \emph{plats} (mais non fidèlement plats en général), ce qui
explique l'intérêt des énoncés qui font appel à une telle restriction.

\item[\emph{b)}] Supposons que les $A_\lambda$ soient des \emph{corps}, de
sorte que $A=\varinjlim A_\lambda$ est aussi un corps. Ce cas se présente
généralement lorsqu'on part de données géométriques au-dessus d'un corps
quelconque $K$, que l'on considère comme extension d'un corps $k$ (par exemple
le sous-corps premier de $K$). Il y a alors intérêt à considérer $K$ comme
limite inductive de ses sous-extensions qui sont \emph{de type fini sur $k$},
ce qui permet dans de nombreuses questions de se ramener au cas où $K$ est
une extension de type fini de $k$. Utilisant aussi la méthode esquissée dans
\emph{a)}, on peut alors se ramener généralement au cas d'un anneau de base
$A$ qui est une \emph{algèbre intègre de type fini sur $k$}.

On notera que dans cet exemple, les morphismes $S_\mu\to S_\lambda$ sont
\emph{fidèlement plats}.

\item[\emph{c)}] Supposons qu'on s'intéresse aux propriétés, locales sur
$Y$, des préschémas de présentation finie au-dessus d'un préschéma quelconque
$Y$, que l'on peut dès lors supposer affine d'anneau $A$. Il y a alors intérêt
à considérer $A$ comme limite inductive de ses sous-anneaux qui sont des
$\mathbf Z$-algèbres \emph{de type fini}, ce qui permet de ramener de
nombreuses questions au cas où $Y$ est le spectre d'une telle algèbre. C'est
là l'explicitation du «~point de vue kroneckerien~» suivant lequel la
Géométrie algébrique se réduit à la Géométrie algébrique des préschémas de
type fini sur $\mathbf Z$ (que l'on qualifie parfois de
\oldpage[IV-3]{7}
«~Géométrie algébrique absolue~»). Cet exemple nous montre en particulier que
dans la plupart des questions «~relatives~» sur un préschéma de base $Y$, on
peut se ramener au cas où $Y$ est \emph{noethérien}.

On notera que dans cet exemple, au contraire des précédents, les morphismes
$S_\mu\to S_\lambda$ n'ont en général aucune propriété particulière de
régularité.
\end{enumerate}

Par la suite, quand nous appliquerons les résultats qui vont suivre à l'une
quelconque des trois situations particulières qui viennent d'être décrites,
nous nous dispenserons de reprendre en détail le procédé qui permet de faire
ces applications, nous contentant de renvoyer à ce qui précède.
\end{env}

\begin{env}[8.1.3]
\label{IV.8.1.3-fr}
Dans l'exemple \emph{a)} de (8.1.2), nous avons vu que si $Y$ est un
préschéma intègre de point générique $y$, $f:X\to Y$ un morphisme de
présentation finie, alors, si la fibre générique $f^{-1}(y)$ est propre sur
$\boldsymbol k(y)$, il y a un voisinage ouvert $U$ de $y$ tel que $f^{-1}(U)$
soit propre sur $U$~; il en résulte \emph{a fortiori} que pour tout $s\in U$,
$f^{-1}(s)$ est propre sur $\boldsymbol k(s)$. Il arrive qu'on ait besoin
d'une réciproque, assurant que si $f^{-1}(s)$ est propre sur
$\boldsymbol k(s)$ pour «~suffisamment~» de points $s\in U$, alors
$f^{-1}(y)$ est propre sur $\boldsymbol k(y)$. Par exemple, supposons que $X$
et $Y$ soient des préschémas algébriques sur un corps $k$ algébriquement clos
(on peut prendre pour $k$ le corps $\mathbf C$ des nombres complexes, pour
fixer les idées)~; on a parfois besoin de savoir que si, pour tout $s\in Y$,
\emph{rationnel sur $k$}, la fibre $f^{-1}(s)$ est propre sur
$\boldsymbol k(s)$, alors $f^{-1}(y)$ est propre sur $\boldsymbol k(y)$, et
par suite $f^{-1}(U)$ est propre sur $U$ pour un voisinage $U$ de
$y$\footnote{On notera qu'un tel énoncé est en fin de compte purement
géométrique, en ce sens qu'il ne fait appel qu'à des points rationnels sur
$k$, et non à des points génériques~; par exemple, lorsque $k=\mathbf C$, cet
énoncé a une signification topologique évidente pour l'analyste, en
interprétant «~propre~» au sens topologique du terme, pour les espaces
sous-jacents aux espaces analytiques formés des points de $X$ et $Y$
rationnels sur $\mathbf C$.}. Or cet énoncé résultera aisément du suivant~:
l'ensemble $E$ des points $s\in Y$ tels que $f^{-1}(s)$ soit propre sur
$\boldsymbol k(s)$ est \emph{constructible} (et par conséquent identique à
$Y$ tout entier s'il contient les points fermés de $Y$, grâce au théorème des
zéros de Hilbert (10.4.8))~; cela équivaut encore à dire que si $f^{-1}(y)$
\emph{n'est pas propre} sur $\boldsymbol k(y)$, alors il existe un voisinage
ouvert $U$ de $y$ tel que $f^{-1}(s)$ ne soit pas propre sur
$\boldsymbol k(s)$ pour tout $s\in U$ (cf. (9.6.1, (iv))). Cet exemple
illustre l'intérêt qu'il y a à développer systématiquement des \emph{critères
de constructibilité} pour les notions les plus importantes~: c'est ce qui
sera fait au \S~9.
\end{env}

\subsection{Limites projectives de préschémas}
\label{subsection:IV.8.2-fr}

\begin{env}[8.2.1]
\label{IV.8.2.1-fr}
Soient $S_0$ un espace annelé, $L$ un ensemble préordonné filtrant croissant,
$(\mathscr A_\lambda,\varphi_{\mu\lambda})$ un système inductif de
$\mathcal O_{S_0}$-Algèbres (non nécessairement commutatives), ayant $L$ pour
ensemble d'indices. On sait que, considéré comme système inductif de
$\mathcal O_{S_0}$-Modules, $(\mathscr A_\lambda,\varphi_{\mu\lambda})$ admet
une limite inductive $\mathscr A$~; désignons par
$\varphi_\lambda:\mathscr A_\lambda\to\mathscr A$ l'homomorphisme canonique
(de $\mathcal O_{S_0}$-Modules). Soit
$m_\lambda:\mathscr A_\lambda\otimes\mathscr A_\lambda\to\mathscr A_\lambda$
l'homomorphisme de $\mathcal O_{S_0}$-Modules qui définit la multiplication
dans la $\mathcal O_{S_0}$-Algèbre $\mathscr A_\lambda$~; l'hypothèse sur les
$\varphi_{\mu\lambda}$ entraîne que les $m_\lambda$ forment un système
inductif d'homomorphismes,
\oldpage[IV-3]{8}
et comme le foncteur $\varinjlim$ commute au produit tensoriel,
$m=\varinjlim m_\lambda$ est un homomorphisme
$\mathscr A\otimes\mathscr A\to\mathscr A$ de
$\mathcal O_{S_0}$-Modules~; par passage à la limite sur les diagrammes
commutatifs exprimant l'associativité de $m_\lambda$ et l'existence de la
section unité dans $\mathscr A_\lambda$, on voit que $m$ définit sur
$\mathscr A$ une structure de $\mathcal O_{S_0}$-Algèbre et que
$\varphi_\lambda$ est un homomorphisme de $\mathcal O_{S_0}$-Algèbres pour
tout $\lambda\in L$. En outre $\mathscr A$ est limite inductive du système
$(\mathscr A_\lambda,\varphi_{\mu\lambda})$ \emph{dans la catégorie des
$\mathcal O_{S_0}$-Algèbres}, autrement dit, pour toute
$\mathcal O_{S_0}$-Algèbre $\mathscr B$, l'application canonique
\begin{equation}
\operatorname{Hom}_{\mathcal O_{S_0}\text{-Alg.}}(\mathscr A,\mathscr B)
\longrightarrow
\varprojlim_\lambda
\operatorname{Hom}_{\mathcal O_{S_0}\text{-Alg.}}
(\mathscr A_\lambda,\mathscr B)
\tag{8.2.1.1}\label{IV.8.2.1.1-fr}
\end{equation}
qui à tout homomorphisme $f:\mathscr A\to\mathscr B$ de
$\mathcal O_{S_0}$-Algèbres, fait correspondre la famille
$(f\circ\varphi_\lambda)$, est \emph{bijective}. En effet, on sait déjà
qu'elle est injective et identifie
$\operatorname{Hom}_{\mathcal O_{S_0}\text{-Alg.}}(\mathscr A,\mathscr B)$
à une partie de
$\varprojlim_\lambda
\operatorname{Hom}_{\mathcal O_{S_0}\text{-Mod.}}
(\mathscr A_\lambda,\mathscr B)$~; tout revient à voir que si $(f_\lambda)$
est un système inductif d'homomorphismes de $\mathcal O_{S_0}$-Algèbres,
$f_\lambda:\mathscr A_\lambda\to\mathscr B$, sa limite inductive
$f:\mathscr A\to\mathscr B$, qui est par définition un homomorphisme de
$\mathcal O_{S_0}$-\emph{Modules}, est aussi un homomorphisme de
$\mathcal O_{S_0}$-Algèbres~; mais cela résulte du passage à la limite
inductive dans le diagramme commutatif d'homomorphismes de
$\mathcal O_{S_0}$-Modules
\[
\xymatrix{
  \mathscr A_\lambda\otimes\mathscr A_\lambda
    \ar[r]^{f_\lambda\otimes f_\lambda} \ar[d]_{m_\lambda}
    & \mathscr B\otimes\mathscr B \ar[d] \\
  \mathscr A_\lambda \ar[r]_{f_\lambda} & \mathscr B
}
\]
et du fait que le foncteur $\varinjlim$ commute aux produits tensoriels.

On notera enfin que si les $\mathscr A_\lambda$ sont des
$\mathcal O_{S_0}$-Algèbres commutatives, il en est de même de $\mathscr A$.
\end{env}

\begin{env}[8.2.2]
\label{IV.8.2.2-fr}
Supposons maintenant que $S_0$ soit un \emph{préschéma}, et que les
$\mathscr A_\lambda$ soient des $\mathcal O_{S_0}$-Algèbres (commutatives)
\emph{quasi-cohérentes}~; on sait alors que
$\mathscr A=\varinjlim\mathscr A_\lambda$ est une
$\mathcal O_{S_0}$-Algèbre quasi-cohérente (I, 4.1.1). Désignons par
$S_\lambda$ (resp. $S$) le spectre de la $\mathcal O_{S_0}$-Algèbre
$\mathscr A_\lambda$ (resp. $\mathscr A$) (II, 1.3.1), et soient
$u_{\lambda\mu}:S_\mu\to S_\lambda$ (pour $\lambda\leq\mu$) et
$u_\lambda:S\to S_\lambda$ les $S_0$-morphismes correspondant aux
homomorphismes $\varphi_{\mu\lambda}$ et $\varphi_\lambda$ respectivement
(II, 1.2.7)~; il est clair que $(S_\lambda,u_{\lambda\mu})$ est un
\emph{système projectif} dans la catégorie des $S_0$-préschémas. On notera
que les $u_{\lambda\mu}$ et $u_\lambda$ sont des morphismes \emph{affines}
(II, 1.6.2), donc \emph{quasi-compacts} et \emph{séparés}.
\end{env}

\begin{proposition}[8.2.3]
\label{IV.8.2.3-fr}
Avec les notations de (8.2.2), les morphismes $u_\lambda:S\to S_\lambda$
font de $S$ une limite projective du système projectif
$(S_\lambda,u_{\lambda\mu})$ dans la catégorie des préschémas. En outre, si
$h:S_0\to T$ est un morphisme, faisant de tout $S_0$-préschéma un
$T$-préschéma, $S$ est aussi limite projective du système
$(S_\lambda,u_{\lambda\mu})$ dans la catégorie des $T$-préschémas.
\end{proposition}

Prouvons d'abord la seconde assertion de l'énoncé dans le cas $T=S_0$.
\oldpage[IV-3]{9}
Tout revient à démontrer que si $X$ est un $S_0$-préschéma quelconque,
l'application canonique
\begin{equation}
\operatorname{Hom}_{S_0}(X,S)\longrightarrow
\varprojlim_\lambda\operatorname{Hom}_{S_0}(X,S_\lambda)
\tag{8.2.3.1}\label{IV.8.2.3.1-fr}
\end{equation}
qui à tout $S_0$-morphisme $v:X\to S$ fait correspondre la famille
$(u_\lambda\circ v)$, est \emph{bijective}. Or, si $g:X\to S_0$ est le
morphisme structural et si on pose $\mathscr B=g_*(\mathcal O_X)$, qui est
une $\mathcal O_{S_0}$-Algèbre, l'application (8.2.3.1) s'identifie
canoniquement à (8.2.1.1) (II, 1.2.7), et la conclusion résulte donc de ce
qu'on a vu dans (8.2.1).

Les autres assertions de (8.2.3) sont conséquences du lemme général suivant~:

\begin{lemma}[8.2.4]
\label{IV.8.2.4-fr}
Soient $\mathbf C$ une catégorie, $T$ un objet de $\mathbf C$,
$\mathbf C_T$ la sous-catégorie des objets de $\mathbf C$ au-dessus de $T$.
Soit $(S_\lambda,u_{\lambda\mu})$ un système projectif dans $\mathbf C_T$~;
alors toute limite projective de ce système dans $\mathbf C_T$ est aussi
limite projective dans $\mathbf C$, et réciproquement.
\end{lemma}

Soit $f_\lambda:S_\lambda\to T$ le morphisme structural. Supposons que $S$
soit limite projective de $(S_\lambda,u_{\lambda\mu})$ dans $\mathbf C$, et
désignons par $u_\lambda:S\to S_\lambda$ les morphismes canoniques
correspondants. Considérons alors un système projectif de $T$-morphismes
$w_\lambda:Y\to S_\lambda$, où $Y\in\mathbf C_T$. Il existe par hypothèse un
morphisme unique (dans $\mathbf C$) $w:Y\to S$ tel que
$w_\lambda=u_\lambda\circ w$ pour tout $\lambda$. L'hypothèse que les
$u_\lambda$ sont des $T$-morphismes entraîne que les morphismes
$f_\lambda\circ u_\lambda:S\to T$ sont tous les mêmes, et ce morphisme $f$
fait donc de $S$ un $T$-objet. Si $g:Y\to T$ est le morphisme structural de
$Y$, on a alors
$f\circ w=f_\lambda\circ u_\lambda\circ w=f_\lambda\circ w_\lambda=g$ pour
tout $\lambda$, ce qui prouve que $w$ est un $T$-morphisme. Inversement,
supposons (avec les mêmes notations) que $S$ soit limite projective de
$(S_\lambda,u_{\lambda\mu})$ dans $\mathbf C_T$, et considérons maintenant un
système projectif de morphismes (de $\mathbf C$) $w_\lambda:Y\to S_\lambda$.
Les morphismes composés $f_\lambda\circ w_\lambda:Y\to T$ sont alors tous
les mêmes~: en effet, pour deux indices quelconques $\lambda$, $\mu$, il y a
un indice $\nu$ tel que $\lambda\leq\nu$ et $\mu\leq\nu$, d'où
$f_\nu=f_\lambda\circ u_{\lambda\nu}=f_\mu\circ u_{\mu\nu}$~; comme
$w_\lambda=u_{\lambda\nu}\circ w_\nu$ et
$w_\mu=u_{\mu\nu}\circ w_\nu$, on a
$f_\lambda\circ w_\lambda=f_\lambda\circ u_{\lambda\nu}\circ w_\nu
=f_\nu\circ w_\nu$, et on voit de même que
$f_\mu\circ w_\mu=f_\nu\circ w_\nu$. Si $g:Y\to T$ est l'unique morphisme
ainsi défini, $g$ fait de $Y$ un $T$-objet, et les $w_\lambda$ sont alors des
$T$-morphismes~; ils ont par suite une limite projective $w:Y\to S$ qui est
un $T$-morphisme, et \emph{a fortiori} un morphisme de $\mathbf C$~;
d'ailleurs la première partie du raisonnement montre que toute limite
projective $w'$ (dans $\mathbf C$) du système projectif $(w_\lambda)$ est
nécessairement aussi un $T$-morphisme, donc égale à $w$, ce qui achève de
prouver le lemme.

\begin{proposition}[8.2.5]
\label{IV.8.2.5-fr}
Sous les conditions de (8.2.2), soient $\alpha$ un élément de $L$,
$X_\alpha$ un $S_\alpha$-préschéma. Pour tout $\lambda\geq\alpha$, posons
$X_\lambda=X_\alpha\times_{S_\alpha}S_\lambda$, et pour
$\alpha\leq\lambda\leq\mu$, posons
$v_{\lambda\mu}=1_{X_\alpha}\times u_{\lambda\mu}$, de sorte que
$(X_\lambda,v_{\lambda\mu})$ est un système projectif de
$X_\alpha$-préschémas, dont l'ensemble d'indices est formé des
$\lambda\geq\alpha$ dans $L$. Posons de même
$X=X_\alpha\times_{S_\alpha}S$ et $v_\lambda=1_{X_\alpha}\times u_\lambda$.
Alors les $X_\alpha$-morphismes $v_\lambda:X\to X_\lambda$ font de $X$ une
limite projective du système projectif $(X_\lambda,v_{\lambda\mu})$ dans la
catégorie des $X_\alpha$-préschémas, ou dans la catégorie de tous les
préschémas.
\end{proposition}

Cela résultera encore du lemme général suivant~:

\begin{lemma}[8.2.6]
\label{IV.8.2.6-fr}
Soient $\mathbf C$ une catégorie dans laquelle les produits fibrés existent,
$q:T'\to T$ un morphisme de $\mathbf C$, $\mathbf C_T$ (resp.
$\mathbf C_{T'}$) la catégorie des objets de $\mathbf C$ au-dessus
\oldpage[IV-3]{10}
de $T$ (resp. $T'$). Soit $(S_\lambda,u_{\lambda\mu})$ un système projectif
(non nécessairement filtrant) dans $\mathbf C_T$, et posons
$S'_\lambda=S_\lambda\times_T T'$,
$u'_{\lambda\mu}=u_{\lambda\mu}\times 1_{T'}$, de sorte que
$(S'_\lambda,u'_{\lambda\mu})$ est un système projectif dans
$\mathbf C_{T'}$. Alors, si $(S,u_\lambda)$ est une limite projective de
$(S_\lambda,u_{\lambda\mu})$ dans $\mathbf C_T$,
$(S\times_T T',u_\lambda\times 1_{T'})$ est une limite projective de
$(S'_\lambda,u'_{\lambda\mu})$ dans $\mathbf C_{T'}$.
\end{lemma}

On a par hypothèse, pour tout $\lambda$, un diagramme commutatif
\[
\xymatrix{
  S' \ar[r]^{u'_\lambda} \ar[d]_p
    & S'_\lambda \ar[r]^{h'_\lambda} \ar[d]_{p_\lambda}
    & T' \ar[d]^q \\
  S \ar[r]_{u_\lambda} & S_\lambda \ar[r]_{h_\lambda} & T
}
\]
où l'on a posé $S'=S\times_T T'$, $u'_\lambda=u_\lambda\times 1_{T'}$,
$h'_\lambda=h_\lambda\times 1_{T'}$. Soient $Y$ un $T'$-objet,
$g':Y\to T'$ le morphisme correspondant, et considérons un système
projectif de $T'$-morphismes $w'_\lambda:Y\to S'_\lambda$. Alors $Y$ est un
$T$-objet pour le morphisme $g=q\circ g'$, et les
$w_\lambda=p_\lambda\circ w'_\lambda$ sont des $T$-morphismes puisque
$h_\lambda\circ w_\lambda=h_\lambda\circ p_\lambda\circ w'_\lambda
=q\circ h'_\lambda\circ w'_\lambda=q\circ g'$ par hypothèse. En outre, on
vérifie aussitôt que $(w_\lambda)$ est un système projectif. Il existe donc
par hypothèse un $T$-morphisme $w:Y\to S$ et un seul tel que
$u_\lambda\circ w=w_\lambda$ pour tout $\lambda$. Par définition du produit
fibré, il y a un unique $T'$-morphisme $w':Y\to S'$ tel que
$p\circ w'=w$. On a alors
$u_\lambda\circ p\circ w'=u_\lambda\circ w=w_\lambda
=p_\lambda\circ w'_\lambda$, ce qui s'écrit aussi
$p_\lambda\circ u'_\lambda\circ w'=p_\lambda\circ w'_\lambda$~; d'autre
part, en écrivant que $w'$ est un $T'$-morphisme, il vient
$h'_\lambda\circ u'_\lambda\circ w'=g'=h'_\lambda\circ w'_\lambda$. La
définition de $S'_\lambda$ comme produit fibré $S_\lambda\times_T T'$ donne
donc $u'_\lambda\circ w'=w'_\lambda$, et il est immédiat que $w'$ est
l'unique $T'$-morphisme vérifiant ces relations, d'où le lemme.

\begin{remark}[8.2.7]
\label{IV.8.2.7-fr}
Étant donné un espace annelé quelconque $S$, les limites inductives par
rapport à un ensemble préordonné \emph{quelconque} $L$ (non nécessairement
filtrant) existent dans la catégorie des $\mathcal O_S$-Algèbres
commutatives, puisque la limite inductive filtrante existe d'après (8.2.1) et
d'autre part, pour deux homomorphismes de $\mathcal O_S$-Algèbres
$\mathscr A\to\mathscr B$, $\mathscr A\to\mathscr C$, le produit tensoriel
$\mathscr B\otimes_{\mathscr A}\mathscr C$ est la «~somme amalgamée~»
correspondante dans cette catégorie.

Lorsque $S$ est un préschéma, on sait que le produit tensoriel
$\mathscr B\otimes_{\mathscr A}\mathscr C$ est une $\mathcal O_S$-Algèbre
quasi-cohérente lorsqu'il en est ainsi de $\mathscr A$, $\mathscr B$ et
$\mathscr C$ (I, 1.3.13)~; on en conclut que, dans la catégorie des
$\mathcal O_S$-Algèbres \emph{quasi-cohérentes}, les limites inductives pour
un ensemble préordonné quelconque d'indices existent toujours. Cela permet de
généraliser la définition d'une limite projective de préschémas et les prop.
(8.2.3) et (8.2.5) au cas où l'ensemble préordonné $L$ n'est pas
nécessairement filtrant.
\end{remark}

\begin{env}[8.2.8]
\label{IV.8.2.8-fr}
Les notations étant celles de (8.2.2), posons
$u_{\lambda\mu}=(\psi_{\lambda\mu},\theta_{\lambda\mu})$ et
$u_\lambda=(\psi_\lambda,\theta_\lambda)$~; $(S_\lambda,\psi_{\lambda\mu})$
est donc un système projectif d'espaces topologiques et $(\psi_\lambda)$ un
système projectif d'applications continues $S\to S_\lambda$, des espaces
sous-jacents respectivement aux préschémas $S$ et $S_\lambda$.
\end{env}

\begin{proposition}[8.2.9]
\label{IV.8.2.9-fr}
Avec les notations de (8.2.8), la limite projective du système projectif
$(\psi_\lambda)$ d'applications continues est un homéomorphisme de l'espace
sous-jacent à $S$ sur la limite projective du système projectif
$(S_\lambda,\psi_{\lambda\mu})$ d'espaces topologiques.
\end{proposition}

\oldpage[IV-3]{11}
Soit $T$ l'espace topologique limite du système
$(S_\lambda,\psi_{\lambda\mu})$ et posons
$\psi=\varprojlim\psi_\lambda:S\to T$. On peut se borner au cas où
$S_0=S_\alpha$ pour un $\alpha\in L$, et $\lambda\geq\alpha$.

\medskip
\noindent
(i) Montrons en premier lieu que la topologie de $S$ est image réciproque par
$\psi$ de la topologie de $T$~; autrement dit, si
$\pi_\lambda:T\to S_\lambda$ est l'application canonique, il faut montrer
que tout ouvert de $S$ est réunion d'ouverts de la forme
$\psi^{-1}(\pi_\lambda^{-1}(U_\lambda))$, où $U_\lambda$ est ouvert dans
$S_\lambda$. Or tout ouvert de $S$ est par définition réunion d'ouverts
obtenus de la façon suivante~: on considère un ouvert affine $U_0$ de $S_0$,
d'anneau $A_0$, de sorte que $\psi_\alpha^{-1}(U_0)$ est l'ouvert affine de
$S$ d'anneau $A=\Gamma(U_0,\mathscr A)$, puis on prend un élément $f\in A$ et
on considère dans $\operatorname{Spec}(A)$, identifié à
$\psi_\alpha^{-1}(U_0)$, l'ouvert $D(f)$. Ce sont ces ouverts $D(f)$ qui
forment une base de la topologie de $S$ (II, 1.3.1). Or, si on pose
$A_\lambda=\Gamma(U_0,\mathscr A_\lambda)$, on a
$A=\varinjlim A_\lambda$ (I, 1.3.9), donc il existe un indice $\lambda$ tel
que $f$ soit l'image canonique d'un élément $f_\lambda\in A_\lambda$~; on a
alors $D(f)=\psi_\lambda^{-1}(D(f_\lambda))$ (I, 1.2.2), et comme
$\psi_\lambda=\pi_\lambda\circ\psi$, notre assertion est démontrée.

\medskip
\noindent
(ii) Prouvons maintenant que $\psi$ est \emph{bijective}, ce qui achèvera la
démonstration. Comme $S$ est un espace de Kolmogoroff, il résulte déjà de (i)
que $\psi$ est injective, et il reste donc à montrer que $\psi$ est
surjective. On peut évidemment remplacer pour cela $T$ par un ouvert
$\pi_\alpha^{-1}(U_0)$, où $U_0$ est ouvert affine dans $S_\alpha=S_0$, donc
on peut se limiter au cas où les $S_\lambda$ et $S$ sont affines, autrement
dit $\mathscr A_\lambda$ est le faisceau associé à une $A_0$-algèbre
$A_\lambda$, et $\mathscr A$ le faisceau d'algèbres associé à
$A=\varinjlim A_\lambda$~; nous noterons encore
$\varphi_{\mu\lambda}:A_\lambda\to A_\mu$ et
$\varphi_\lambda:A_\lambda\to A$ les homomorphismes canoniques. Par
définition, un élément de $T$ est une famille
$(\mathfrak p_\lambda)_{\lambda\in L}$, où $\mathfrak p_\lambda$ est un
idéal premier de $A_\lambda$ et où l'on a
$\mathfrak p_\lambda=\varphi_{\mu\lambda}^{-1}(\mathfrak p_\mu)$ pour
$\lambda\leq\mu$. On sait alors (5.13.3 et 5.13.1) qu'il y a un idéal premier
$\mathfrak p$ de $A$ tel que
$\mathfrak p_\lambda=\varphi_\lambda^{-1}(\mathfrak p)$ pour tout
$\lambda\in L$, ce qui achève la démonstration.

En particulier, on a ainsi démontré le

\begin{corollary}[8.2.10]
\label{IV.8.2.10-fr}
Soit $(A_\lambda)_{\lambda\in L}$ un système inductif filtrant d'anneaux, et
soient $A=\varinjlim A_\lambda$, $\varphi_\lambda:A_\lambda\to A$ les
homomorphismes canoniques. L'application canonique
$\mathfrak p\mapsto(\varphi_\lambda^{-1}(\mathfrak p))$ est un
homéomorphisme de $\operatorname{Spec}(A)$ sur l'espace topologique
$\varprojlim\operatorname{Spec}(A_\lambda)$.
\end{corollary}

\begin{corollary}[8.2.11]
\label{IV.8.2.11-fr}
Avec les notations de (8.2.8), pour tout ouvert quasi-compact $U$ de $S$, il
existe un indice $\lambda$ et un ouvert quasi-compact $U_\lambda$ de
$S_\lambda$ tel que $U=\psi_\lambda^{-1}(U_\lambda)$.
\end{corollary}

Cela résulte de ce que, par définition de la topologie limite projective, les
$\psi_\lambda^{-1}(U_\lambda)$ ($U_\lambda$ ouvert quasi-compact de
$S_\lambda$) forment une base de la topologie de $S$, et de ce que l'ensemble
d'indices $L$ est filtrant.

\begin{corollary}[8.2.12]
\label{IV.8.2.12-fr}
Avec les notations de (8.2.8), la limite inductive du système inductif
d'homomorphismes
$\theta_\lambda^\sharp:\psi_\lambda^*(\mathcal O_{S_\lambda})\to
\mathcal O_S$ de faisceaux d'anneaux sur $S$ est un isomorphisme
\begin{equation}
\varinjlim_\lambda\psi_\lambda^*(\mathcal O_{S_\lambda})
\xrightarrow{\sim}\mathcal O_S.
\tag{8.2.12.1}\label{IV.8.2.12.1-fr}
\end{equation}
\end{corollary}

On peut évidemment supposer les $S_\lambda$ affines~; avec les notations de
la démonstration de (8.2.9), tout revient à voir que la limite inductive du
système d'applications canoniques
$(A_\lambda)_{\mathfrak p_\lambda}\to A_{\mathfrak p}$ est un isomorphisme,
ce qui n'est autre que (5.13.3, (ii)).

\oldpage[IV-3]{12}
\begin{proposition}[8.2.13]
\label{IV.8.2.13-fr}
Supposons que les morphismes $u_{\lambda\mu}:S_\mu\to S_\lambda$ soient des
immersions ouvertes, de sorte que $S_\mu$ s'identifie à un sous-préschéma
induit sur un ouvert de $S_\lambda$ pour $\lambda\leq\mu$. Alors, pour tout
$\alpha\in L$, l'espace sous-jacent au préschéma $S$ s'identifie au
sous-espace de $S_\alpha$ intersection des $S_\lambda$ pour
$\lambda\geq\alpha$, et le faisceau structural $\mathcal O_S$ au faisceau
induit $(\mathrm G,\mathrm{II},1.5)$ par $\mathcal O_{S_\alpha}$ sur cette
intersection~; plus généralement, pour tout $\mathcal O_{S_\alpha}$-Module
$\mathscr F_\alpha$, $u_\alpha^*(\mathscr F_\alpha)$ s'identifie au
$\mathcal O_S$-Module induit par $\mathscr F_\alpha$ sur $S$.
\end{proposition}

La première assertion résulte de (8.2.9), vu la définition d'une limite
projective d'espaces topologiques~; en outre tous les
$\psi_\lambda^*(\mathcal O_{S_\lambda})$ sont égaux au faisceau induit par
$\mathcal O_{S_\alpha}$ sur $S$ par définition $(0_{\mathrm I},3.7.1)$ et,
avec les notations de la démonstration de (8.2.9), les homomorphismes
$(A_\lambda)_{\mathfrak p_\lambda}\to
(A_\mu)_{\mathfrak p_\mu}$ sont bijectifs pour un système
$(\mathfrak p_\lambda)$ d'idéaux premiers correspondant à un même point de
$S$~; l'assertion relative à $\mathcal O_S$ se déduit donc de (8.2.12). La
dernière assertion résulte alors de la définition de
$u_\alpha^*(\mathscr F_\alpha)$ $(0_{\mathrm I},4.3.1)$.

\begin{remark}[8.2.14]
\label{IV.8.2.14-fr}
Les résultats de (8.2.9) et (8.2.12) montrent que $S$ est limite projective du
système projectif $(S_\lambda,u_{\lambda\mu})$ dans la catégorie de
\emph{tous} les espaces annelés (ou de tous les espaces annelés en anneaux
locaux). En effet, soit $Y$ un espace annelé, et considérons un système
projectif de morphismes d'espaces annelés $w_\lambda:Y\to S_\lambda$. Si
l'on pose $w_\lambda=(\rho_\lambda,\omega_\lambda)$, les $\rho_\lambda$
forment un système projectif d'applications continues et, en vertu de
(8.2.9), leur limite projective $\rho$ s'identifie à une application continue
$Y\to S$ telle que $\rho_\lambda=\psi_\lambda\circ\rho$. D'autre part, les
\[
\omega_\lambda^\sharp:
\rho_\lambda^*(\mathcal O_{S_\lambda})\to\mathcal O_Y
\]
forment un système inductif d'homomorphismes de faisceaux d'anneaux~; comme
on peut écrire
$\rho_\lambda^*(\mathcal O_{S_\lambda})=
\rho^*(\psi_\lambda^*(\mathcal O_{S_\lambda}))$ et que le foncteur $\rho^*$
est exact, la limite inductive des
$\rho_\lambda^*(\mathcal O_{S_\lambda})$ est $\rho^*(\mathcal O_S)$ en
vertu de (8.2.12), et il y a donc un homomorphisme unique
$\omega^\sharp:\rho^*(\mathcal O_S)\to\mathcal O_Y$ tel que
$\omega_\lambda^\sharp=\omega^\sharp\circ
\rho^*(\theta_\lambda^\sharp)$, ce qui démontre notre assertion.
\end{remark}

\subsection{Parties constructibles dans une limite projective de préschémas}
\label{subsection:IV.8.3-fr}

\begin{env}[8.3.1]
\label{IV.8.3.1-fr}
Dans toute la suite de ce paragraphe, nous supposons remplies les conditions
de (8.2.2), dont nous conservons les notations.
\end{env}

\begin{theorem}[8.3.2]
\label{IV.8.3.2-fr}
Pour tout $\lambda$, soient $E_\lambda$, $F_\lambda$ deux parties de
$S_\lambda$. Posons
\begin{equation}
E=\bigcap_\lambda u_\lambda^{-1}(E_\lambda),\qquad
F=\bigcup_\lambda u_\lambda^{-1}(F_\lambda).
\tag{8.3.2.1}\label{IV.8.3.2.1-fr}
\end{equation}
On suppose vérifiées les conditions suivantes~:
\begin{enumerate}
\item[(i)] Pour tout $\lambda$, $E_\lambda$ est pro-constructible et
$F_\lambda$ est ind-constructible (1.9.4).
\item[(ii)] Pour $\lambda\leq\mu$, on a
\[
E_\mu\subset u_{\lambda\mu}^{-1}(E_\lambda),\qquad
F_\mu\supset u_{\lambda\mu}^{-1}(F_\lambda).
\]
\item[(iii)] Il existe $\alpha\in L$ tel que $S_\alpha$ soit quasi-compact
(ce qui entraîne que $S_\lambda$ est quasi-compact pour
$\lambda\geq\alpha$).
\end{enumerate}
Alors, les propriétés suivantes sont équivalentes~:
\begin{enumerate}
\item[\emph{a)}] $E\subset F$.
\oldpage[IV-3]{13}
\item[\emph{b)}] Il existe $\lambda\geq\alpha$ tel que
$u_\lambda^{-1}(E_\lambda)\subset u_\lambda^{-1}(F_\lambda)$ (et alors on a
$u_\mu^{-1}(E_\mu)\subset u_\mu^{-1}(F_\mu)$ pour $\mu\geq\lambda$).
\item[\emph{c)}] Il existe $\lambda\geq\alpha$ tel que
$E_\lambda\subset F_\lambda$ (et alors on a $E_\mu\subset F_\mu$ pour
$\mu\geq\lambda$).
\end{enumerate}
\end{theorem}

Les remarques entre parenthèses dans \emph{b)} et \emph{c)} résultent de (ii).
Posons
\[
G_\lambda=E_\lambda\cap(S_\lambda-F_\lambda),\qquad
G=E\cap(S-F).
\]
Alors $G_\lambda$ est une partie pro-constructible de $S_\lambda$ (1.9.5,
(i)), et en vertu de (8.3.2.1) et de (ii), on a
\[
G_\mu\subset u_{\lambda\mu}^{-1}(G_\lambda)
\quad\text{pour }\lambda\leq\mu,
\qquad
G=\bigcap_\lambda u_\lambda^{-1}(G_\lambda).
\]
On est donc ramené à prouver le cas particulier de (8.3.2) correspondant à
$F_\lambda=\emptyset$ pour tout $\lambda$~:

\begin{corollary}[8.3.3]
\label{IV.8.3.3-fr}
Pour tout $\lambda$, soit $E_\lambda$ une partie pro-constructible de
$S_\lambda$ telle que, pour $\lambda\leq\mu$, on ait
$E_\mu\subset u_{\lambda\mu}^{-1}(E_\lambda)$. Supposons qu'il existe
$\alpha\in L$ tel que $S_\alpha$ soit quasi-compact. Alors les conditions
suivantes sont équivalentes~:
\begin{enumerate}
\item[\emph{a)}]
$E=\bigcap_\lambda u_\lambda^{-1}(E_\lambda)=\emptyset$.
\item[\emph{b)}] Il existe un $\lambda$ tel que
$u_\lambda^{-1}(E_\lambda)=\emptyset$ (et alors
$u_\mu^{-1}(E_\mu)=\emptyset$ pour $\mu\geq\lambda$).
\item[\emph{c)}] Il existe $\lambda$ tel que $E_\lambda=\emptyset$ (et alors
$E_\mu=\emptyset$ pour $\mu\geq\lambda$).
\end{enumerate}
\end{corollary}

Il est clair que \emph{c)} implique \emph{a)}. Prouvons que \emph{a)} entraîne
\emph{b)}~: puisque $S_\lambda$ est quasi-compact, il en est de même de $S$
(8.2.2)~; $E_\lambda$ étant pro-constructible, il en est de même de
$u_\lambda^{-1}(E_\lambda)$ (1.9.5, (vi))~; alors la famille filtrante
décroissante d'ensembles pro-constructibles $u_\lambda^{-1}(E_\lambda)$ a
une intersection vide, donc (1.9.9) l'un d'eux est vide.

Enfin, montrons que \emph{b)} entraîne \emph{c)}. Comme $S_\alpha$ est
quasi-compact et $L$ filtrant, on peut remplacer $S_\alpha$ par un ouvert
affine, donc supposer (8.2.2) que $S$ et les $S_\lambda$ pour
$\lambda\geq\alpha$ sont affines~; on a alors (1.9.2.1), pour
$\lambda\geq\alpha$,
\[
u_\lambda(S)=\bigcap_{\mu\geq\lambda}u_{\lambda\mu}(S_\mu),
\]
d'où
\[
E_\lambda\cap u_\lambda(S)=
\bigcap_{\mu\geq\lambda}
\bigl(E_\lambda\cap u_{\lambda\mu}(S_\mu)\bigr).
\]
Or, comme $u_\lambda$ et les $u_{\lambda\mu}$ sont quasi-compacts,
$u_\lambda(S)$ et $u_{\lambda\mu}(S_\mu)$ sont pro-constructibles dans
$S_\lambda$ (1.9.5, (vii)), donc les ensembles
$E_\lambda\cap u_{\lambda\mu}(S_\mu)$ pour $\mu\geq\lambda$ forment une
famille filtrante décroissante de parties pro-constructibles de $S_\lambda$.
Comme $S_\lambda$ est quasi-compact, l'hypothèse \emph{b)} entraîne que
l'intersection de cette famille est vide, donc (1.9.9) un des ensembles
$E_\lambda\cap u_{\lambda\mu}(S_\mu)$ est vide, donc
$E_\mu\subset u_{\lambda\mu}^{-1}(E_\lambda)$ est vide. C.Q.F.D.

\begin{corollary}[8.3.4]
\label{IV.8.3.4-fr}
Pour tout $\lambda$, soit $F_\lambda$ une partie ind-constructible de
$S_\lambda$ telle que pour $\lambda\leq\mu$ on ait
$u_{\lambda\mu}^{-1}(F_\lambda)\subset F_\mu$. Supposons qu'il existe
$\alpha\in L$ tel que $S_\alpha$ soit quasi-compact. Alors les conditions
suivantes sont équivalentes~:
\begin{enumerate}
\item[\emph{a)}] L'ensemble
$F=\bigcup_\lambda u_\lambda^{-1}(F_\lambda)$ est égal à $S$.
\item[\emph{b)}] Il existe un $\lambda$ tel que
$u_\lambda^{-1}(F_\lambda)=S$ (et alors
$u_\mu^{-1}(F_\mu)=S$ pour $\mu\geq\lambda$).
\oldpage[IV-3]{14}
\item[\emph{c)}] Il existe un $\lambda$ tel que $F_\lambda=S_\lambda$ (et
alors $F_\mu=S_\mu$ pour $\mu\geq\lambda$).
\end{enumerate}
\end{corollary}

Cela se déduit aussitôt de (8.3.3) par passage aux complémentaires.

\begin{corollary}[8.3.5]
\label{IV.8.3.5-fr}
Pour tout $\lambda$, soient $E_\lambda$, $F_\lambda$ deux parties
constructibles de $S_\lambda$ telles que, pour $\mu\geq\lambda$, on ait
$E_\mu\subset u_{\lambda\mu}^{-1}(E_\lambda)$ et
$F_\mu\supset u_{\lambda\mu}^{-1}(F_\lambda)$. Supposons qu'il existe un
indice $\alpha$ tel que $S_\alpha$ soit quasi-compact. Alors, pour que
$E\subset F$ (resp. $E=F$), il faut et il suffit qu'il existe $\lambda$ tel
que $E_\lambda\subset F_\lambda$ (resp. $E_\lambda=F_\lambda$), auquel cas
on a aussi $E_\mu\subset F_\mu$ (resp. $E_\mu=F_\mu$) pour
$\mu\geq\lambda$.
\end{corollary}

C'est un cas particulier de (8.3.2).

En particulier~:

\begin{corollary}[8.3.6]
\label{IV.8.3.6-fr}
Supposons qu'il existe un $\alpha$ tel que $S_\alpha$ soit quasi-compact.
Pour que $S=\emptyset$, il faut et il suffit qu'il existe un $\lambda$ tel
que $S_\lambda=\emptyset$.
\end{corollary}

\begin{corollary}[8.3.7]
\label{IV.8.3.7-fr}
On a, pour tout $\lambda$,
\begin{equation}
u_\lambda(S)=\bigcap_{\mu\geq\lambda}u_{\lambda\mu}(S_\mu).
\tag{8.3.7.1}\label{IV.8.3.7.1-fr}
\end{equation}
\end{corollary}

Il est clair que le premier membre de (8.3.7.1) est contenu dans le second.
Soit $s$ un point de $S_\lambda$ et posons
$X_\lambda=\operatorname{Spec}(k(s))$~; considérons le système projectif
$(X_\mu,v_{\mu\nu})$ où
$X_\mu=X_\lambda\times_{S_\lambda}S_\mu$ et
$v_{\mu\nu}=1\times u_{\mu\nu}$ pour
$\lambda\leq\mu\leq\nu$~; sa limite projective est
$X=X_\lambda\times_{S_\lambda}S$ et $v_\lambda=1\times u_\lambda$ est
l'application canonique $X\to X_\lambda$ (8.2.5). Si
$s\in\bigcap_{\mu\geq\lambda}u_{\lambda\mu}(S_\mu)$, cela entraîne que
$X_\mu\neq\emptyset$ pour tout $\mu\geq\lambda$ (I, 3.4.8)~; il résulte
alors de (8.3.6) que $X\neq\emptyset$, donc $s\in u_\lambda(S)$ par
(I, 3.4.8).

\begin{proposition}[8.3.8]
\label{IV.8.3.8-fr}
\begin{enumerate}
\item[(i)] Pour que le morphisme $u_\lambda:S\to S_\lambda$ soit dominant
(resp. surjectif), il faut et il suffit que pour $\mu\geq\lambda$ le
morphisme $u_{\lambda\mu}:S_\mu\to S_\lambda$ soit dominant (resp.
surjectif).
\item[(ii)] Si, pour un indice $\lambda$, les morphismes
$u_{\lambda\mu}:S_\mu\to S_\lambda$ sont plats (resp. fidèlement plats)
pour tout $\mu\geq\lambda$, alors le morphisme
$u_\lambda:S\to S_\lambda$ est plat (resp. fidèlement plat).
\item[(iii)] Supposons que les morphismes $u_\mu:S\to S_\mu$ soient
surjectifs pour $\mu\geq\lambda$. Pour que $u_\lambda$ soit un morphisme
ouvert (resp. universellement ouvert), il faut et il suffit que, pour tout
$\mu\geq\lambda$, $u_{\lambda\mu}$ soit un morphisme ouvert (resp.
universellement ouvert).
\end{enumerate}
\end{proposition}

\noindent
(i) Comme $u_\lambda(S)\subset u_{\lambda\mu}(S_\mu)$ pour
$\mu\geq\lambda$, la nécessité des conditions est triviale, et il résulte
aussitôt de (8.3.7.1) que si les $u_{\lambda\mu}$ sont surjectifs, il en est
de même de $u_\lambda$. Supposons maintenant les $u_{\lambda\mu}$ dominants
pour $\mu\geq\lambda$, et considérons dans $S_\lambda$ un ouvert
quasi-compact non vide $U_\lambda$~; alors les
$U_\mu=u_{\lambda\mu}^{-1}(U_\lambda)$ pour $\mu\geq\lambda$ forment un
système projectif dont la limite projective est
$U=u_\lambda^{-1}(U_\lambda)$ (8.2.5). Par hypothèse les $U_\mu$ sont tous
non vides, donc il en est de même de $U$ par (8.3.6), ce qui prouve que
$u_\lambda$ est dominant.

\medskip
\noindent
(ii) En vertu de (i), il suffit de considérer le cas où les
$u_{\lambda\mu}$ sont plats~; on peut alors se borner au cas où $S_\lambda$
est affine, donc aussi les $S_\mu$ pour $\mu\geq\lambda$ et $S$, et
l'assertion résulte alors de (2.1.2) et $(0_{\mathrm I},6.2.3)$.

\medskip
\noindent
(iii) En vertu de (8.2.5) et de (I, 3.5.2), il suffit de traiter le cas des
morphismes ouverts. Comme $u_\lambda=u_{\lambda\mu}\circ u_\mu$ et que
$u_\mu$ est surjectif, on sait que si $u_\lambda$ est ouvert il en est de
même de $u_{\lambda\mu}$ pour $\mu\geq\lambda$ (Bourbaki,
\emph{Top. gén.}, chap. $\mathrm I^{\mathrm{er}}$, $3^{\mathrm e}$ éd.,
\S~5, n\textsuperscript{o} 1, prop. 1).

\oldpage[IV-3]{15}
Inversement, pour montrer que $u_\lambda$ est ouvert lorsque tous les
$u_{\lambda\mu}$ le sont pour $\mu\geq\lambda$, il suffit de voir que pour
tout ouvert quasi-compact $V$ de $S$, $u_\lambda(V)$ est ouvert dans
$S_\lambda$~; mais il existe alors $\mu\geq\lambda$ tel que
$V=u_\mu^{-1}(V_\mu)$, où $V_\mu$ est ouvert dans $S_\mu$ (8.2.11)~; comme
$u_\mu$ est surjectif, on a $V_\mu=u_\mu(V)$ et
$u_\lambda(V)=u_{\lambda\mu}(u_\mu(V))$ est donc ouvert par hypothèse.

On notera qu'il peut se faire que tous les $u_{\lambda\mu}$ soient ouverts
sans qu'il en soit ainsi de $u_\lambda$ si les $u_\mu$ ne sont pas
surjectifs. On en a un exemple en considérant un anneau intègre $A$ qui n'est
pas un corps, et son corps des fractions $K$, qui est limite inductive des
$A_f$, où $f$ parcourt $A-\{0\}$~; si l'on pose
$S_1=\operatorname{Spec}(A)$, $S_f=\operatorname{Spec}(A_f)$, on a
$S=\varprojlim S_f=\operatorname{Spec}(K)$, et le morphisme $S\to S_1$ n'est
pas ouvert, bien que les morphismes $S_f\to S_1$ le soient.

\begin{env}[8.3.9]
\label{IV.8.3.9-fr}
Pour tout préschéma $X$, nous noterons comme d'ordinaire $\mathfrak P(X)$
l'ensemble des parties de l'ensemble sous-jacent à $X$, par $\mathfrak C(X)$
(resp. $\mathfrak D_c(X)$, $\mathfrak F_c(X)$, $\mathfrak{LF}_c(X)$)
l'ensemble des parties constructibles (resp. constructibles et ouvertes,
resp. constructibles et fermées, resp. constructibles et localement fermées)
de $X$. Il est clair que
$(\mathfrak P(S_\lambda),u_{\lambda\mu}^{-1})$ est un \emph{système inductif
d'ensembles} et que les applications
$u_\lambda^{-1}:\mathfrak P(S_\lambda)\to\mathfrak P(S)$ forment un système
inductif d'applications, d'où, en passant à la limite inductive, une
application canonique
\begin{equation}
u_{\mathfrak P}:\varinjlim\mathfrak P(S_\lambda)\to\mathfrak P(S).
\tag{8.3.9.1}\label{IV.8.3.9.1-fr}
\end{equation}
En outre, il résulte de (1.8.2) que $u_{\lambda\mu}^{-1}$ applique
$\mathfrak C(S_\lambda)$ (resp. $\mathfrak D_c(S_\lambda)$,
$\mathfrak F_c(S_\lambda)$, $\mathfrak{LF}_c(S_\lambda)$) dans
$\mathfrak C(S_\mu)$ (resp. $\mathfrak D_c(S_\mu)$,
$\mathfrak F_c(S_\mu)$, $\mathfrak{LF}_c(S_\mu)$) et que
$u_\lambda^{-1}$ applique $\mathfrak C(S_\lambda)$ (resp.
$\mathfrak D_c(S_\lambda)$, $\mathfrak F_c(S_\lambda)$,
$\mathfrak{LF}_c(S_\lambda)$) dans $\mathfrak C(S)$ (resp.
$\mathfrak D_c(S)$, $\mathfrak F_c(S)$, $\mathfrak{LF}_c(S)$). On a donc par
restriction de (8.3.9.1), des applications canoniques
\begin{align}
\varinjlim\mathfrak C(S_\lambda)&\longrightarrow\mathfrak C(S)
\tag{8.3.9.2}\label{IV.8.3.9.2-fr}\\
\varinjlim\mathfrak D_c(S_\lambda)&\longrightarrow\mathfrak D_c(S)
\tag{8.3.9.3}\label{IV.8.3.9.3-fr}\\
\varinjlim\mathfrak F_c(S_\lambda)&\longrightarrow\mathfrak F_c(S)
\tag{8.3.9.4}\label{IV.8.3.9.4-fr}\\
\varinjlim\mathfrak{LF}_c(S_\lambda)&\longrightarrow\mathfrak{LF}_c(S).
\tag{8.3.9.5}\label{IV.8.3.9.5-fr}
\end{align}
\end{env}

\begin{env}[8.3.10]
\label{IV.8.3.10-fr}
Soit $g_\alpha:X_\alpha\to S_\alpha$ un morphisme~; avec les notations de
(8.2.5) on a comme ci-dessus une application canonique
$v_{\mathfrak P}:\varinjlim\mathfrak P(X_\lambda)\to\mathfrak P(X)$~;
d'autre part, on a des morphismes de projection
$g_\lambda:X_\lambda\to S_\lambda$ pour tout $\lambda\geq\alpha$ et un
morphisme de projection $g:X\to S$. Il est clair que les
$g_\lambda^{-1}:\mathfrak P(S_\lambda)\to\mathfrak P(X_\lambda)$ forment un
système inductif d'applications, et que les diagrammes
\[
\xymatrix{
  \mathfrak P(S_\lambda) \ar[r]^{u_\lambda^{-1}} \ar[d]_{g_\lambda^{-1}}
    & \mathfrak P(S) \ar[d]^{g^{-1}} \\
  \mathfrak P(X_\lambda) \ar[r]_{v_\lambda^{-1}} & \mathfrak P(X)
}
\]
\oldpage[IV-3]{16}
sont commutatifs~; on en déduit donc par passage à la limite inductive un
diagramme commutatif
\begin{equation}
\xymatrix{
  \varinjlim\mathfrak P(S_\lambda) \ar[r]^{u_{\mathfrak P}}
    \ar[d]_{\varinjlim g_\lambda^{-1}}
    & \mathfrak P(S) \ar[d]^{g^{-1}} \\
  \varinjlim\mathfrak P(X_\lambda) \ar[r]_{v_{\mathfrak P}}
    & \mathfrak P(X)
}
\tag{8.3.10.1}\label{IV.8.3.10.1-fr}
\end{equation}
et il résulte de (1.8.2) que l'on a des diagrammes analogues en remplaçant
$\mathfrak P$ par $\mathfrak C$, $\mathfrak D_c$, $\mathfrak F_c$ ou
$\mathfrak{LF}_c$.
\end{env}

Il résulte de (8.3.5) que sous l'hypothèse que pour un $\alpha\in L$,
$S_\alpha$ est \emph{quasi-compact}, l'application canonique (8.3.9.2) est
\emph{injective}. En outre~:

\begin{theorem}[8.3.11]
\label{IV.8.3.11-fr}
Supposons qu'il existe un $\alpha\in L$ tel que $S_\alpha$ soit quasi-compact
et quasi-séparé. Alors les applications canoniques (8.3.9.2), (8.3.9.3),
(8.3.9.4) et (8.3.9.5) sont \emph{bijectives}.
\end{theorem}

En vertu de la remarque précédente, il reste à prouver que ces applications
sont \emph{surjectives}~; comme toute partie constructible de $S$ est réunion
finie d'ensembles de la forme $U\cap\complement V$, où $U$ et $V$ sont
ouverts et constructibles, il suffira de prouver que (8.3.9.3) est surjective
pour qu'il en soit de même de (8.3.9.2) (et aussi de (8.3.9.4), par passage
aux complémentaires). Or, comme les morphismes $S_\lambda\to S_\alpha$ et
$S\to S_\alpha$ sont affines, donc séparés, les $S_\lambda$ pour
$\lambda\geq\alpha$ et $S$ sont quasi-compacts et quasi-séparés (1.2.2), et
l'on sait que les parties ouvertes constructibles dans un tel préschéma ne
sont autres que les parties ouvertes quasi-compactes (1.8.1). La conclusion
résulte donc de (8.2.11) sauf pour l'application (8.3.9.5). Pour prouver que
cette dernière est surjective, considérons une partie $Z$ localement fermée
et constructible dans $S$~; $Z$ est donc quasi-compact
$(0_{\mathrm{III}},9.1.4)$. Comme tout point $x\in Z$ admet par hypothèse un
voisinage ouvert quasi-compact $V_x$ dans $S$ tel que $Z\cap V_x$ soit fermé
dans $V_x$, on peut couvrir $Z$ par un nombre fini des $V_x$~; autrement dit,
il y a un ouvert quasi-compact $U$ contenant $Z$ et tel que $Z$ soit fermé
dans $U$~; comme $Z$ est constructible dans $S$, il l'est aussi dans $U$
$(0_{\mathrm{III}},9.1.8)$. On sait (8.2.11) qu'il y a un indice $\lambda$ et
un ouvert quasi-compact $U_\lambda$ dans $S_\lambda$ tels que
$U=u_\lambda^{-1}(U_\lambda)$. Appliquant à $U$ (qui est limite projective
des $U_\mu=u_{\lambda\mu}^{-1}(U_\lambda)$ pour $\mu\geq\lambda$) le fait
que l'application (8.3.9.4) est surjective, on voit qu'il existe un
$\mu\geq\lambda$ et un ensemble fermé constructible $Z_\mu$ dans $U_\mu$
tels que $Z=u_\mu^{-1}(Z_\mu)$. Mais comme l'immersion canonique
$U_\mu\to S_\mu$ est quasi-compacte par hypothèse (1.2.7), elle est de
présentation finie (1.6.2, (i)), et $Z_\mu$ est aussi une partie constructible
\emph{de} $S_\mu$ en vertu de (1.8.4) et (1.8.1)~; comme $Z_\mu$ est
évidemment localement fermée dans $S_\mu$, cela achève la démonstration.

\begin{corollary}[8.3.12]
\label{IV.8.3.12-fr}
Supposons qu'il existe un $\alpha$ tel que $S_\alpha$ soit quasi-compact, et
soit, pour tout $\lambda$, $Z_\lambda$ une partie constructible de
$S_\lambda$ telle que
$Z_\mu=u_{\lambda\mu}^{-1}(Z_\lambda)$ pour $\mu\geq\lambda$. Si
$Z=u_\lambda^{-1}(Z_\lambda)$,
\oldpage[IV-3]{17}
alors, pour que $Z$ soit ouverte (resp. fermée, resp. localement fermée) dans
$S$, il faut et il suffit qu'il existe $\lambda\geq\alpha$ tel que
$Z_\lambda$ le soit dans $S_\lambda$.
\end{corollary}

Couvrons $S_\alpha$ par un nombre fini d'ouverts affines
$U_\alpha^{(i)}$~; alors les $U^{(i)}=u_\alpha^{-1}(U_\alpha^{(i)})$ forment
un recouvrement ouvert affine de $S$, et pour que $Z$ soit ouvert (resp.
fermé, resp. localement fermé) dans $S$, il faut et il suffit que chacun des
$Z\cap U^{(i)}$ le soit dans $U^{(i)}$. Comme $L$ est filtrant et que chacun
des $Z\cap U^{(i)}$ est constructible dans $U^{(i)}$
$(0_{\mathrm{III}},9.1.8)$, on peut se borner à démontrer le corollaire
lorsque $S_\alpha$ est affine, donc quasi-compact et quasi-séparé~; mais alors
il résulte de (8.3.11).

\begin{proposition}[8.3.13]
\label{IV.8.3.13-fr}
Supposons que les morphismes $u_{\lambda\mu}:S_\mu\to S_\lambda$ soient
plats pour $\lambda\leq\mu$, et qu'il existe $\alpha$ tel que $S_\alpha$ soit
quasi-compact. Pour tout $\lambda$, soient $Z'_\lambda$, $Z''_\lambda$ deux
parties pro-constructibles de $S_\lambda$, telles que
$Z'_\mu=u_{\lambda\mu}^{-1}(Z'_\lambda)$ et
$Z''_\mu=u_{\lambda\mu}^{-1}(Z''_\lambda)$ pour $\mu\geq\lambda$~; on
suppose en outre que $\overline{Z'_\alpha}$ est constructible dans
$S_\alpha$. Soient $Z'=u_\lambda^{-1}(Z'_\lambda)$,
$Z''=u_\lambda^{-1}(Z''_\lambda)$~; pour que
$Z''\subset\overline{Z'}$, il faut et il suffit qu'il existe
$\lambda\geq\alpha$ tel que
$Z''_\lambda\subset\overline{Z'_\lambda}$.
\end{proposition}

En effet, on sait que $u_\lambda:S\to S_\lambda$ est aussi un morphisme plat
pour tout $\lambda$ (8.3.8)~; comme $Z'_\lambda$ est pro-constructible,
l'adhérence de $Z'_\mu$ pour $\mu\geq\lambda$ (resp. de $Z'$) dans $S_\mu$
(resp. $S$) est égale à
$u_{\lambda\mu}^{-1}(\overline{Z'_\lambda})$ (resp.
$u_\lambda^{-1}(\overline{Z'_\lambda})$) (2.3.10). Comme les
$u_{\lambda\mu}^{-1}(\overline{Z'_\lambda})$ et
$u_\lambda^{-1}(\overline{Z'_\lambda})$ sont constructibles (1.8.2), la
conclusion résulte de (8.3.2).

\subsection{Critères d'irréductibilité et de connexion pour les limites
projectives de préschémas}
\label{subsection:IV.8.4-fr}

\begin{proposition}[8.4.1]
\label{IV.8.4.1-fr}
Supposons qu'il existe un indice $\alpha$ tel que $S_\alpha$ soit
quasi-compact.
\begin{enumerate}
\item[(i)] Si $S$ n'est pas irréductible et si de plus l'espace sous-jacent à
$S$ est noethérien et $S_\alpha$ quasi-séparé, il existe
$\lambda\geq\alpha$ tel que, pour $\mu\geq\lambda$, $S_\mu$ ne soit pas
irréductible.
\item[(ii)] Si $S$ n'est pas connexe, il existe $\lambda\geq\alpha$ tel que,
pour $\mu\geq\lambda$, $S_\mu$ ne soit pas connexe.
\end{enumerate}
\end{proposition}

Supposons que $S$ soit réunion de deux ensembles fermés $S'$, $S''$ distincts
de $S$ (resp. fermés disjoints et non vides). Dans le cas (i), $S'$ et $S''$
sont constructibles puisque l'espace $S$ est noethérien. En vertu de
(8.3.11), il existe donc un $\lambda\geq\alpha$ et deux ensembles fermés
constructibles $S'_\lambda$, $S''_\lambda$ de $S_\lambda$ tels que
$S'=u_\lambda^{-1}(S'_\lambda)$,
$S''=u_\lambda^{-1}(S''_\lambda)$~; comme $S=S'\cup S''$, il résulte aussi
de (8.3.11) que l'on peut supposer que
$S_\lambda=S'_\lambda\cup S''_\lambda$~; comme $S'_\lambda$ et
$S''_\lambda$ sont distincts de $S_\lambda$, cela prouve que $S_\lambda$
n'est pas irréductible.

Dans le cas (ii), $S'$ et $S''$ sont ouverts quasi-compacts, donc, en vertu
de (8.2.11), il existe un $\lambda\geq\alpha$ et deux ensembles ouverts
quasi-compacts $S'_\lambda$, $S''_\lambda$ de $S_\lambda$ tels que
$S'=u_\lambda^{-1}(S'_\lambda)$,
$S''=u_\lambda^{-1}(S''_\lambda)$. D'ailleurs, comme $S'$ et $S''$ sont
ouverts et fermés dans $S$, ils sont à la fois pro-constructibles et
ind-constructibles (1.9.6), donc constructibles (1.9.11), et il résulte donc
de (8.3.5) que l'on peut supposer $\lambda$ pris tel que
$S_\lambda=S'_\lambda\cup S''_\lambda$ et
$S'_\lambda\cap S''_\lambda=\emptyset$, ce qui montre que $S_\lambda$
n'est pas connexe.

\begin{proposition}[8.4.2]
\label{IV.8.4.2-fr}
Supposons que l'espace sous-jacent à $S$ soit noethérien et que l'une des deux
conditions suivantes soit vérifiée~:
\begin{enumerate}
\item[\emph{a)}] Pour $\lambda\leq\mu$,
$u_{\lambda\mu}:S_\mu\to S_\lambda$ est dominant, et il existe $\alpha$ tel
que $S_\alpha$ soit quasi-compact.
\oldpage[IV-3]{18}
\item[\emph{b)}] Il existe $\alpha$ tel que l'espace sous-jacent à
$S_\alpha$ soit noethérien, et pour $\mu\geq\lambda$, $u_{\lambda\mu}$ est
un homéomorphisme de $S_\mu$ sur un sous-espace de $S_\lambda$.
\end{enumerate}
Sous ces conditions, il existe un $\lambda$ tel que, pour tout
$\mu\geq\lambda$~:
\begin{enumerate}
\item[(i)] Pour toute composante irréductible $Y_i$ de $S$
$(1\leq i\leq m)$, $\overline{u_\mu(Y_i)}$ est une composante irréductible
de $S_\mu$, et l'application $Y_i\mapsto\overline{u_\mu(Y_i)}$ est une
bijection de l'ensemble des composantes irréductibles de $S$ sur l'ensemble
des composantes irréductibles de $S_\mu$.
\item[(ii)] Pour toute composante connexe $C_j$ de $S$
$(1\leq j\leq n)$, $\overline{u_\mu(C_j)}$ est une composante connexe de
$S_\mu$, et l'application $C_j\mapsto\overline{u_\mu(C_j)}$ est une
bijection de l'ensemble des composantes connexes de $S$ sur l'ensemble des
composantes connexes de $S_\mu$.
\end{enumerate}
\end{proposition}

Nous établirons d'abord le

\begin{lemma}[8.4.2.1]
\label{IV.8.4.2.1-fr}
Sous la condition \emph{a)} ou \emph{b)} de (8.4.2), il existe $\lambda$ tel
que, pour $\mu\geq\lambda$, $u_\mu:S\to S_\mu$ soit dominant.
\end{lemma}

Dans le cas \emph{a)}, cela a déjà été prouvé sans supposer l'espace $S$
noethérien (8.3.8, (i)). Dans le cas \emph{b)}, posons
$Z_\alpha=\overline{u_\alpha(S)}$~; en tant que partie fermée de l'espace
noethérien $S_\alpha$, $Z_\alpha$ est constructible, et comme
$u_\alpha^{-1}(Z_\alpha)=S$, il résulte de (8.3.5) qu'il existe
$\lambda\geq\alpha$ tel que, pour $\mu\geq\lambda$, on ait
$Z_\mu=u_{\alpha\mu}^{-1}(Z_\alpha)=S_\mu$. Mais comme $u_{\alpha\mu}$ est
un homéomorphisme de $S_\mu$ sur un sous-espace de $S_\alpha$, et comme
l'application composée $S\to Z_\mu\to Z_\alpha$ est dominante, il en est de
même de $S\to Z_\mu=S_\mu$.

Ce lemme étant démontré, on peut supposer que pour tout $\lambda$,
$u_\lambda$ est un morphisme dominant.

\medskip
\noindent
(i) Chacun des $S_\lambda$ est réunion des
$\overline{u_\lambda(Y_i)}$, qui sont irréductibles. D'autre part, si $U_i$
est l'ouvert de $S$ complémentaire de la réunion des $Y_j$ d'indice $j\neq i$
$(1\leq i\leq m)$, les $U_i$ sont deux à deux disjoints et
$Y_i=\overline{U_i}$ $(0_{\mathrm I},2.1.6)$. Comme l'espace sous-jacent à
$S$ est noethérien, les $U_i$ sont quasi-compacts, donc il existe un indice
$\lambda$ et des ouverts $U_{i\lambda}$ de $S_\lambda$ tels que
$U_i=u_\lambda^{-1}(U_{i\lambda})$ pour $1\leq i\leq m$ (8.2.11). On en
conclut que si l'on pose
$U_{i\mu}=u_{\lambda\mu}^{-1}(U_{i\lambda})$ pour $\lambda\leq\mu$, les
$U_{i\mu}$ sont deux à deux disjoints, car les
$U_i=u_\mu^{-1}(U_{i\mu})$ le sont et $u_\mu$ est dominant. Par suite,
aucune des adhérences $\overline{U_{i\mu}}$ n'est contenue dans une autre et
$u_\mu(U_i)$ est dense dans $U_{i\mu}$ puisque $u_\mu$ est dominant~; on a
donc
$\overline{U_{i\mu}}=\overline{u_\mu(Y_i)}$, ce qui prouve que les
$\overline{U_{i\mu}}$ sont les composantes irréductibles de $S_\mu$
$(0_{\mathrm I},2.1.7)$ et achève la démonstration.

\medskip
\noindent
(ii) Comme l'espace $S$ est noethérien, les $C_j$ sont ouverts et fermés dans
$S$ et quasi-compacts~; le même raisonnement que dans (i) montre donc qu'il
existe un $\lambda$ et des ouverts $V_{j\lambda}$ de $S_\lambda$ tels que
$C_j=u_\lambda^{-1}(V_{j\lambda})$ pour $1\leq j\leq n$. On voit aussi,
comme dans (i), que si l'on pose
$V_{j\mu}=u_{\lambda\mu}^{-1}(V_{j\lambda})$ pour $\mu\geq\lambda$, les
$V_{j\mu}$ sont deux à deux disjoints, et $u_\mu(C_j)$ est dense dans
$V_{j\mu}$~; cela entraîne que $V_{j\mu}$ est connexe. En outre, il résulte
de (8.3.4) que pour $\mu$ assez grand, la réunion des $V_{j\mu}$ est
$S_\mu$, puisque tout ouvert dans un préschéma est ind-constructible (1.9.6).
Les $V_{j\mu}$ sont donc les composantes connexes de $S_\mu$, ce qui achève
la démonstration.

On notera que si les morphismes $u_{\lambda\mu}$ sont des \emph{immersions},
ils vérifieront en particulier la condition \emph{b)} de (8.4.2).

\oldpage[IV-3]{19}
\begin{corollary}[8.4.3]
\label{IV.8.4.3-fr}
Supposons vérifiée l'une des conditions \emph{a)}, \emph{b)} de (8.4.2),
l'espace sous-jacent à $S$ étant noethérien~; alors, pour que $S$ soit
irréductible, il faut et il suffit qu'il existe $\lambda$ tel que $S_\mu$ le
soit pour tout $\mu\geq\lambda$.
\end{corollary}

\begin{proposition}[8.4.4]
\label{IV.8.4.4-fr}
Supposons qu'il existe un $\alpha$ tel que $S_\alpha$ soit quasi-compact et
que, pour $\lambda\leq\mu$, $u_{\lambda\mu}:S_\mu\to S_\lambda$ soit dominant.
Alors, pour que $S$ soit connexe, il faut et il suffit qu'il existe $\lambda$
tel que $S_\mu$ le soit pour tout $\mu\geq\lambda$.
\end{proposition}

\begin{proof}
La condition est suffisante en vertu de (8.4.1)~; d'autre part, on a vu
(8.3.8, (i)) que $u_\lambda:S\to S_\lambda$ est dominant pour $\lambda$ assez
grand, donc, si $S$ est connexe, il en est de même de $S_\lambda$, puisque
$u_\lambda(S)$ est dense dans $S_\lambda$ et connexe.
\end{proof}

\begin{corollary}[8.4.5]
\label{IV.8.4.5-fr}
Soient $k$ un corps, $X$ un $k$-préschéma quasi-compact. Pour que $X$ soit
géométriquement connexe (4.5.2), il faut et il suffit que, pour toute extension
finie et séparable $K$ de $k$, $X\otimes_k K$ soit connexe.
\end{corollary}

\begin{proof}
La condition est trivialement nécessaire. Pour voir qu'elle est suffisante,
nous devons prouver que si $\Omega$ est une clôture algébrique de $k$,
$X\otimes_k\Omega$ est connexe (4.5.1). Or, $\Omega$ est limite inductive
filtrante des sous-extensions finies $k'$ de $k$, et pour
$k\subset k'\subset k''\subset\Omega$, le morphisme
$X\otimes_k k''\to X\otimes_k k'$ est surjectif. On est donc ramené, en vertu
de (8.4.4), à prouver que $X\otimes_k k'$ est connexe pour toute extension
finie $k'$ de $k$. Mais si $K$ est la plus grande extension séparable contenue
dans $k'$, le morphisme $X\otimes_k k'\to X\otimes_k K$ est fini, surjectif et
radiciel, donc (2.4.5) un homéomorphisme, et comme $X\otimes_k K$ est connexe
par hypothèse, il en est de même de $X\otimes_k k'$.
\end{proof}

\begin{remark}[8.4.6]
\label{IV.8.4.6-fr}
\begin{enumerate}[label=(\roman*)]
  \item La démonstration de (8.4.2) montre que la conclusion de cette
  proposition est valable si l'on suppose que l'espace sous-jacent à $S$ est
  noethérien, qu'il existe $\alpha$ tel que $S_\alpha$ soit quasi-compact, et
  enfin que les $u_\lambda:S\to S_\lambda$ sont dominants.

  \item Par contre, la conclusion de (8.4.2) peut être en défaut lorsque les
  $u_\lambda$ ne sont pas dominants pour $\lambda$ assez grand, même lorsque
  les $S_\lambda$ et $S$ sont noethériens, ainsi que le montre l'exemple
  suivant. Prenons pour ensemble d'indices $\mathbf N$, tous les $S_n$ étant
  égaux à
  $\operatorname{Spec}(A\times K)=\operatorname{Spec}(A)\amalg
  \operatorname{Spec}(K)$, où $K$ est un corps, $A$ une $K$-algèbre
  quelconque, et tous les morphismes $u_{n,n+1}$ étant égaux au même morphisme
  correspondant à l'homomorphisme $(x,y)\mapsto(j(y),y)$ de $A\times K$ dans
  lui-même, où $j:K\to A$ est l'homomorphisme canonique. On vérifie aisément
  que la limite inductive de ce système d'anneaux est $K$, l'homomorphisme
  canonique $u_n$ correspondant à la seconde projection $A\times K\to K$. On
  voit donc que $S=\operatorname{Spec}(K)$ est irréductible bien qu'aucun des
  $S_n$ ne soit connexe.
\end{enumerate}
\end{remark}

\subsection{Modules de présentation finie sur une limite projective de préschémas}
\label{subsection:IV.8.5-fr}

\begin{env}[8.5.1]
\label{IV.8.5.1-fr}
Nous conservons toujours les notations de (8.2.2)~; nous nous bornerons en
outre au cas où $S_0$ est l'un des $S_\lambda$, auquel on peut toujours se
ramener.

\emph{Lorsque, dans cette section, nous considérerons une famille
$(\mathcal F_\lambda)$, où, pour tout $\lambda\in L$, $\mathcal F_\lambda$ est
un $\mathcal O_{S_\lambda}$-Module, il sera sous-entendu que cette famille
vérifie la condition}
\begin{equation}
  \mathcal F_\mu=u_{\lambda\mu}^*(\mathcal F_\lambda)
  \qquad\text{pour }\lambda\leq\mu.
  \tag{8.5.1.1}\label{IV.8.5.1.1-fr}
\end{equation}

\oldpage[IV-3]{20}
Nous poserons alors
\begin{equation}
  \mathcal F=u_\lambda^*(\mathcal F_\lambda)
  \tag{8.5.1.2}\label{IV.8.5.1.2-fr}
\end{equation}
qui est un $\mathcal O_S$-Module ne dépendant pas de l'indice $\lambda\in L$,
en vertu de l'hypothèse (8.5.1.1).

Soient maintenant $(\mathcal F_\lambda)$, $(\mathcal G_\lambda)$ deux telles
familles de $\mathcal O_{S_\lambda}$-Modules. Il est clair que les applications
$f_\lambda\mapsto u_{\lambda\mu}^*(f_\lambda)$ de
$\operatorname{Hom}_{S_\lambda}(\mathcal F_\lambda,\mathcal G_\lambda)$ dans
$\operatorname{Hom}_{S_\mu}(\mathcal F_\mu,\mathcal G_\mu)$ définissent alors
un \emph{système inductif} de groupes abéliens
$(\operatorname{Hom}_{S_\lambda}(\mathcal F_\lambda,\mathcal G_\lambda),
u_{\lambda\mu}^*)$, et que les applications
$f_\lambda\mapsto u_\lambda^*(f_\lambda)$ forment un \emph{système inductif}
d'homomorphismes de groupes abéliens
\[
  u_\lambda^*:\operatorname{Hom}_{S_\lambda}
  (\mathcal F_\lambda,\mathcal G_\lambda)
  \longrightarrow\operatorname{Hom}_S(\mathcal F,\mathcal G)~;
\]
d'où, en passant à la limite inductive, un homomorphisme canonique de groupes
abéliens
\begin{equation}
  u_{\mathcal F,\mathcal G}^*:
  \varinjlim\operatorname{Hom}_{S_\lambda}
  (\mathcal F_\lambda,\mathcal G_\lambda)
  \longrightarrow\operatorname{Hom}_S(\mathcal F,\mathcal G).
  \tag{8.5.1.3}\label{IV.8.5.1.3-fr}
\end{equation}

Notons que lorsque $\mathcal F_\lambda=\mathcal O_{S_\lambda}$, la condition
(8.5.1.1) est vérifiée, et que l'on a $\mathcal F=\mathcal O_S$~;
l'homomorphisme (8.5.1.3) donne donc $(0_{\mathrm I},5.1.1)$ un homomorphisme
canonique de groupes abéliens
\begin{equation}
  u_{\mathcal G}:\varinjlim\Gamma(S_\lambda,\mathcal G_\lambda)
  \longrightarrow\Gamma(S,\mathcal G).
  \tag{8.5.1.4}\label{IV.8.5.1.4-fr}
\end{equation}
\end{env}

\begin{theorem}[8.5.2]
\label{IV.8.5.2-fr}
\begin{enumerate}[label=(\roman*)]
  \item Supposons $S_0$ quasi-compact (resp. quasi-compact et quasi-séparé) et
  que, pour un $\lambda\in L$, $\mathcal F_\lambda$ soit quasi-cohérent et de
  type fini (resp. de présentation finie) et $\mathcal G_\lambda$
  quasi-cohérent. Alors l'homomorphisme $u_{\mathcal F,\mathcal G}^*$ est
  injectif (resp. bijectif).

  \item Supposons $S_0$ quasi-compact et quasi-séparé. Pour tout
  $\mathcal O_S$-Module quasi-cohérent $\mathcal F$ de présentation finie, il
  existe un $\lambda\in L$ et un $\mathcal O_{S_\lambda}$-Module
  quasi-cohérent $\mathcal F_\lambda$ de présentation finie tels que
  $\mathcal F$ soit isomorphe à $u_\lambda^*(\mathcal F_\lambda)$.
\end{enumerate}
\end{theorem}

\begin{proof}
\begin{enumerate}[label=(\roman*)]
  \item On peut évidemment se borner au cas où $S_0=S_\lambda$ puisque les
  morphismes $u_{0\lambda}:S_\lambda\to S_0$ sont affines, donc quasi-compacts
  et séparés. Considérons d'abord le cas où
  $S_0=\operatorname{Spec}(A_0)$ est \emph{affine}. Alors l'assertion (i)
  équivaut au

  \begin{lemma}[8.5.2.1]
  \label{IV.8.5.2.1-fr}
  Soient $A_0$ un anneau, $(A_\lambda)_{\lambda\in L}$ un système inductif de
  $A_0$-algèbres, $A=\varinjlim A_\lambda$~; soient $M_0$, $N_0$ deux
  $A_0$-modules, et posons
  $M_\lambda=M_0\otimes_{A_0}A_\lambda$,
  $N_\lambda=N_0\otimes_{A_0}A_\lambda$,
  $M=M_0\otimes_{A_0}A=\varinjlim M_\lambda$,
  $N=N_0\otimes_{A_0}A=\varinjlim N_\lambda$. Si $M_0$ est de type fini
  (resp. de présentation finie), l'homomorphisme canonique
  \begin{equation}
    \varinjlim\operatorname{Hom}_{A_\lambda}(M_\lambda,N_\lambda)
    \longrightarrow\operatorname{Hom}_A(M,N)
    \tag{8.5.2.2}\label{IV.8.5.2.2-fr}
  \end{equation}
  est injectif (resp. bijectif).
  \end{lemma}

  On sait en effet (Bourbaki, \emph{Alg.}, chap. II, 3\textsuperscript{e} éd.,
  \S~5, n\textsuperscript{o}~1) que l'on a des isomorphismes canoniques
  fonctoriels
  \[
    \operatorname{Hom}_{A_\lambda}(M_\lambda,N_\lambda)
    \xrightarrow{\sim}\operatorname{Hom}_{A_0}(M_0,N_\lambda),
    \qquad
    \operatorname{Hom}_A(M,N)
    \xrightarrow{\sim}\operatorname{Hom}_{A_0}(M_0,N)
  \]

  \oldpage[IV-3]{21}
  si bien que l'homomorphisme (8.5.2.2) n'est autre, à des isomorphismes
  canoniques près, que l'homomorphisme canonique
  \begin{equation}
    \varinjlim\operatorname{Hom}_{A_0}(M_0,N_\lambda)
    \longrightarrow
    \operatorname{Hom}_{A_0}(M_0,\varinjlim N_\lambda),
    \tag{8.5.2.3}\label{IV.8.5.2.3-fr}
  \end{equation}
  qui, à tout système inductif d'homomorphismes de $A_0$-modules
  $\theta_\lambda:M_0\to N_\lambda$, fait correspondre sa limite inductive.

  Or, si $M_0$ est de type fini (resp. de présentation finie), on a une suite
  exacte $A_0^m\to M_0\to0$ (resp. $A_0^n\to A_0^m\to M_0\to0$)~; comme il est
  clair que (8.5.2.3) est bijectif lorsque $M_0$ est de la forme $A_0^q$, il
  suffit d'utiliser l'exactitude à gauche du foncteur
  $M_0\mapsto\operatorname{Hom}_{A_0}(M_0,P)$ et l'exactitude du foncteur
  $\varinjlim$ (dans la catégorie des groupes abéliens) pour conclure.

  Passons au cas où $S_0$ est quasi-compact, et soit $(U_i)$ un recouvrement
  fini de $S_0$ par des ouverts affines~; pour tout $\lambda$, les
  $U_{i\lambda}=u_{0\lambda}^{-1}(U_i)$ forment un recouvrement ouvert affine de
  $S_\lambda$, et les $V_i=u_0^{-1}(U_i)$ un recouvrement ouvert affine de $S$.
  Pour voir que $u_{\mathcal F,\mathcal G}^*$ est injectif, il faut prouver que
  si $f_\lambda:\mathcal F_\lambda\to\mathcal G_\lambda$ est tel que
  $f=u_\lambda^*(f_\lambda)=0$, alors il existe $\mu\geq\lambda$ tel que
  $f_\mu=u_{\lambda\mu}^*(f_\lambda)=0$. En vertu du lemme (8.5.2.1), pour
  chaque $i$ il existe un $\lambda_i$ tel que
  $f_\mu|U_{i\mu}=0$ pour $\mu\geq\lambda_i$. Il suffit donc de prendre $\mu$
  plus grand que tous les $\lambda_i$.

  Supposons en outre $S_0$ quasi-séparé et $\mathcal F_0$ de présentation finie,
  et soit $f:\mathcal F\to\mathcal G$ un homomorphisme de
  $\mathcal O_S$-Modules. En vertu du lemme (8.5.2.1), pour tout $i$, il existe
  un indice $\lambda_i$ et un homomorphisme
  $f_{\lambda_i}^{(i)}:\mathcal F_{\lambda_i}|U_{i,\lambda_i}
  \to\mathcal G_{\lambda_i}|U_{i,\lambda_i}$ tels que
  $u_{\lambda_i}^*(f_{\lambda_i}^{(i)})=f|V_i$. Comme $L$ est filtrant, on peut
  en outre supposer tous les $\lambda_i$ égaux à un même $\lambda$. Notons
  maintenant que $S_\lambda$ est quasi-séparé (1.2.3) et
  $\mathcal F_\lambda$ est un $\mathcal O_{S_\lambda}$-Module quasi-cohérent de
  présentation finie $(0_{\mathrm I},5.2.5)$~; comme, pour tout couple d'indices
  $i,j$, $U_{ij\lambda}=U_{i\lambda}\cap U_{j\lambda}$ est quasi-compact et que
  l'on a
  \[
    u_\lambda^*(f_\lambda^{(i)}|U_{ij\lambda})
    =u_\lambda^*(f_\lambda^{(j)}|U_{ij\lambda})
    =f|(V_i\cap V_j)
  \]
  par définition, il résulte de ce qu'on a vu plus haut qu'il existe un indice
  $\lambda_{ij}$ tel que
  \[
    u_{\lambda\mu}^*(f_\lambda^{(i)}|U_{ij\lambda})
    =u_{\lambda\mu}^*(f_\lambda^{(j)}|U_{ij\lambda})
    \quad\text{pour }\mu\geq\lambda_{ij}~;
  \]
  prenant $\mu$ plus grand que tous les $\lambda_{ij}$, on voit donc que
  $u_{\lambda\mu}^*(f_\lambda^{(i)})$ et
  $u_{\lambda\mu}^*(f_\lambda^{(j)})$ coïncident dans
  $U_{i\mu}\cap U_{j\mu}$ pour tout couple $(i,j)$, et par suite définissent
  un homomorphisme $f_\mu:\mathcal F_\mu\to\mathcal G_\mu$ tel que
  $f=u_\mu^*(f_\mu)$.
\end{enumerate}
\end{proof}

Avant de passer à la démonstration de (ii), notons les corollaires suivants de
(i)~:

\begin{corollary}[8.5.2.4]
\label{IV.8.5.2.4-fr}
Supposons $S_0$ quasi-compact, $\mathcal F_\lambda$ quasi-cohérent de type fini,
$\mathcal G_\lambda$ quasi-cohérent de présentation finie. Soit
$f_\lambda:\mathcal F_\lambda\to\mathcal G_\lambda$ un homomorphisme. Pour que
$f=u_\lambda^*(f_\lambda):\mathcal F\to\mathcal G$ soit un isomorphisme, il faut
et il suffit qu'il existe $\mu\geq\lambda$ tel que
$f_\mu=u_{\lambda\mu}^*(f_\lambda):\mathcal F_\mu\to\mathcal G_\mu$ soit un
isomorphisme.
\end{corollary}

\begin{proof}
On peut toujours supposer $S_\lambda=S_0$~; la question étant locale sur $S_0$,
on peut en outre ($S_0$ étant quasi-compact et $L$ filtrant) se ramener au cas
où $S_0$ est affine, donc quasi-séparé. La condition étant trivialement
suffisante, il reste à montrer qu'elle est nécessaire~: or, il y a par
hypothèse un $\mathcal O_S$-homomorphisme $g:\mathcal G\to\mathcal F$ tel que
$g\circ f=1_{\mathcal F}$ et $f\circ g=1_{\mathcal G}$. Comme $\mathcal G$ est
de présentation finie, il existe $\nu\geq\lambda$ et un homomorphisme
$g_\nu:\mathcal G_\nu\to\mathcal F_\nu$ tels que $g=u_\nu^*(g_\nu)$ en vertu de
(8.5.2, (i))~; on a par suite
$u_\nu^*(g_\nu\circ f_\nu)=1_{\mathcal F}$

\oldpage[IV-3]{22}
et $u_\nu^*(f_\nu\circ g_\nu)=1_{\mathcal G}$~; compte tenu de ce que
$\mathcal F_\nu$ et $\mathcal G_\nu$ sont de type fini, on en conclut par
(8.5.2, (i)) qu'il existe $\mu\geq\nu$ tel que l'on ait
$g_\mu\circ f_\mu=1_{\mathcal F_\mu}$ et
$f_\mu\circ g_\mu=1_{\mathcal G_\mu}$, d'où le corollaire.
\end{proof}

\begin{corollary}[8.5.2.5]
\label{IV.8.5.2.5-fr}
Supposons $S_0$ quasi-compact et quasi-séparé. Supposons que
$\mathcal F_\lambda$, $\mathcal G_\lambda$ soient des
$\mathcal O_{S_\lambda}$-Modules quasi-cohérents de présentation finie. Pour
que $\mathcal F$ et $\mathcal G$ soient isomorphes, il faut et il suffit qu'il
existe $\mu\geq\lambda$ tel que $\mathcal F_\mu$ et $\mathcal G_\mu$ soient
isomorphes. De plus, pour tout isomorphisme
$f:\mathcal F\xrightarrow{\sim}\mathcal G$, il existe un $\nu\geq\mu$ et un
isomorphisme $f_\nu:\mathcal F_\nu\xrightarrow{\sim}\mathcal G_\nu$ tels que
$f=u_\nu^*(f_\nu)$.
\end{corollary}

\begin{proof}
Cela résulte de (8.5.2.4) et de (8.5.2, (i)) puisque tout homomorphisme
$f:\mathcal F\to\mathcal G$ est de la forme $u_\mu^*(f_\mu)$ pour un
$\mu\geq\lambda$ et un homomorphisme
$f_\mu:\mathcal F_\mu\to\mathcal G_\mu$.
\end{proof}

\par\medskip\noindent\textup{(ii)} Considérons encore en premier lieu le cas où
$S_0=\operatorname{Spec}(A_0)$ est
affine. Alors l'assertion équivaut au lemme (5.13.7.1).

Dans le cas général, en partant d'un recouvrement ouvert affine fini $(U_i)$
de $S_0$, on déduit de (5.13.7.1) que pour tout $i$, il existe un indice
$\lambda(i)$ et un $\mathcal O_{U_{i,\lambda(i)}}$-Module quasi-cohérent de
présentation finie $\mathcal F_i$ tels que
$u_{\lambda(i)}^*(\mathcal F_i)=\mathcal F|V_i$ (avec les notations de (i)). En
outre, comme $L$ est filtrant, on peut supposer que tous les $\lambda_i$ sont
égaux à un même $\lambda$. Comme
$U_{ij\lambda}=U_{i\lambda}\cap U_{j\lambda}$ est quasi-compact et
quasi-séparé (1.2.7), il résulte de (8.5.2.5) que pour tout couple $(i,j)$, il
existe un indice $\lambda_{ij}\geq\lambda$ et un isomorphisme
\[
  \theta_{ij}:u_{\lambda,\lambda_{ij}}^*(\mathcal F_i|U_{ij\lambda})
  \xrightarrow{\sim}
  u_{\lambda,\lambda_{ij}}^*(\mathcal F_j|U_{ij\lambda})
\]
tel que $u_{\lambda_{ij}}^*(\theta_{ij})$ soit l'automorphisme identique de
$\mathcal F|(V_i\cap V_j)$~; on peut encore supposer tous les
$\lambda_{ij}$ égaux à $\lambda$. Changeant les notations, on peut donc
supposer qu'il existe pour tout couple $(i,j)$ un isomorphisme
\[
  \theta_{ij}:\mathcal F_i|U_{ij\lambda}
  \xrightarrow{\sim}\mathcal F_j|U_{ij\lambda},
\]
tels que $u_\lambda^*(\theta_{ij})$ soit l'automorphisme identique de
$\mathcal F|(V_i\cap V_j)$. Enfin, pour trois indices quelconques $i,j,k$, si
l'on pose
$U_{ijk,\lambda}=U_{i\lambda}\cap U_{j\lambda}\cap U_{k\lambda}$,
$U_{ijk,\lambda}$ est quasi-compact, et si $\theta'_{ij}$, $\theta'_{jk}$ et
$\theta'_{ik}$ désignent les restrictions de $\theta_{ij}$, $\theta_{jk}$ et
$\theta_{ik}$ à $U_{ijk,\lambda}$, on a
\[
  u_\lambda^*(\theta'_{ij}\circ\theta'_{jk})=u_\lambda^*(\theta'_{ik}).
\]
Il y a donc, en vertu de (i), un indice $\mu\geq\lambda$ tel que l'on ait
\[
  u_{\lambda\mu}^*(\theta'_{ik})
  =u_{\lambda\mu}^*(\theta'_{ij}\circ\theta'_{jk}),
\]
donc les isomorphismes
\[
  u_{\lambda\mu}^*(\theta'_{ij}):
  u_{\lambda\mu}^*(\mathcal F_i)|U_{ij\mu}
  \xrightarrow{\sim}
  u_{\lambda\mu}^*(\mathcal F_j)|U_{ij\mu}
\]
vérifient la condition de recollement, et définissent par suite sur $S_\mu$ un
$\mathcal O_{S_\mu}$-Module quasi-cohérent $\mathcal F_\mu$ de présentation
finie $(0_{\mathrm I},3.3.1)$ tel que $\mathcal F$ soit isomorphe à
$u_\mu^*(\mathcal F_\mu)$.

\begin{env}[8.5.3]
\label{IV.8.5.3-fr}
Le résultat de (8.5.2) s'exprime encore en disant que si $S_0$ est
quasi-compact et quasi-séparé, la catégorie des $\mathcal O_S$-Modules
quasi-cohérents de présentation finie est déterminée à équivalence près par la
donnée des catégories des $\mathcal O_{S_\lambda}$-Modules quasi-cohérents de
présentation finie, des foncteurs $u_{\lambda\mu}^*$ entre ces catégories, et
des isomorphismes de transition
$u_{\mu\nu}^*\circ u_{\lambda\mu}^*\xrightarrow{\sim}u_{\lambda\nu}^*$. De façon
imagée, on peut dire que la donnée d'un $\mathcal O_S$-Module quasi-cohérent de
présentation finie $\mathcal F$ équivaut «~fonctoriellement~» à la donnée d'un
$\mathcal O_{S_\lambda}$-Module de présentation finie $\mathcal F'_\lambda$
pour $\lambda$ grand et que, si un $\mathcal O_{S_\mu}$-Module quasi-cohérent de
présentation finie $\mathcal F''_\mu$ a aussi $\mathcal F$ comme image
réciproque, alors $\mathcal F'_\lambda$ et $\mathcal F''_\mu$ ont même image
réciproque dans un $S_\nu$ convenable ($\nu\geq\lambda$, $\nu\geq\mu$).

Nous allons interpréter de ce point de vue diverses notions liées aux
$\mathcal O_S$-Modules quasi-cohérents.
\end{env}

\oldpage[IV-3]{23}
\begin{corollary}[8.5.4]
\label{IV.8.5.4-fr}
Supposons $S_0$ quasi-compact et quasi-séparé~; alors, pour tout
$\mathcal O_{S_\lambda}$-Module quasi-cohérent $\mathcal G_\lambda$,
l'homomorphisme canonique (8.5.1.4) est bijectif.
\end{corollary}

\begin{proof}
En effet, il suffit d'appliquer (8.5.2, (i)) au cas où
$\mathcal F_\lambda=\mathcal O_{S_\lambda}$, qui est de présentation finie.
\end{proof}

\begin{proposition}[8.5.5]
\label{IV.8.5.5-fr}
Supposons $S_0$ quasi-compact, et supposons que $\mathcal F_\lambda$ soit un
$\mathcal O_{S_\lambda}$-Module quasi-cohérent de présentation finie. Pour que
$\mathcal F$ soit localement libre (resp. localement libre de rang $n$), il
faut et il suffit qu'il existe $\mu\geq\lambda$ tel que $\mathcal F_\mu$ le
soit.
\end{proposition}

\begin{proof}
La condition étant trivialement suffisante, prouvons qu'elle est nécessaire.
Si $\mathcal F$ est localement libre (resp. localement libre de rang $n$), il
existe un recouvrement ouvert affine fini $(V_i)$ de $S$ tel que
$\mathcal F|V_i$ soit isomorphe à $\mathcal O_S^{n_i}|V_i$ (resp.
$\mathcal O_S^n|V_i$) pour tout $i$. En vertu de (8.2.11), il existe un
$\nu\geq\lambda$ et pour chaque $i$ un ouvert quasi-compact $U_{i\nu}$ de
$S_\nu$ tels que $V_i=u_\nu^{-1}(U_{i\nu})$. Comme $S_\nu$ est quasi-compact,
chaque $U_{i\nu}$ est réunion finie d'ouverts affines~; on est donc ramené au
cas où $S_0$ est affine et $\mathcal F=\mathcal O_S^n$~; on sait alors qu'il
existe $\mu\geq\lambda$ tel que $\mathcal F_\mu$ soit isomorphe à
$\mathcal O_{S_\mu}^n$ (8.5.2.5).
\end{proof}

\begin{proposition}[8.5.6]
\label{IV.8.5.6-fr}
Supposons $S_0$ quasi-compact, et considérons une suite
\[
  \mathcal F_\lambda\longrightarrow\mathcal G_\lambda
  \longrightarrow\mathcal H_\lambda\longrightarrow0
\]
d'homomorphismes de $\mathcal O_{S_\lambda}$-Modules quasi-cohérents, où
$\mathcal F_\lambda$ et $\mathcal G_\lambda$ sont de type fini et
$\mathcal H_\lambda$ de présentation finie. Pour que la suite correspondante
$\mathcal F\to\mathcal G\to\mathcal H\to0$ soit exacte, il faut et il suffit
qu'il existe $\mu\geq\lambda$ tel que la suite
$\mathcal F_\mu\to\mathcal G_\mu\to\mathcal H_\mu\to0$ le soit (auquel cas il
en est de même de la suite
$\mathcal F_\nu\to\mathcal G_\nu\to\mathcal H_\nu\to0$ pour
$\nu\geq\mu$).
\end{proposition}

\begin{proof}
Le fait que la condition soit suffisante et la dernière assertion résultent
de ce que le foncteur $u_\lambda^*$ (resp. $u_{\mu\nu}^*$) est exact à droite.
Pour démontrer que la condition est nécessaire, notons qu'il résulte de
l'hypothèse et de (8.5.2, (i)) qu'il existe $\nu\geq\lambda$ tel que le composé
$\mathcal F_\nu\to\mathcal G_\nu\to\mathcal H_\nu$ soit nul. Si l'on pose
$\mathcal H'_\nu=\operatorname{Coker}(\mathcal F_\nu\to\mathcal G_\nu)$, on a
donc un homomorphisme $f_\nu:\mathcal H'_\nu\to\mathcal H_\nu$~; par
hypothèse, $u_\nu^*(f_\nu)$ est un isomorphisme, et il résulte donc de
(8.5.2.4) qu'il existe $\mu\geq\nu$ tel que $u_{\nu\mu}^*(f_\nu)$ soit un
isomorphisme, ce qui achève la démonstration.
\end{proof}

\begin{corollary}[8.5.7]
\label{IV.8.5.7-fr}
Supposons $S_0$ quasi-compact, $\mathcal F_\lambda$ quasi-cohérent,
$\mathcal G_\lambda$ quasi-cohérent de type fini, et soit
$f_\lambda:\mathcal F_\lambda\to\mathcal G_\lambda$ un homomorphisme. Pour que
$f=u_\lambda^*(f_\lambda):\mathcal F\to\mathcal G$ soit surjectif, il faut et
il suffit qu'il existe $\mu\geq\lambda$ tel que
$f_\mu=u_{\lambda\mu}^*(f_\lambda):\mathcal F_\mu\to\mathcal G_\mu$ le soit.
\end{corollary}

\begin{proof}
C'est le cas particulier de (8.5.6) appliqué à la suite
$0\to\mathcal H_\lambda\to0\to0$, où
$\mathcal H_\lambda=\operatorname{Coker}(f_\lambda)$, qui est quasi-cohérent et
de type fini (compte tenu de ce que l'on a
$\mathcal H=\operatorname{Coker}(f)$ et
$\mathcal H_\mu=\operatorname{Coker}(f_\mu)$, en vertu de l'exactitude à droite
des foncteurs $u_\lambda^*$ et $u_{\lambda\mu}^*$).
\end{proof}

\begin{corollary}[8.5.8]
\label{IV.8.5.8-fr}
Supposons $S_0$ quasi-compact et les morphismes
$u_{\lambda\mu}:S_\mu\to S_\lambda$ plats. Alors~:
\begin{enumerate}[label=(\roman*)]
  \item Soit
  $\mathcal F_\lambda\xrightarrow{f_\lambda}\mathcal G_\lambda
  \xrightarrow{g_\lambda}\mathcal H_\lambda$ une suite d'homomorphismes de
  $\mathcal O_{S_\lambda}$-Modules quasi-cohérents, tels que
  $\operatorname{Im}f_\lambda$ et $\operatorname{Ker}g_\lambda$ soient de type
  fini. Pour que la suite correspondante
  $\mathcal F\xrightarrow{f}\mathcal G\xrightarrow{g}\mathcal H$ soit

  \oldpage[IV-3]{24}
  exacte, il faut et il suffit qu'il existe $\mu\geq\lambda$ tel que la suite
  $\mathcal F_\mu\xrightarrow{f_\mu}\mathcal G_\mu
  \xrightarrow{g_\mu}\mathcal H_\mu$ soit exacte.

  \item Soit $f_\lambda:\mathcal F_\lambda\to\mathcal G_\lambda$ un
  homomorphisme de $\mathcal O_{S_\lambda}$-Modules quasi-cohérents tel que
  $\operatorname{Ker}f_\lambda$ soit de type fini. Pour que
  $f=u_\lambda^*(f_\lambda):\mathcal F\to\mathcal G$ soit injectif, il faut et
  il suffit qu'il existe $\mu\geq\lambda$ tel que
  $f_\mu=u_{\lambda\mu}^*(f_\lambda):\mathcal F_\mu\to\mathcal G_\mu$ le soit.
\end{enumerate}
\end{corollary}

\begin{proof}
\begin{enumerate}[label=(\roman*)]
  \item Compte tenu de (8.3.8, (ii)), notons que, par platitude,
  $\operatorname{Im}f$ et $\operatorname{Ker}g$ (resp.
  $\operatorname{Im}f_\mu$ et $\operatorname{Ker}g_\mu$ pour
  $\mu\geq\lambda$) sont les images réciproques de
  $\operatorname{Im}f_\lambda$ et $\operatorname{Ker}g_\lambda$
  $(0_{\mathrm I},6.7.2)$. Supposons que la suite
  $\mathcal F\xrightarrow{f}\mathcal G\xrightarrow{g}\mathcal H$ soit exacte.
  Puisque $\operatorname{Im}f_\lambda$ est de type fini, il existe
  $\mu\geq\lambda$ tel que le composé
  $\mathcal F_\mu\xrightarrow{f_\mu}\mathcal G_\mu
  \xrightarrow{g_\mu}\mathcal H_\mu$ soit nul, en vertu de (8.5.2, (i)).
  Changeant les notations, on peut donc déjà supposer que
  $g_\lambda\circ f_\lambda=0$. Alors comme l'homomorphisme
  $\mathcal F\to\operatorname{Ker}g$ est surjectif et que
  $\operatorname{Ker}g_\lambda$ est de type fini, il résulte de (8.5.7) qu'il
  existe $\mu\geq\lambda$ tel que l'homomorphisme
  $\mathcal F_\mu\to\operatorname{Ker}g_\mu$ soit surjectif, ce qui achève de
  prouver (i).

  \item L'assertion est le cas particulier de (i) appliqué à la suite
  $0\to\mathcal F_\lambda\to\mathcal G_\lambda$.
\end{enumerate}
\end{proof}

\begin{lemma}[8.5.9]
\label{IV.8.5.9-fr}
Supposons $S_0$ quasi-compact, $\mathcal F_\lambda$ quasi-cohérent de type
fini~; soient $\mathcal G'_\lambda$ et $\mathcal G''_\lambda$ deux quotients
quasi-cohérents de $\mathcal F_\lambda$, $\mathcal G'_\lambda$ étant en outre
supposé de présentation finie. Pour que $\mathcal G''$ soit quotient de
$\mathcal G'$, il faut et il suffit qu'il existe $\mu\geq\lambda$ tel que
$\mathcal G''_\mu$ soit quotient de $\mathcal G'_\mu$.
\end{lemma}

\begin{proof}
Par hypothèse, il y a deux homomorphismes surjectifs
$p'_\lambda:\mathcal F_\lambda\to\mathcal G'_\lambda$,
$p''_\lambda:\mathcal F_\lambda\to\mathcal G''_\lambda$~; en vertu de
l'exactitude à droite de $u_\lambda^*$ et de $u_{\lambda\mu}^*$,
$p'=u_\lambda^*(p'_\lambda)$, $p''=u_\lambda^*(p''_\lambda)$,
$p'_\mu=u_{\lambda\mu}^*(p'_\lambda)$,
$p''_\mu=u_{\lambda\mu}^*(p''_\lambda)$ sont aussi surjectifs~; en outre, s'il
existe un homomorphisme $f:\mathcal G'\to\mathcal G''$ (resp.
$f_\mu:\mathcal G'_\mu\to\mathcal G''_\mu$) tel que $p''=f\circ p'$ (resp.
$p''_\mu=f_\mu\circ p'_\mu$), cet homomorphisme est nécessairement unique, ce
qui montre que la question est locale sur $S_\lambda$, et qu'on peut par suite
($S_0$ étant quasi-compact et $L$ filtrant) supposer $S_\lambda$ affine, donc
quasi-séparé. Il est clair que la condition de l'énoncé est suffisante.
Inversement, comme $\mathcal G'_\lambda$ est de présentation finie,
$S_\lambda$ quasi-compact et quasi-séparé, il résulte de (8.5.2, (i)) que s'il
existe un homomorphisme $f:\mathcal G'\to\mathcal G''$ tel que
$p''=f\circ p'$, il existe un $\mu\geq\lambda$ et un homomorphisme
$f_\mu:\mathcal G'_\mu\to\mathcal G''_\mu$ tels que
$f=u_\mu^*(f_\mu)$ et $p''_\mu=f_\mu\circ p'_\mu$, d'où le lemme.
\end{proof}

\begin{env}[8.5.10]
\label{IV.8.5.10-fr}
Dans la fin de ce numéro, pour tout Module quasi-cohérent $\mathcal F$ sur un
préschéma, notons $\mathfrak Q(\mathcal F)$ l'ensemble des Modules quotients de
$\mathcal F$ qui sont de présentation finie. Si $\mathcal F_\lambda$ est
quasi-cohérent et $\mathcal G_\lambda\in\mathfrak Q(\mathcal F_\lambda)$, il
résulte du fait que $u_{\lambda\mu}^*$ et $u_\lambda^*$ sont exacts à droite,
que l'on a $\mathcal G_\mu\in\mathfrak Q(\mathcal F_\mu)$ pour
$\mu\geq\lambda$ et $\mathcal G\in\mathfrak Q(\mathcal F)$~; il est clair que
$(\mathfrak Q(\mathcal F_\lambda),u_{\lambda\mu}^*)$ est un système inductif
d'ensembles, et que les
$u_\lambda^*:\mathfrak Q(\mathcal F_\lambda)\to\mathfrak Q(\mathcal F)$
forment un système inductif d'applications, d'où, en passant à la limite
inductive, une application canonique
\begin{equation}
  u_{\mathfrak Q}:\varinjlim\mathfrak Q(\mathcal F_\lambda)
  \longrightarrow\mathfrak Q(\mathcal F).
  \tag{8.5.10.1}\label{IV.8.5.10.1-fr}
\end{equation}

En outre, si $(\mathcal F'_\lambda)$ est une seconde famille de
$\mathcal O_{S_\lambda}$-Modules quasi-cohérents et si, pour tout $\lambda$,
$\mathcal F'_\lambda$ est un quotient de $\mathcal F_\lambda$, alors
$\mathcal F'$ est un quotient de $\mathcal F$ et l'on a un diagramme
commutatif

\oldpage[IV-3]{25}
\[
\begin{tikzcd}
\varinjlim\mathfrak Q(\mathcal F_\lambda)
  \ar[r,"u_{\mathfrak Q}"] \ar[d]
  & \mathfrak Q(\mathcal F) \ar[d] \\
\varinjlim\mathfrak Q(\mathcal F'_\lambda)
  \ar[r,"u_{\mathfrak Q}"']
  & \mathfrak Q(\mathcal F')
\end{tikzcd}
\tag{8.5.10.2}\label{IV.8.5.10.2-fr}
\]
\end{env}

\begin{proposition}[8.5.11]
\label{IV.8.5.11-fr}
Supposons $S_0$ quasi-compact (resp. quasi-compact et quasi-séparé). Supposons
$\mathcal F_\lambda$ quasi-cohérent de type fini (resp. de présentation finie)
pour tout $\lambda$~; alors l'application canonique (8.5.10.1) est injective
(resp. bijective).
\end{proposition}

\begin{proof}
La première assertion résulte du lemme plus précis (8.5.9). Pour démontrer la
seconde, considérons un $\mathcal O_S$-Module quotient $\mathcal G$ de
$\mathcal F$ qui soit de présentation finie. Il résulte de (8.5.2, (ii)) qu'il
existe un $\lambda'\in L$ et un $\mathcal O_{S_{\lambda'}}$-Module
$\mathcal G_{\lambda'}$ quasi-cohérent de présentation finie tels que
$\mathcal G=u_{\lambda'}^*(\mathcal G_{\lambda'})$~; comme $L$ est filtrant, on
peut supposer $\lambda'=\lambda$ (en remplaçant $\lambda$ et $\lambda'$ par un
indice qui les majore). Considérons alors l'homomorphisme canonique
$p:\mathcal F\to\mathcal G$~; il résulte de (8.5.2, (i)) qu'il existe un
$\mu\geq\lambda$ et un homomorphisme
$p_\mu:\mathcal F_\mu\to\mathcal G_\mu$ tels que $p=u_\mu^*(p_\mu)$. En outre,
en vertu de (8.5.7), on peut supposer $\mu$ pris assez grand pour que $p_\mu$
soit surjectif, ce qui termine la démonstration.
\end{proof}

\subsection{Sous-préschémas de présentation finie d'une limite projective de préschémas}
\label{subsection:IV.8.6-fr}

\begin{env}[8.6.1]
\label{IV.8.6.1-fr}
Étant donné un préschéma $Y$, désignons dans ce numéro par
$\mathfrak{Spr}(Y)$ l'ensemble ordonné (I, 4.1.10) des sous-préschémas de $Y$
qui sont de présentation finie sur $Y$ (1.6.1), par $\mathfrak{Spro}(Y)$
(resp. $\mathfrak{Sprf}(Y)$) la partie de $\mathfrak{Spr}(Y)$ formée des
sous-préschémas induits sur des ouverts (resp. des sous-préschémas fermés) de
$Y$, de présentation finie sur $Y$~; il revient au même de dire qu'un
sous-préschéma $Z$ de $Y$ appartient à $\mathfrak{Spro}(Y)$ (resp.
$\mathfrak{Sprf}(Y)$) ou qu'il est induit sur un ouvert et que l'espace
sous-jacent $Z$ est rétrocompact dans $Y$ (resp. qu'il est fermé et que l'idéal
$\mathcal J$ de $\mathcal O_Y$ qui le définit est de type fini, ce qui signifie
aussi que $j_*(\mathcal O_Z)\in\mathfrak Q(\mathcal O_Y)$ si $j:Z\to Y$ est
l'injection canonique) (1.6.1 et 1.4.5).
\end{env}

\begin{env}[8.6.2]
\label{IV.8.6.2-fr}
Nous conservons toujours les notations de (8.2.2) et supposons que $S_0$ est
un des $S_\lambda$. Soit $Y_\lambda$ un sous-préschéma de $S_\lambda$~; alors
$Y_\mu=u_{\lambda\mu}^{-1}(Y_\lambda)$ (resp.
$Y=u_\lambda^{-1}(Y_\lambda)$) est un sous-préschéma de $S_\mu$ pour
$\mu\geq\lambda$ (resp. de $S$)~; il est induit sur un ouvert (resp. il est
fermé) si $Y_\lambda$ l'est (I, 4.3.2) et de présentation finie sur $S_\mu$
(resp. $S$) si $Y_\lambda$ est de présentation finie sur $S_\lambda$
(1.6.2, (iii)). Par suite
$(\mathfrak{Spr}(S_\lambda),u_{\lambda\mu}^{-1})$ (resp.
$(\mathfrak{Spro}(S_\lambda),u_{\lambda\mu}^{-1})$,
$(\mathfrak{Sprf}(S_\lambda),u_{\lambda\mu}^{-1})$) est un système inductif et
les applications
$u_\lambda^{-1}:\mathfrak{Spr}(S_\lambda)\to\mathfrak{Spr}(S)$ (resp.
$\mathfrak{Spro}(S_\lambda)\to\mathfrak{Spro}(S)$,
$\mathfrak{Sprf}(S_\lambda)\to\mathfrak{Sprf}(S)$) forment un système inductif
d'applications~; d'où, en passant à la limite inductive, des applications
canoniques
\begin{align}
  \varinjlim\mathfrak{Spr}(S_\lambda)
    &\longrightarrow\mathfrak{Spr}(S)
    \tag{8.6.2.1}\label{IV.8.6.2.1-fr}\\
  \varinjlim\mathfrak{Spro}(S_\lambda)
    &\longrightarrow\mathfrak{Spro}(S)
    \tag{8.6.2.2}\label{IV.8.6.2.2-fr}\\
  \varinjlim\mathfrak{Sprf}(S_\lambda)
    &\longrightarrow\mathfrak{Sprf}(S).
    \tag{8.6.2.3}\label{IV.8.6.2.3-fr}
\end{align}
\end{env}

\oldpage[IV-3]{26}

Rappelons (I, 4.1.10) que $\mathfrak{Spr}(X)$, pour tout préschéma $X$, est un
ensemble \emph{ordonné} par la relation «~$Z$ est un sous-préschéma du
sous-préschéma $Y$~» qui s'écrit $Z\leq Y$. Les applications
$u_{\lambda\mu}^{-1}$ et $u_\lambda^{-1}$ sont \emph{croissantes} pour les
relations d'ordre correspondantes dans $\mathfrak{Spr}(S_\lambda)$,
$\mathfrak{Spr}(S_\mu)$, $\mathfrak{Spr}(S)$. En outre, on définit une relation
d'ordre dans l'ensemble $\varinjlim\mathfrak{Spr}(S_\lambda)$ en écrivant que
$\eta\leq\eta'$ pour deux éléments de cet ensemble lorsqu'il existe un
$\lambda$ et deux éléments $Y_\lambda$, $Y'_\lambda$ de
$\mathfrak{Spr}(S_\lambda)$, dont $\eta$ et $\eta'$ sont les images canoniques,
et qui sont tels que $Y_\lambda\leq Y'_\lambda$~; on vérifie aisément que cela
ne dépend pas des représentants $Y_\lambda$, $Y'_\lambda$ considérés, et que
l'on a bien ainsi une \emph{relation d'ordre}. Cela étant, le fait que les
$u_\lambda^{-1}$ sont croissantes entraîne aussitôt que l'application
canonique (8.6.2.1) est \emph{croissante}~; il en est de même évidemment de
(8.6.2.2) et (8.6.2.3).

\begin{proposition}[8.6.3]
\label{IV.8.6.3-fr}
Supposons $S_0$ quasi-compact (resp. quasi-compact et quasi-séparé). Alors les
applications (8.6.2.1), (8.6.2.2), (8.6.2.3) sont injectives (resp.
bijectives).
\end{proposition}

\begin{proof}
Compte tenu des remarques de (8.6.1), les assertions relatives à (8.6.2.3)
découlent de (8.5.11) appliqué à $\mathcal F_\lambda=\mathcal O_{S_\lambda}$~;
de même, les assertions relatives à (8.6.2.2) sont des cas particuliers de
(8.3.5) et (8.3.11), compte tenu de ce que les $S_\lambda$ et $S$ sont
quasi-compacts. Reste à considérer l'application (8.6.2.1). Prouvons d'abord
qu'elle est surjective lorsque $S_0$ est quasi-compact et quasi-séparé. Soit
$Z$ un sous-préschéma de $S$, de présentation finie sur $S$~; comme $S$ est
quasi-compact, il en est de même de $Z$, donc il existe un ouvert
quasi-compact $U$ de $S$ tel que $Z$ soit un sous-préschéma fermé de $U$, de
présentation finie sur $U$. Il existe alors un indice $\lambda$ et un ouvert
quasi-compact $U_\lambda$ de $S_\lambda$ tels que
$U=u_\lambda^{-1}(U_\lambda)$ (8.2.11)~; comme $S_\lambda$ est quasi-séparé, il
en est de même de $U_\lambda$ (I, 1.2.7), et par suite on peut se borner au cas
où $U=S$~; mais dans ce cas, on est ramené au fait que (8.6.2.3) est
surjective.

Enfin, pour voir que (8.6.2.1) est injective lorsque $S_0$ est quasi-compact,
il suffira de prouver le résultat plus précis suivant~:
\end{proof}

\begin{corollary}[8.6.3.1]
\label{IV.8.6.3.1-fr}
Supposons $S_0$ quasi-compact et soient $Z'_\lambda$, $Z''_\lambda$ deux
sous-préschémas de $S_\lambda$, de présentation finie sur $S_\lambda$. Pour
que $Z'=u_\lambda^{-1}(Z'_\lambda)$ soit majoré par
$Z''=u_\lambda^{-1}(Z''_\lambda)$ (I, 4.1.10), il faut et il suffit qu'il
existe $\mu\geq\lambda$ tel que
$Z'_\mu=u_{\lambda\mu}^{-1}(Z'_\lambda)$ soit majoré par
$Z''_\mu=u_{\lambda\mu}^{-1}(Z''_\lambda)$.
\end{corollary}

\begin{proof}
Il est trivial que la condition est suffisante. Pour voir qu'elle est
nécessaire, notons d'abord que les ensembles sous-jacents $Z'_\lambda$ et
$Z''_\lambda$ sont localement constructibles dans $S_\lambda$ par hypothèse
(I, 1.8.4), donc l'hypothèse entraîne l'existence d'un $\nu\geq\lambda$ tel
que $Z'_\nu\subset Z''_\nu$ (8.3.5)~; remplaçant $\lambda$ par $\nu$, on peut
donc déjà supposer que l'on a $Z'_\lambda\subset Z''_\lambda$. En outre, par
hypothèse (I, 1.6.1), les sous-espaces $Z'_\lambda$ et $Z''_\lambda$ de
$S_\lambda$ sont quasi-compacts~; pour tout point $x\in Z'_\lambda$, il y a un
voisinage ouvert quasi-compact $V(x)$ dans $S_\lambda$ tels que
$V(x)\cap Z'_\lambda$ et $V(x)\cap Z''_\lambda$ soient fermés dans $V(x)$. En
recouvrant $S_\lambda$ par un nombre fini de voisinages $V(x_i)$ on voit donc
qu'il y a un ouvert quasi-compact $U_\lambda$ de $S_\lambda$ contenant
$Z'_\lambda$ et tel que $Z'_\lambda$ et $Z''_\lambda\cap U_\lambda$ soient
fermés dans $U_\lambda$. Si l'on désigne par $Y_\lambda$ le sous-préschéma
induit par $Z''_\lambda$ sur $U_\lambda\cap Z''_\lambda$, il est clair qu'avec
les notations habituelles, $Y_\mu$ (resp. $Y$) est induit par $Z''_\mu$ sur
$U_\mu\cap Z''_\mu$ pour $\mu\geq\lambda$ (resp. par $Z''$ sur
$U\cap Z''$)~; en outre $Z'$ est majoré par $Y$ (I, 4.4.1), et comme il suffit
de prouver que $Z'_\mu$ est majoré par $Y_\mu$ pour $\mu$
\oldpage[IV-3]{27}
assez grand, on voit finalement qu'on est ramené (en remplaçant $S_\lambda$
par $U_\lambda$) au cas où $Z'_\lambda$ et $Z''_\lambda$ sont des
sous-préschémas fermés de $S_\lambda$. Mais alors, cela a déjà été démontré
puisque (8.6.2.3) est croissante et injective.
\end{proof}

\begin{corollary}[8.6.4]
\label{IV.8.6.4-fr}
Supposons $S_0$ quasi-compact, et soit $Z_\lambda$ un sous-préschéma de
$S_\lambda$, de présentation finie sur $S_\lambda$. Pour que
$Z=u_\lambda^{-1}(Z_\lambda)$ soit un sous-préschéma induit sur un ouvert
(resp. un sous-préschéma fermé) de $S$, il faut et il suffit qu'il existe
$\mu\geq\lambda$ tel que $Z_\mu=u_{\lambda\mu}^{-1}(Z_\lambda)$ soit induit
sur un ouvert (resp. un sous-préschéma fermé) dans $S_\mu$.
\end{corollary}

\begin{proof}
Soit $(U_\lambda^{(i)})$ un recouvrement ouvert affine fini de $S_\lambda$, et
posons $U_\mu^{(i)}=u_{\lambda\mu}^{-1}(U_\lambda^{(i)})$ pour
$\mu\geq\lambda$ et $U^{(i)}=u_\lambda^{-1}(U_\lambda^{(i)})$. Si $Z$ est
ouvert (resp. fermé) dans $S$, $Z\cap U^{(i)}$ l'est dans $U^{(i)}$, et
inversement si chacun des $Z_\mu\cap U_\mu^{(i)}$ est ouvert (resp. fermé)
dans $U_\mu^{(i)}$, $Z_\mu$ l'est dans $S_\mu$. Puisque $L$ est filtrant, il
suffit de démontrer le corollaire lorsque $S_\lambda$ est affine, donc
quasi-séparé. Mais alors le résultat découle de ce que les applications
(8.6.2.1), (8.6.2.2) et (8.6.2.3) sont bijectives.
\end{proof}

\subsection{Critères pour qu'une limite projective de préschémas soit un préschéma réduit (resp. intègre)}
\label{subsection:IV.8.7-fr}

Nous conservons les hypothèses et notations de (8.2.2) et supposons toujours
que $S_0$ est un des $S_\lambda$.

\begin{proposition}[8.7.1]
\label{IV.8.7.1-fr}
Supposons que $S$ soit non réduit. Alors il existe $\lambda_0$ tel que pour
$\lambda\geq\lambda_0$, $S_\lambda$ soit non réduit.
\end{proposition}

\begin{proof}
La question étant locale sur $S_0$, on peut supposer
$S_0=\operatorname{Spec}(A_0)$ affine, d'où $S=\operatorname{Spec}(A)$, où
$A=\varinjlim A_\lambda$ est la limite inductive d'un système inductif de
$A_0$-algèbres $(A_\lambda)$. On sait alors (5.13.2) que le nilradical de $A$
est la limite inductive de ceux des $A_\lambda$~; s'il n'est pas nul, un des
$A_\lambda$ contient donc un élément nilpotent $a_\lambda\neq0$ dont l'image
dans $A$ est un élément nilpotent et $\neq0$, et l'image de $a_\lambda$ dans
les $A_\mu$ pour $\mu\geq\lambda$ est par suite un élément nilpotent et
$\neq0$.
\end{proof}

\begin{proposition}[8.7.2]
\label{IV.8.7.2-fr}
Supposons vérifiée l'une des hypothèses suivantes~:
\begin{enumerate}
\item[\emph{a)}] $S_0$ est quasi-compact, le nilradical $\mathcal N_0$ de
$\mathcal O_{S_0}$ est un Idéal de type fini (ce qui sera par exemple vérifié
lorsque $S_0$ est noethérien), et les morphismes
$u_{\lambda\mu}:S_\mu\to S_\lambda$ sont des immersions ouvertes.
\item[\emph{b)}] Les morphismes $u_{\lambda\mu}:S_\mu\to S_\lambda$ sont
fidèlement plats.
\end{enumerate}
Sous ces conditions, pour que $S$ soit réduit, il faut et il suffit qu'il
existe $\lambda_0$ tel que $S_\lambda$ soit réduit pour
$\lambda\geq\lambda_0$.

En outre, dans le cas \emph{b)}, si $S$ est réduit, tous les $S_\lambda$ le
sont.
\end{proposition}

\begin{proof}
La dernière assertion résulte de ce que le morphisme $S\to S_\lambda$ est
alors fidèlement plat pour tout $\lambda$ (8.3.8) et de (2.1.13). D'autre part,
(8.7.1) prouve que la condition de l'énoncé est suffisante (sans hypothèse sur
$S_0$ ni sur les $u_{\lambda\mu}$). Il reste donc à voir que la condition est
nécessaire dans l'hypothèse \emph{a)}~; alors (8.2.13), l'espace sous-jacent à
$S$ s'identifie à l'intersection des espaces sous-jacents aux $S_\lambda$ (les
$S_\lambda$ étant identifiés à des sous-préschémas induits sur des ouverts de
$S_0$), et le faisceau structural $\mathcal O_S$ s'identifie au
\oldpage[IV-3]{28}
faisceau induit sur $S$ par tous les $\mathcal O_{S_\lambda}$~; en particulier
pour tout $s\in S$, l'anneau local $\mathcal O_s$ est \emph{le même} pour $S$
et pour tous les $S_\lambda$. Si $\mathcal N_\lambda$ est le Nilradical de
$\mathcal O_{S_\lambda}$, le Nilradical $\mathcal N$ de $\mathcal O_S$ a donc
en chaque point de $S$ la même fibre $\mathcal N_s$ (nilradical de
$\mathcal O_s$) que $u_\lambda^*(\mathcal N_\lambda)$ (induit sur $S$ par
$\mathcal N_\lambda$). L'hypothèse que $S$ est réduit entraîne par suite
$u_\lambda^*(\mathcal N_\lambda)=0$~; comme $\mathcal N_0$ est supposé de
type fini, il en est de même des $\mathcal N_\lambda$, et la conclusion résulte
donc de (8.5.7).
\end{proof}

\begin{corollary}[8.7.3]
\label{IV.8.7.3-fr}
Supposons vérifiée l'une des hypothèses suivantes~:
\begin{enumerate}
\item[\emph{a)}] $S_0$ est un préschéma noethérien et les morphismes
$u_{\lambda\mu}:S_\mu\to S_\lambda$ sont des immersions ouvertes.
\item[\emph{b)}] Les morphismes $u_{\lambda\mu}:S_\mu\to S_\lambda$ sont
fidèlement plats.
\end{enumerate}
Alors, pour que $S$ soit intègre, il faut et il suffit qu'il existe
$\lambda_0$ tel que $S_\lambda$ soit intègre pour $\lambda\geq\lambda_0$.
\end{corollary}

\begin{proof}
Dire qu'un préschéma est intègre signifie qu'il est à la fois réduit et
irréductible~; le corollaire résulte donc de (8.7.2) et de (8.4.3).
\end{proof}

\begin{remark}[8.7.4]
\label{IV.8.7.4-fr}
Si l'on ne fait pas d'hypothèse sur les $u_{\lambda\mu}$, il peut se faire que
$S$ soit intègre bien que tous les $S_\lambda$ soient non réduits et non
connexes, comme le montre l'exemple (8.4.4), où on prend l'anneau $A$ non
réduit.
\end{remark}

\subsection{Préschémas de présentation finie sur une limite projective de préschémas}
\label{subsection:IV.8.8-fr}

\begin{env}[8.8.1]
\label{IV.8.8.1-fr}
Conservant toujours les notations et hypothèses de (8.2.2), nous supposerons
donnés dans cette section deux $S_\alpha$-préschémas $X_\alpha$, $Y_\alpha$,
ce qui définit (8.2.5) deux systèmes projectifs de préschémas
$(X_\lambda,v_{\lambda\mu})$ et $(Y_\lambda,w_{\lambda\mu})$ en posant
$X_\lambda=X_\alpha\times_{S_\alpha}S_\lambda$,
$Y_\lambda=Y_\alpha\times_{S_\alpha}S_\lambda$,
$v_{\lambda\mu}=1_{X_\alpha}\times u_{\lambda\mu}$,
$w_{\lambda\mu}=1_{Y_\alpha}\times u_{\lambda\mu}$ (pour
$\alpha\leq\lambda\leq\mu$), dont les limites projectives sont respectivement
$X=X_\alpha\times_{S_\alpha}S$, $Y=Y_\alpha\times_{S_\alpha}S$, les
morphismes canoniques $v_\lambda:X\to X_\lambda$ et
$w_\lambda:Y\to Y_\lambda$ étant respectivement égaux à
$1_{X_\alpha}\times u_\lambda$ et $1_{Y_\alpha}\times u_\lambda$. Pour
$\alpha\leq\lambda\leq\mu$, on a une application canonique
$e_{\mu\lambda}:\operatorname{Hom}_{S_\lambda}(X_\lambda,Y_\lambda)\to
\operatorname{Hom}_{S_\mu}(X_\mu,Y_\mu)$, qui à tout $S_\lambda$-morphisme
$f_\lambda:X_\lambda\to Y_\lambda$ fait correspondre
$f_\mu=f_\lambda\times1_{S_\mu}:X_\lambda\times_{S_\lambda}S_\mu\to
Y_\lambda\times_{S_\lambda}S_\mu$, et il est clair que
$(\operatorname{Hom}_{S_\lambda}(X_\lambda,Y_\lambda),e_{\mu\lambda})$ est un
\emph{système inductif} d'ensembles. De même, on a une application canonique
$e_\lambda:\operatorname{Hom}_{S_\lambda}(X_\lambda,Y_\lambda)\to
\operatorname{Hom}_S(X,Y)$ qui à $f_\lambda$ fait correspondre
$f=f_\lambda\times1_S:X_\lambda\times_{S_\lambda}S\to
Y_\lambda\times_{S_\lambda}S$ et $(e_\lambda)$ est un système inductif
d'applications~; d'où, en passant à la limite inductive, une application
canonique, fonctorielle en $S_\alpha$, $X_\alpha$ et $Y_\alpha$~:
\begin{equation}
e:\varinjlim\operatorname{Hom}_{S_\lambda}(X_\lambda,Y_\lambda)
\longrightarrow\operatorname{Hom}_S(X,Y).
\tag{8.8.1.1}\label{IV.8.8.1.1-fr}
\end{equation}
\end{env}

\begin{theorem}[8.8.2]
\label{IV.8.8.2-fr}
\begin{enumerate}
\item[(i)] Supposons $X_\alpha$ quasi-compact (resp. quasi-compact et
quasi-séparé), et $Y_\alpha$ localement de type fini (resp. localement de
présentation finie) sur $S_\alpha$. Alors l'application (8.8.1.1) est
injective (resp. bijective).
\item[(ii)] Supposons $S_0$ quasi-compact et quasi-séparé. Pour tout préschéma
$X$ de présentation finie sur $S$, il existe un $\lambda\in L$, un préschéma
$X_\lambda$ de présentation finie sur $S_\lambda$, et un $S$-isomorphisme
$X\xrightarrow{\sim}X_\lambda\times_{S_\lambda}S$.
\end{enumerate}
\end{theorem}

\begin{proof}
\oldpage[IV-3]{29}
(i) Considérons d'abord le cas où $S_0=\operatorname{Spec}(A_0)$,
$X_\alpha=\operatorname{Spec}(B_\alpha)$ et
$Y_\alpha=\operatorname{Spec}(C_\alpha)$ sont affines~; alors les
$S_\lambda=\operatorname{Spec}(A_\lambda)$ et $S=\operatorname{Spec}(A)$ sont
aussi affines, avec $A=\varinjlim A_\lambda$, et les assertions de (i) sont
équivalentes au
\end{proof}

\begin{lemma}[8.8.2.1]
\label{IV.8.8.2.1-fr}
Soient $A_0$ un anneau, $(A_\lambda)_{\lambda\in L}$ un système inductif de
$A_0$-algèbres, $A=\varinjlim A_\lambda$~; soient $B_\alpha$ une
$A_\alpha$-algèbre, $C_\alpha$ une $A_\alpha$-algèbre de type fini (resp. de
présentation finie). Alors l'homomorphisme canonique
\begin{equation}
\varinjlim_\lambda
\operatorname{Hom}_{A_\lambda\text{-alg.}}
  (C_\alpha\otimes_{A_\alpha}A_\lambda,
   B_\alpha\otimes_{A_\alpha}A_\lambda)
\longrightarrow
\operatorname{Hom}_{A\text{-alg.}}
  (C_\alpha\otimes_{A_\alpha}A,
   B_\alpha\otimes_{A_\alpha}A)
\tag{8.8.2.2}\label{IV.8.8.2.2-fr}
\end{equation}
est injectif (resp. bijectif).
\end{lemma}

\begin{proof}
On sait que l'on a des isomorphismes canoniques fonctoriels
\[
\operatorname{Hom}_{A_\lambda\text{-alg.}}
 (C_\alpha\otimes_{A_\alpha}A_\lambda,
  B_\alpha\otimes_{A_\alpha}A_\lambda)
\xrightarrow{\sim}
\operatorname{Hom}_{A_\alpha\text{-alg.}}
 (C_\alpha,B_\alpha\otimes_{A_\alpha}A_\lambda)
\]
\[
\operatorname{Hom}_{A\text{-alg.}}
 (C_\alpha\otimes_{A_\alpha}A,
  B_\alpha\otimes_{A_\alpha}A)
\xrightarrow{\sim}
\operatorname{Hom}_{A_\alpha\text{-alg.}}
 (C_\alpha,B_\alpha\otimes_{A_\alpha}A)
\]
en vertu de la propriété universelle du produit tensoriel de deux algèbres. Il
suffit donc de prouver le
\end{proof}

\begin{lemma}[8.8.2.3]
\label{IV.8.8.2.3-fr}
Soient $E$ un anneau, $G$ une $E$-algèbre de type fini (resp. de présentation
finie), $(F_\lambda)$ un système inductif de $E$-algèbres. Alors
l'homomorphisme canonique
\[
\varinjlim\operatorname{Hom}_{E\text{-alg.}}(G,F_\lambda)
\longrightarrow
\operatorname{Hom}_{E\text{-alg.}}(G,\varinjlim F_\lambda)
\]
qui, à tout système inductif d'homomorphismes
$\theta_\lambda:G\to F_\lambda$ de $E$-algèbres, fait correspondre sa limite
inductive, est injectif (resp. bijectif).
\end{lemma}

\begin{proof}
Supposons d'abord la $E$-algèbre $G$ de type fini, et soit
$(t_i)_{1\leq i\leq n}$ un système de générateurs de cette $E$-algèbre~;
montrons que si $(\theta'_\lambda)$, $(\theta''_\lambda)$ sont deux systèmes
inductifs d'homomorphismes $G\to F_\lambda$ tels que
$\varinjlim\theta'_\lambda=\varinjlim\theta''_\lambda$, il existe un
$\mu\geq\lambda$ tel que $\theta'_\mu=\theta''_\mu$. En effet, si
$\varphi_{\mu\lambda}:F_\lambda\to F_\mu$ et
$\varphi_\lambda:F_\lambda\to F=\varinjlim F_\lambda$ sont les
homomorphismes canoniques du système inductif $(F_\lambda)$, par hypothèse,
pour chaque indice $i$, il existe un indice $\lambda_i$ tel que
$\varphi_{\lambda_i}(\theta'_{\lambda_i}(t_i))=
\varphi_{\lambda_i}(\theta''_{\lambda_i}(t_i))$, et l'on peut supposer tous
les $\lambda_i$ égaux à un même $\lambda$~; il en résulte de la même manière
l'existence d'un $\mu\geq\lambda$ tel que
$\varphi_{\mu\lambda}(\theta'_\lambda(t_i))=
\varphi_{\mu\lambda}(\theta''_\lambda(t_i))$ pour $1\leq i\leq n$,
c'est-à-dire $\theta'_\mu(t_i)=\theta''_\mu(t_i)$ pour $1\leq i\leq n$,
d'où $\theta'_\mu=\theta''_\mu$.

Supposons en second lieu $G$ de présentation finie, de sorte que l'on a
\[
G=E[T_1,\ldots,T_n]/\mathfrak J,
\]
où $\mathfrak J$ est un idéal de type fini, $t_i$ étant la classe de $T_i$
mod. $\mathfrak J$. Soit $(P_j)_{1\leq j\leq m}$ un système de générateurs de
$\mathfrak J$. Supposons donné un homomorphisme de $E$-algèbres
$\theta:G\to F$~; posons $b_i=\theta(t_i)$~; par définition, on a donc
$P_j(b_1,b_2,\ldots,b_n)=\theta(P_j(t_1,\ldots,t_n))=0$ pour
$1\leq j\leq m$. Or, il existe un $\lambda$ et des éléments
$x_1,\ldots,x_n$ dans $F_\lambda$ tels que $b_i=\varphi_\lambda(x_i)$ pour
$1\leq i\leq n$~; on a donc
$\varphi_\lambda(P_j(x_1,\ldots,x_n))=P_j(b_1,\ldots,b_n)=0$, et par suite il
existe $\mu\geq\lambda$ tel que
$\varphi_{\mu\lambda}(P_j(x_1,\ldots,x_n))=
P_j(\varphi_{\mu\lambda}(x_1),\ldots,\varphi_{\mu\lambda}(x_n))=0$ pour
$1\leq j\leq m$~; on en conclut qu'il existe un homomorphisme de
$E$-algèbres $\theta_\mu:G\to F_\mu$ tel que
$\theta_\mu(t_i)=\varphi_{\mu\lambda}(x_i)$
\oldpage[IV-3]{30}
pour $1\leq i\leq n$~; on en déduit pour tout $\nu\geq\mu$ un
homomorphisme de $E$-algèbres
$\theta_\nu=\varphi_{\nu\mu}\circ\theta_\mu$ de $G$ dans $F_\nu$, et il est
clair que $\theta$ est la limite inductive de ce système inductif
d'homomorphismes, ce qui achève de démontrer le lemme.

Passons maintenant au cas où $X_\alpha$ est quasi-compact et $Y_\alpha$
localement de type fini sur $S_\alpha$. Posons
$Z_\alpha=X_\alpha\times_{S_\alpha}Y_\alpha$ et introduisons le système
projectif correspondant des
$Z_\lambda=Z_\alpha\times_{S_\alpha}S_\lambda=
X_\lambda\times_{S_\lambda}Y_\lambda$ et sa limite
$Z=Z_\alpha\times_{S_\alpha}S=X\times_S Y$~; les bijections canoniques
(I, 3.3.14) donnent des diagrammes commutatifs
\[
\begin{tikzcd}[column sep=large,row sep=large]
\operatorname{Hom}_{S_\lambda}(X_\lambda,Y_\lambda)
  \ar[r,"e_{\mu\lambda}"] \ar[d,"\wr"] &
\operatorname{Hom}_{S_\mu}(X_\mu,Y_\mu)
  \ar[r,"e_\mu"] \ar[d,"\wr"] &
\operatorname{Hom}_S(X,Y) \ar[d,"\wr"] \\
\operatorname{Hom}_{X_\lambda}(X_\lambda,Z_\lambda)
  \ar[r,"e_{\mu\lambda}"'] &
\operatorname{Hom}_{X_\mu}(X_\mu,Z_\mu)
  \ar[r,"e_\mu"'] &
\operatorname{Hom}_X(X,Z)
\end{tikzcd}
\]
et par suite on est ramené à prouver que (8.8.1.1) est injective dans le cas
particulier où $S_\alpha=X_\alpha$ (compte tenu de (I, 3.4)). De plus, comme
$X_\alpha$ est quasi-compact, donc réunion finie d'ouverts affines, on peut
supposer $X_\alpha$ affine ($L$ étant filtrant). Supposons donc donnés deux
$X_\alpha$-morphismes $f'_\alpha:X_\alpha\to Y_\alpha$,
$f''_\alpha:X_\alpha\to Y_\alpha$ tels que $f'$ et $f''$ soient des
$X$-morphismes égaux de $X$ dans $Y$~; il s'agit de prouver que
$f'_\mu=f''_\mu$ pour $\mu\geq\alpha$ assez grand. Comme $X_\alpha$ est
quasi-compact, il en est de même de
$f'_\alpha(X_\alpha)\cup f''_\alpha(X_\alpha)$, et comme $Y_\alpha$ est
localement de type fini sur $X_\alpha$,
$f'_\alpha(X_\alpha)\cup f''_\alpha(X_\alpha)$ est contenu dans une réunion
finie d'ouverts affines $V_{i\alpha}$ de $Y_\alpha$, de type fini sur
$X_\alpha$. Posons
$U'_{i\alpha}=f_\alpha^{'-1}(V_{i\alpha})$,
$U''_{i\alpha}=f_\alpha^{''-1}(V_{i\alpha})$,
$U_{i\alpha}=U'_{i\alpha}\cap U''_{i\alpha}$,
$U_\alpha=\bigcup_iU_{i\alpha}$. L'hypothèse $f'=f''$ entraîne
$v_\alpha^{-1}(U'_{i\alpha})=v_\alpha^{-1}(U''_{i\alpha})$, ces deux ensembles
étant respectivement égaux à $f^{'-1}(w_\alpha^{-1}(V_{i\alpha}))$ et à
$f^{''-1}(w_\alpha^{-1}(V_{i\alpha}))$. Comme les $V_{i\alpha}$ recouvrent
$Y_\alpha$, on a $v_\alpha^{-1}(U_\alpha)=f^{'-1}(Y)=X$. Comme $X_\alpha$
est quasi-compact et que toute partie ouverte de $X_\lambda$ est
ind-constructible, il résulte de (8.3.4) qu'il y a un indice
$\lambda\geq\alpha$ tel que les
$U_{i\mu}=v_{\alpha\mu}^{-1}(U_{i\alpha})$ forment un \emph{recouvrement} de
$X_\mu$. Remplaçant $\alpha$ par $\mu$, on peut donc supposer que les
$U_{i\alpha}$ recouvrent $X_\alpha$~; cela entraîne que pour tout
$x\in X_\alpha$, il y a un voisinage ouvert affine $W(x)$ contenu dans un des
$U_{i\alpha}$, autrement dit tel que les restrictions de $f'_\alpha$ et
$f''_\alpha$ à $W(x)$ appliquent $W(x)$ dans un \emph{même} ouvert affine
$V_{i\alpha}$. Comme $X_\alpha$ est quasi-compact, il est recouvert par un
nombre fini d'ouverts affines $W(x_k)$~; en vertu du lemme (8.8.2.1) et du fait
que $L$ est filtrant, il existe un $\lambda\geq\alpha$ tel que les
restrictions de $f'_\lambda$ et $f''_\lambda$ à chacun des ouverts
$v_{\alpha\lambda}^{-1}(W(x_k))$ soient égales, d'où
$f'_\lambda=f''_\lambda$.

Supposons maintenant $X_\alpha$ quasi-compact et quasi-séparé et $Y_\alpha$
localement de présentation finie sur $S_\alpha$, et prouvons que (8.8.1.1) est
surjective. Supposons donc donné un $X$-morphisme $f:X\to Y$. Comme $X$ est
quasi-compact, il en est de même
\oldpage[IV-3]{31}
de $f(X)$, et par suite il y a un ouvert quasi-compact $Y'$ dans $Y$ qui
contient $f(X)$~; il existe par suite un $\lambda\geq\alpha$ et un ouvert
quasi-compact $Y'_\lambda$ dans $Y_\lambda$ tels que
$Y'=w_\lambda^{-1}(Y'_\lambda)$ (8.2.11). Remplaçant au besoin $\alpha$ par
$\lambda$ et $Y_\lambda$ par $Y'_\lambda$, on peut donc se borner au cas où
$Y_\alpha$ est quasi-compact, donc aussi $Y$ et les $Y_\lambda$. Comme $Y$ est
quasi-compact, il est réunion finie d'ouverts affines $V_i$, et par suite $X$
est réunion des ouverts $f^{-1}(V_i)$. Comme tout point de $X$ a un voisinage
ouvert quasi-compact contenu dans un des $f^{-1}(V_i)$ et que $X$ est
quasi-compact, on peut, en tenant compte de (8.2.11) et remplaçant au besoin
$\alpha$ par un indice $\lambda\geq\alpha$, supposer que $Y$ est réunion finie
d'ouverts $V_i=w_\alpha^{-1}(V_{i\alpha})$, où les $V_{i\alpha}$ sont des
ouverts affines de $Y_\alpha$~; par suite $X$ est réunion des ouverts
$f^{-1}(V_i)$. Comme tout point de $X$ a un voisinage ouvert quasi-compact
contenu dans l'un des $f^{-1}(V_i)$ et que $X$ est quasi-compact, on peut
recouvrir $X$ par un nombre fini de tels voisinages, et (en répétant au besoin
certains des $V_i$) supposer que l'on a un recouvrement $(U_i)$ de $X$ par des
ouverts quasi-compacts ayant même ensemble d'indices que $(V_i)$ et tel que
$f(U_i)\subset V_i$ pour tout $i$. En outre, à l'aide de (8.2.11) (et en
remplaçant au besoin $\alpha$ par un indice $\lambda\geq\alpha$), on peut
supposer que l'on a $U_i=v_\alpha^{-1}(U_{i\alpha})$ où les $U_{i\alpha}$ sont
des ouverts quasi-compacts dans $X_\alpha$~; de plus, utilisant (8.3.4) comme
ci-dessus, on peut supposer que $(U_{i\alpha})$ est un \emph{recouvrement} de
$X_\alpha$.

Cela étant, il suffira de montrer qu'il existe un $\lambda\geq\alpha$ et pour
chaque $i$ un morphisme $f_{i\lambda}:U_{i\lambda}\to V_{i\lambda}$ (avec
$U_{i\lambda}=v_{\alpha\lambda}^{-1}(U_{i\alpha})$ et
$V_{i\lambda}=w_{\alpha\lambda}^{-1}(V_{i\alpha})$) tels que le morphisme
correspondant $f_i=e_\lambda(f_{i\lambda}):U_i\to V_i$ soit égal à la
\emph{restriction} de $f$ à $U_i$. En effet, s'il en est ainsi, comme
$X_\lambda$ est quasi-séparé (I, 1.2.3), les
$U_{i\lambda}\cap U_{j\lambda}$ sont \emph{quasi-compacts} et le résultat
d'unicité démontré plus haut (qui s'applique puisque $Y_\lambda$ est
localement de type fini sur $S_\lambda$) prouve qu'il existe
$\mu\geq\lambda$ tel que $f_{i\mu}$ et $f_{j\mu}$ coïncident dans
$U_{i\mu}\cap U_{j\mu}$ pour tout couple $(i,j)$, et par suite définissent un
$S_\mu$-morphisme $f_\mu:X_\mu\to Y_\mu$ tel que $e_\mu(f_\mu)=f$.

On est ainsi ramené au cas où $Y_\alpha$ est \emph{affine}, et comme en outre
on peut supposer que les $V_{i\alpha}$ ont des images dans $S_\alpha$
contenues dans des ouverts affines, on peut aussi supposer que $S_\alpha$ est
affine, soient donc $S_\alpha=\operatorname{Spec}(A_\alpha)$,
$Y_\alpha=\operatorname{Spec}(C_\alpha)$, $C_\alpha$ étant une
$A_\alpha$-algèbre de présentation finie, $S=\operatorname{Spec}(A)$,
$Y=\operatorname{Spec}(C)$, avec $A=\varinjlim A_\lambda$,
$C=\varinjlim(C_\alpha\otimes_{A_\alpha}A_\lambda)=
C_\alpha\otimes_{A_\alpha}A$. On a alors
\[
\operatorname{Hom}_S(X,Y)=
\operatorname{Hom}_{A\text{-alg.}}(C,\Gamma(X,\mathcal O_X))=
\operatorname{Hom}_{A_\alpha\text{-alg.}}(C_\alpha,\Gamma(X,\mathcal O_X))
\]
(I, 2.2.4) et de même
\[
\operatorname{Hom}_{S_\lambda}(X_\lambda,Y_\lambda)=
\operatorname{Hom}_{A_\lambda\text{-alg.}}
(C_\alpha\otimes_{A_\alpha}A_\lambda,\Gamma(X_\lambda,\mathcal O_{X_\lambda}))=
\operatorname{Hom}_{A_\alpha\text{-alg.}}
(C_\alpha,\Gamma(X_\lambda,\mathcal O_{X_\lambda})).
\]
Mais comme $X_\alpha$ est quasi-compact et quasi-séparé, on sait (8.5.4) que
l'on a
$\varinjlim\Gamma(X_\lambda,\mathcal O_{X_\lambda})=
\Gamma(X,\mathcal O_X)$~; puisque $C_\alpha$ est une $A_\alpha$-algèbre de
présentation finie, le fait que (8.8.1.1) est bijectif résulte alors de
(8.8.2.3).

Avant de passer à la démonstration de (ii), notons les corollaires suivants de
(i)~:
\end{proof}

\begin{corollary}[8.8.2.4]
\label{IV.8.8.2.4-fr}
Supposons $S_0$ quasi-compact, $X_\alpha$ de présentation finie sur $S_\alpha$,
$Y_\alpha$ de type fini sur $S_\alpha$ et quasi-séparé sur $S_\alpha$ (ce qui
sera par exemple le cas si $Y_\alpha$ est également
\oldpage[IV-3]{32}
de présentation finie sur $S_\alpha$). Soit
$f_\alpha:X_\alpha\to Y_\alpha$ un $S_\alpha$-morphisme. Pour que
$f:X\to Y$ soit un isomorphisme, il faut et il suffit qu'il existe
$\lambda\geq\alpha$ tel que $f_\lambda:X_\lambda\to Y_\lambda$ soit un
isomorphisme.
\end{corollary}

\begin{proof}
La condition est évidemment suffisante. Pour prouver qu'elle est nécessaire,
notons d'abord que la question étant locale sur $S_0$ (puisque $L$ est
filtrant), on peut toujours supposer $S_0$ affine, donc quasi-séparé. Il y a
par hypothèse un $S$-morphisme $g:Y\to X$ tel que $g\circ f=1_X$ et
$f\circ g=1_Y$. Comme $X_\alpha$ est de présentation finie sur $S_\alpha$,
et $Y_\alpha$ quasi-compact et quasi-séparé (puisque $S_\alpha$ est
quasi-compact et quasi-séparé), il existe un $\mu\geq\alpha$ et un
$S_\mu$-morphisme $g_\mu:Y_\mu\to X_\mu$ tels que
$g=g_\mu\times1_S$ en vertu de (8.8.2, (i)). D'autre part, il résulte aussi de
(8.8.2, (i)) et des relations $g\circ f=1_X$ et $f\circ g=1_Y$ qu'il existe
$\nu\geq\mu$ tel que l'on ait $g_\nu\circ f_\nu=1_{X_\nu}$ et
$f_\nu\circ g_\nu=1_{Y_\nu}$, puisque $X_\alpha$ et $Y_\alpha$ sont de type
fini sur $S_\alpha$ et quasi-compacts. Cela signifie que
$f_\nu:X_\nu\to Y_\nu$ est un isomorphisme, d'où le corollaire.
\end{proof}

\begin{corollary}[8.8.2.5]
\label{IV.8.8.2.5-fr}
Supposons $S_0$ quasi-compact et quasi-séparé, $X_\alpha$ et $Y_\alpha$ de
présentation finie sur $S_\alpha$. Pour que $X$ et $Y$ soient
$S$-isomorphes, il faut et il suffit qu'il existe $\lambda\geq\alpha$ tels que
$X_\lambda$ et $Y_\lambda$ soient $S_\lambda$-isomorphes. De plus, pour tout
$S$-isomorphisme $f:X\xrightarrow{\sim}Y$, il existe $\mu\geq\lambda$ et un
$S_\mu$-isomorphisme $f_\mu:X_\mu\to Y_\mu$ tels que
$f=f_\mu\times1_S$.
\end{corollary}

\begin{proof}
La condition est évidemment suffisante~; inversement, s'il existe un
$S$-isomorphisme $f:X\to Y$, il résulte de (8.8.2, (i)) que $f$ est de la
forme $f_\mu\times1_S$ pour un $\mu\geq\alpha$ et un homomorphisme
$f_\mu:X_\mu\to Y_\mu$~; mais comme $f$ est un isomorphisme, il résulte de
(8.8.2.4) qu'il existe $\nu\geq\mu$ tel que
$f_\nu:X_\nu\to Y_\nu$ soit un isomorphisme.
\end{proof}

\begin{proof}[Démonstration de (8.8.2), (ii)]
Considérons encore en premier lieu le cas où $S_0=\operatorname{Spec}(A_0)$ et
$X=\operatorname{Spec}(B)$ sont affines~; alors l'assertion de (ii) équivaut
au lemme (I, 8.4.2).

Pour démontrer (ii) dans le cas général, notons que $S$ est quasi-compact et
quasi-séparé~; comme $X$ est de présentation finie sur $S$ et $S$ affine sur
$S_0$, il existe donc un recouvrement fini $(U_i)$ de $S_0$ et, si $(W_i)$ est
le recouvrement ouvert affine de $S$ formé des images réciproques des $U_i$
par le morphisme $S\to S_0$, un recouvrement fini $(X_r)$ de $X$ par des
ouverts affines tel que l'image par $X\to S$ de chaque $X_r$ soit contenue
dans un $W_{i(r)}$~; l'anneau $A(X_r)$ est alors une algèbre de présentation
finie sur l'anneau $A(W_{i(r)})$ (I, 1.4.6). En vertu du lemme (I, 8.4.2) et du
fait que $L$ est filtrant, il existe un indice $\lambda\in L$ et, pour chaque
indice $r$, un schéma affine $Z_{r\lambda}$, de présentation finie sur l'image
réciproque $W_{i(r),\lambda}$ de $U_{i(r)}$ dans $S_\lambda$, et un
$S$-isomorphisme
$g_r:Z_{r\lambda}\times_{S_\lambda}S\xrightarrow{\sim}X_r$. Soit $Z_{rs}$
l'image réciproque par $g_r$ du sous-préschéma induit sur l'ouvert
$X_r\cap X_s$ de $X_r$, qui est quasi-compact puisque $X$ est quasi-séparé,
et désignons par $g'_{rs}$ la restriction
$Z_{rs}\xrightarrow{\sim}X_r\cap X_s$ de l'isomorphisme $g_r$. En vertu de
(8.8.2.5), il existe un $\mu\geq\lambda$ et, pour tout couple $(r,s)$, un
ouvert quasi-compact $Z_{rs\mu}$ de
$Z_{r\mu}=v_{\lambda\mu}^{-1}(Z_{r\lambda})$ tels que $Z_{rs}$ soit l'image
réciproque de $Z_{rs\mu}$~; d'ailleurs, comme $S_\mu$ est quasi-séparé, et
$W_{i(r),\mu}$ ouvert quasi-compact dans $S_\mu$, chacun des $Z_{rs\mu}$ est
de présentation finie sur $S_\mu$ (I, 1.6.2). Considérons alors, pour tout
couple $(r,s)$, l'isomorphisme
$h_{sr}=g_{sr}^{'-1}\circ g'_{rs}$ de $Z_{rs}$ sur $Z_{sr}$~; il résulte de
(8.8.2.4) qu'il existe un $\nu\geq\mu$ et, pour tout couple $(r,s)$, un
isomorphisme $h_{sr\nu}:Z_{rs\nu}\to Z_{sr\nu}$ tels que
$h_{sr}=h_{sr\nu}\times1_S$. Enfin, pour tout triplet $(r,s,t)$ désignons par
$h'_{sr}$ la restriction de $h_{sr}$ à
\oldpage[IV-3]{33}
$Z_{rs}\cap Z_{rt}$~; il résulte aussitôt des définitions que $h'_{sr}$ est un
isomorphisme de $Z_{rs}\cap Z_{rt}$ sur $Z_{sr}\cap Z_{st}$, et que l'on a la
relation $h'_{ts}\circ h'_{sr}=h'_{tr}$. En vertu de (8.8.2, (i)), il existe
donc $\rho\geq\nu$ tel que, pour tout triplet $(r,s,t)$, on ait
$h'_{ts\rho}\circ h'_{sr\rho}=h'_{tr\rho}$. On en conclut que les
isomorphismes $h_{sr\rho}$ vérifient la condition de recollement
$(0_{\mathrm I}, 4.1.7)$ et définissent donc un préschéma $X_\rho$ tel que $X$
soit isomorphe à $X_\rho\times_{S_\rho}S$. En outre, les $Z_{r\rho}$ sont de
présentation finie sur $S_\rho$ et, si on les identifie à des ouverts de
$X_\rho$, les intersections $Z_{r\rho}\cap Z_{s\rho}$, isomorphes aux
$Z_{rs\rho}$, sont quasi-compactes, donc (I, 1.6.3) $X_\rho$ est de
présentation finie sur $S_\rho$.
\end{proof}

\begin{env}[8.8.3]
\label{IV.8.8.3-fr}
Le contenu essentiel de (8.8.2) s'exprime encore en disant que si $S_0$ est
quasi-compact et quasi-séparé, la catégorie des $S$-préschémas de présentation
finie est déterminée à équivalence près par la donnée des catégories des
$S_\lambda$-préschémas de présentation finie, des foncteurs
$\rho_{\mu\lambda}:X_\lambda\rightsquigarrow
X_\lambda\times_{S_\lambda}S_\mu$ (pour $\lambda\leq\mu$) entre ces catégories
et des isomorphismes de transitivité
$\rho_{\nu\mu}\circ\rho_{\mu\lambda}\xrightarrow{\sim}\rho_{\nu\lambda}$. De
façon imagée, on peut dire que la donnée d'un $S$-préschéma de présentation
finie $X$ équivaut «~fonctoriellement~» à la donnée d'un
$S_\lambda$-préschéma de présentation finie $X'_\lambda$~; si un
$S_\mu$-préschéma de présentation finie $X''_\mu$ est tel que $X$ soit
$S$-isomorphe à $X''_\mu\times_{S_\mu}S$, il existe un $\nu$ tel que
$\lambda\leq\nu$, $\mu\leq\nu$ et tel que
$X'_\lambda\times_{S_\lambda}S_\nu$ et
$X''_\mu\times_{S_\mu}S_\nu$ soient $S_\nu$-isomorphes. Le fait que $L$ est
filtrant entraîne en outre que si on se donne une famille \emph{finie}
$(X^{(i)})_{i\in I}$ de $S$-préschémas de présentation finie, et une famille
\emph{finie} $(f^{(j)})_{j\in J}$ de $S$-morphismes entre ces préschémas
($f^{(j)}$ étant donc de la forme
$X^{(\sigma(j))}\to X^{(\tau(j))}$ où $\sigma$ et $\tau$ sont deux
applications de $J$ dans $I$), alors il y a un indice $\lambda\in L$, une
famille $(X_\lambda^{(i)})_{i\in I}$ de $S_\lambda$-préschémas de présentation
finie et une famille $(f_\lambda^{(j)})_{j\in J}$ de $S_\lambda$-morphismes
tels que $X^{(i)}$ s'identifie à
$X_\lambda^{(i)}\times_{S_\lambda}S$ et $f^{(j)}$ à
$f_\lambda^{(j)}\times1_S$ quels que soient $i$ et $j$~; en outre, si on a
une relation $f^{(j)}=f^{(k)}$, on peut supposer $\lambda$ choisi de telle
sorte que $f_\lambda^{(j)}=f_\lambda^{(k)}$.

Considérons en particulier trois $S$-préschémas de présentation finie $X$,
$Y$, $Z$ et deux $S$-morphismes $f:X\to Z$, $g:Y\to Z$, de sorte que le
produit $X\times_ZY$ relatif à ces morphismes est encore un $S$-préschéma de
présentation finie (I, 1.6.2). Alors il résulte de ce qui précède et de
(I, 3.3.11) que l'on peut écrire
$X\times_ZY=(X_\lambda\times_{Z_\lambda}Y_\lambda)\times_{S_\lambda}S$ pour
un $\lambda$ convenable, $X_\lambda$, $Y_\lambda$, $Z_\lambda$ étant des
$S_\lambda$-préschémas de présentation finie~; on peut donc dire que la
détermination des $S$-préschémas de présentation finie par la donnée des
$S_\lambda$-préschémas de présentation finie est «~\emph{compatible avec les
produits fibrés}~». On a vu d'autre part (8.6.3) que si $g$ est une immersion,
on peut supposer qu'il en est de même de $g_\lambda:Y_\lambda\to Z_\lambda$~;
identifiant $Y$ (resp. $Y_\lambda$) avec un sous-préschéma de $Z$ (resp.
$Z_\lambda$), on voit donc que l'on peut écrire, pour un $\lambda$ convenable,
$f^{-1}(Y)=f_\lambda^{-1}(Y_\lambda)\times_{S_\lambda}S$ (I, 4.4.1)~; il y a
donc aussi «~\emph{compatibilité avec la formation des images réciproques de
sous-préschémas}~». Plus particulièrement, si $f_1$, $f_2$ sont deux
$S$-morphismes de $X$ dans $Y$, on appelle \emph{noyau} de ce couple de
morphismes l'image réciproque $N$ du sous-préschéma diagonal de $Y\times_SY$
par le $S$-morphisme $(f_1,f_2)_S$~; on déduit de ce qui précède que l'on aura
pour $\lambda$ convenable $N=N_\lambda\times_{S_\lambda}S$, où $N_\lambda$
est le noyau du couple de $S_\lambda$-morphismes
$(f_{1\lambda},f_{2\lambda})$ de $X_\lambda$ dans $Y_\lambda$. Ces remarques
s'étendent à des produits finis quelconques et aux «~noyaux~» de systèmes
finis quelconques de $S$-morphismes d'un $S$-préschéma de présentation finie
\oldpage[IV-3]{34}
dans un autre~; on en conclut que d'une façon générale la formation des
$S$-préschémas de présentation finie par la donnée des
$S_\lambda$-préschémas de présentation finie est «~\emph{compatible avec les
limites projectives finies}~», une telle limite étant par définition le noyau
d'un système fini de morphismes d'un même $S$-préschéma dans un produit de
$S$-préschémas (T, 1.8).

Nous nous permettrons couramment par la suite de faire les traductions
impliquées par les propriétés précédentes (ainsi que par (8.3.11), (8.5.2) et
(8.6.3)) sans toujours nous astreindre à les justifier pas à pas comme
ci-dessus. Par exemple, la donnée d'un «~\emph{préschéma en groupes}~» $G$ de
présentation finie sur $S$ équivaut à la donnée d'un préschéma en groupes
$G_\lambda$ de présentation finie sur un $S_\lambda$ pour un $\lambda$ assez
grand~; écrire en effet la condition d'associativité pour le $S$-morphisme
«~loi de composition~» $G\times_SG\to G$ du préschéma en groupes revient à
écrire que le noyau de deux $S$-morphismes de la forme
$G\times_SG\times_SG\to G$ est égal à $G\times_SG\times_SG$ (II, 8.3.9), et
les autres conditions intervenant dans la définition d'un préschéma en
groupes s'interprètent de même.
\end{env}

\subsection{Premières applications à l'élimination des hypothèses noethériennes}
\label{subsection:IV.8.9-fr}

\begin{proposition}[8.9.1]
\label{IV.8.9.1-fr}
Soient $A$ un anneau, $X$ un $A$-préschéma.
\begin{enumerate}
\item[(i)] Les conditions suivantes sont équivalentes~:
\begin{enumerate}
\item[\emph{a)}] $X$ est de présentation finie sur $A$.
\item[\emph{b)}] Il existe un anneau noethérien $A_0$, un préschéma $X_0$ de
type fini sur $A_0$, un homomorphisme d'anneaux $A_0\to A$, et un
$A$-isomorphisme $X_0\otimes_{A_0}A\xrightarrow{\sim}X$.
\item[\emph{c)}] Il existe un sous-anneau $A_0$ de $A$, qui est une
$\mathbf Z$-algèbre de type fini, un préschéma $X_0$ de type fini sur $A_0$ et
un $A$-isomorphisme $X_0\otimes_{A_0}A\xrightarrow{\sim}X$.
\end{enumerate}
\item[(ii)] Si $\mathcal F$ est un $\mathcal O_X$-Module quasi-cohérent de
présentation finie, on peut supposer le sous-anneau $A_0$ tel qu'il existe un
$\mathcal O_{X_0}$-Module cohérent $\mathcal F_0$ tel que $\mathcal F$ soit
isomorphe à $\mathcal F_0\otimes_{A_0}A$~;
$\operatorname{Supp}(\mathcal F)$ est constructible et fermé dans $X$, et il
y a un sous-préschéma fermé $Z$ de $X$, ayant
$\operatorname{Supp}(\mathcal F)$ comme espace sous-jacent, tel que
l'immersion canonique $Z\to X$ soit de présentation finie.
\item[(iii)] Si $Y$ est un second $A$-préschéma de présentation finie, et
$f:X\to Y$ un $A$-morphisme, on peut supposer le sous-anneau $A_0$ de $A$ tel
qu'il existe un préschéma $Y_0$ de type fini sur $A_0$, un $A$-isomorphisme
$Y_0\otimes_{A_0}A\xrightarrow{\sim}Y$ et un $A_0$-morphisme
$f_0:X_0\to Y_0$ (nécessairement de type fini) tels que $f$ s'identifie à
$f_0\otimes1_A$.
\end{enumerate}
\end{proposition}

\begin{proof}
(i) Comme $A$ est limite inductive de ses sous-anneaux qui sont de type fini
sur $\mathbf Z$, \emph{a)} implique \emph{c)} en vertu de (8.8.2, (ii))~;
\emph{c)} implique \emph{b)} car une $\mathbf Z$-algèbre de type fini est un
anneau noethérien~; enfin, si $A_0$ est noethérien, un $A_0$-préschéma de type
fini est de présentation finie (I, 1.6.1), donc \emph{b)} implique \emph{a)}
en vertu de (I, 1.6.2, (iii)).

L'assertion (ii) se déduit immédiatement de (8.5.2, (ii)), (8.3.11) et
(I, 1.8.2) et l'assertion (iii) de (8.8.2, (i)).
\end{proof}

\begin{env}[8.9.2]
\label{IV.8.9.2-fr}
La proposition (8.9.1) et les résultats de (8.5.2) et (8.8.2) permettent
d'étendre dans de nombreux cas à des morphismes de présentation finie
$X\to Y$ des résultats prouvés par des techniques noethériennes en supposant
$Y$ \emph{localement noethérien}.
\oldpage[IV-3]{35}
Nous en verrons de nombreux exemples par la suite et nous nous bornerons ici
à donner quelques résultats de ce type.
\end{env}

\begin{proposition}[8.9.3]
\label{IV.8.9.3-fr}
Soient $A$ un anneau, $M$ un $A$-module de présentation finie. Alors tout
$A$-endomorphisme surjectif de $M$ est bijectif.
\end{proposition}

\begin{proof}
En effet, considérons $A$ comme limite inductive de ses sous-$\mathbf Z$-algèbres
de type fini. Il résulte de (8.5.2.6) qu'il existe une de ces sous-algèbres
$A_0$ et un $A_0$-module $M_0$ de présentation finie tels que $M$ soit
$A$-isomorphe à $M_0\otimes_{A_0}A$~; en outre, si $f:M\to M$ est un
$A$-endomorphisme surjectif, on peut supposer (8.5.2, (i)) qu'il existe un
$A_0$-endomorphisme $f_0:M_0\to M_0$ tel que $f=f_0\otimes1_A$~; enfin
(8.5.7) on peut supposer que $f_0$ est surjectif. Mais comme $A_0$ est
noethérien et $M_0$ un $A_0$-module de type fini, $M_0$ est un $A_0$-module
noethérien, donc (Bourbaki, \emph{Alg.}, chap. VIII, \S~2, n\textsuperscript{o}~2,
lemme 3) $f_0$ est bijectif, et il en est par suite de même de $f$.
\end{proof}

\begin{proposition}[8.9.4]
\label{IV.8.9.4-fr}
(«~théorème de platitude générique~»). --- Soient $Y$ un préschéma intègre,
$u:X\to Y$ un morphisme de type fini et localement de présentation finie,
$\mathcal F$ un $\mathcal O_X$-Module quasi-cohérent de présentation finie. Il
existe alors un ouvert non vide $U$ de $Y$ tel que
$\mathcal F|u^{-1}(U)$ soit plat sur $U$.
\end{proposition}

\begin{proof}
Le raisonnement du début de (6.9.1) ramène à prouver le

\begin{lemma}[8.9.4.1]
\label{IV.8.9.4.1-fr}
Soient $A$ un anneau intègre, $B$ une $A$-algèbre de présentation finie,
$M$ un $B$-module de présentation finie. Alors il existe un $f\neq0$ dans
$A$ tel que $M_f$ soit un $A_f$-module libre.
\end{lemma}

En effet, d'après (8.9.1) il y a une sous-$\mathbf Z$-algèbre de type fini
$A_0$ de $A$, une $A_0$-algèbre de type fini $B_0$ et un $B_0$-module de type
fini $M_0$ tels que $B$ soit $A$-isomorphe à $B_0\otimes_{A_0}A$ et que $M$
soit $B$-isomorphe à $M_0\otimes_{A_0}A$~; d'après (6.9.2), il existe
$f_0\neq0$ dans $A_0$ tel que $(M_0)_{f_0}$ soit un
$(A_0)_{f_0}$-module libre. Comme
$M_{f_0}=(M_0)_{f_0}\otimes_{A_0}A$ et
$A_{f_0}=(A_0)_{f_0}\otimes_{A_0}A$, $M_{f_0}$ est un $A_{f_0}$-module libre.
\end{proof}

\begin{corollary}[8.9.5]
\label{IV.8.9.5-fr}
Soient $S$ un préschéma quasi-compact et quasi-séparé, $u:X\to S$ un
morphisme de présentation finie, $\mathcal F$ un $\mathcal O_X$-Module
quasi-cohérent de présentation finie. Il existe alors une partition
$(S_i)_{1\leq i\leq n}$ de $S$ en un nombre fini d'ensembles localement fermés
dans $S$ telle que, pour $1\leq i\leq n$, il existe un sous-préschéma $S'_i$
de $S$, ayant pour espace sous-jacent $S_i$, de présentation finie sur $S$,
tel que si l'on pose $X_i=X\times_SS'_i$, le $\mathcal O_{X_i}$-Module
$\mathcal F_i:=\mathcal F\otimes_{\mathcal O_S}\mathcal O_{S'_i}$ soit plat
sur $S'_i$.
\end{corollary}

\begin{proof}
Considérons un recouvrement fini $(U_j)_{1\leq j\leq m}$ de $S$ par des
ouverts affines et définissons par récurrence $T_1=U_1$,
$T_j=U_j-\bigcup_{h<j}(U_j\cap U_h)$ pour $j\geq2$~; chaque $T_j$ est fermé
dans l'ouvert affine $U_j$, et les $T_j$ sont deux à deux sans point commun~;
en outre les $U_j\cap U_h$ sont quasi-compacts puisque $S$ est quasi-séparé,
et par suite ($S$ étant aussi quasi-compact), sont constructibles dans $S$
(I, 1.8.1), donc il en est de même des $T_j$. Il suffira évidemment de
démontrer le corollaire pour un sous-préschéma convenable $T'_j$ de $S$ ayant
$T_j$ pour espace sous-jacent, de présentation finie sur $S$ et pour le
morphisme et le Module déduits respectivement de $u$ et $\mathcal F$ par le
changement de base $T'_j\to S$. Notons pour cela que si l'on prend sur $U_j$
la structure de préschéma induite par celle de $S$, l'immersion ouverte
$U_j\to S$ est quasi-compacte puisque $S$ est quasi-séparé (I, 1.2.7), donc
elle est de
\oldpage[IV-3]{36}
présentation finie (I, 1.6.2, (i)). Il suffira donc que $T'_j$ soit de
présentation finie sur $U_j$, autrement dit, on peut déjà se borner au cas où
$U_j=S$ et $T_j=T$ est fermé constructible dans $S$.

Soit $S=\operatorname{Spec}(A)$, et considérons $A$ comme limite inductive de
ses sous-$\mathbf Z$-algèbres de type fini, de sorte que
$S=\varprojlim S_\lambda$, où les $S_\lambda$ sont affines et noethériens. En
vertu de (8.3.11), il existe un $\lambda$ et une partie fermée (nécessairement
constructible) $T_\lambda$ de $S_\lambda$ tels que
$T=u_\lambda^{-1}(T_\lambda)$ ($u_\lambda:S\to S_\lambda$ étant le morphisme
canonique). On munira $T_\lambda$ d'une structure de sous-préschéma de
$S_\lambda$ et on prendra $T'=T_\lambda\times_{S_\lambda}S$~; comme
l'immersion canonique $T_\lambda\to S_\lambda$ est de présentation finie
(I, 1.6.1), il en est de même de l'immersion $T'\to S$. Comme $T'$ est affine,
on voit finalement qu'on peut se borner au cas où $T'=S$ est affine.

Alors (8.9.1), avec les mêmes notations, il existe un $\lambda$, un morphisme
de type fini $u_\lambda:X_\lambda\to S_\lambda$ et un
$\mathcal O_{X_\lambda}$-Module cohérent $\mathcal F_\lambda$ tels que $X$
soit isomorphe à $X_\lambda\times_{S_\lambda}S$, $\mathcal F$ à
$\mathcal F_\lambda\otimes_{\mathcal O_{X_\lambda}}\mathcal O_X$. On peut
alors appliquer à $S_\lambda$, $X_\lambda$ et $\mathcal F_\lambda$ le
résultat de (6.9.3) et il y a des sous-préschémas $S'_{\lambda i}$ de
$S_\lambda$, dont les ensembles sous-jacents forment une partition de
$S_\lambda$ et qui sont tels que si l'on pose
$X_{\lambda i}=X_\lambda\times_{S_\lambda}S'_{\lambda i}$,
$\mathcal F_{\lambda i}=\mathcal F_\lambda\otimes_{\mathcal O_{S_\lambda}}
\mathcal O_{S'_{\lambda i}}$ soit plat sur $S'_{\lambda i}$. Les
$S'_i=S'_{\lambda i}\times_{S_\lambda}S$ sont alors des sous-préschémas de
$S$ répondant à la question en vertu de (2.1.4).
\end{proof}

\subsection{Propriétés de permanence des morphismes par passage à la limite projective}
\label{subsection:IV.8.10-fr}

Dans cette section, nous conservons les hypothèses générales et les notations
de (8.8.1).

\begin{proposition}[8.10.1]
\label{IV.8.10.1-fr}
S'il existe $\lambda$ tel que, pour $\mu\geq\lambda$, $f_\mu$ soit un morphisme
ouvert (resp. universellement ouvert), alors $f$ est un morphisme ouvert
(resp. universellement ouvert).
\end{proposition}

\begin{proof}
Comme $X=X_\lambda\times_{Y_\lambda}Y$, l'assertion relative aux morphismes
universellement ouverts est un cas particulier de (2.4.3, (iii)). Supposons
donc $f_\mu$ ouvert pour $\mu\geq\lambda$~; il suffit de voir que pour tout
ouvert quasi-compact $U$ de $X$, $f(U)$ est ouvert dans $Y$. Or, il existe
$\mu\geq\lambda$ tel que $U=v_\mu^{-1}(U_\mu)$, où $U_\mu$ est un ouvert
quasi-compact dans $X_\mu$ (2.3.11)~; on a alors
$f(v_\mu^{-1}(U_\mu))=w_\mu^{-1}(f_\mu(U_\mu))$ (I, 3.4.8), donc $f(U)$ est
ouvert dans $Y$.
\end{proof}

\begin{corollary}[8.10.2]
\label{IV.8.10.2-fr}
Soit $f:X\to Y$ un morphisme. Pour que $f$ soit universellement ouvert, il
suffit que, pour tout entier $n>0$, si l'on pose
$Y_n=Y\otimes_{\mathbf Z}\mathbf Z[T_1,\ldots,T_n](=\mathbf V_Y^n)$, et
$X_n=X\times_YY_n$, la projection canonique
$f_n=f\times1_{Y_n}:X_n\to Y_n$ soit un morphisme ouvert.
\end{corollary}

\begin{proof}
Pour démontrer que $f$ est universellement ouvert, il suffit de prouver qu'il
en est ainsi de la restriction $f^{-1}(U)\to U$ de $f$ pour tout ouvert affine
$U$ de $Y$~; comme, par hypothèse, si
$U_n=U\otimes_{\mathbf Z}\mathbf Z[T_1,\ldots,T_n]$ est l'image réciproque de
$U$ dans $Y_n$, le morphisme $f_n^{-1}(U_n)\to U_n$, restriction de $f_n$, est
ouvert, on voit qu'on peut se borner au cas où $Y=\operatorname{Spec}(A)$ est
affine. En outre, il suffit aussi évidemment de montrer que pour tout
morphisme $Y'\to Y$, où $Y'=\operatorname{Spec}(A')$ est lui aussi affine,
$f'=f_{(Y')}$ est ouvert. Supposons d'abord que $A'$ soit une $A$-algèbre de
type fini, donc quotient d'une algèbre de polynômes $A[T_1,\ldots,T_n]$~;
alors $Y'$ est un sous-préschéma fermé de $Y_n$ et $f'$ la
\oldpage[IV-3]{37}
restriction de $f_n$ à $f_n^{-1}(Y')$~; mais pour tout ouvert $V$ de $X_n$,
on a $f_n(V\cap f_n^{-1}(Y'))=f_n(V)\cap Y'$, et comme par hypothèse $f_n(V)$
est ouvert dans $Y_n$, cela montre que $f'$ est aussi un morphisme ouvert.
Lorsque $A'$ est quelconque, il peut être considéré comme limite inductive de
ses sous-$A$-algèbres de type fini $A'_\lambda$, et le fait que $f'$ soit un
morphisme ouvert résulte de ce qui précède et de (8.10.1).
\end{proof}

\begin{proposition}[8.10.3]
\label{IV.8.10.3-fr}
Supposons qu'il existe $\alpha$ tel que~: 1\textsuperscript{o} $S_\alpha$ soit
quasi-compact~; 2\textsuperscript{o} les morphismes $X_\alpha\to S_\alpha$,
$Y_\alpha\to S_\alpha$ soient quasi-compacts et le morphisme
$Y_\alpha\to S_\alpha$ quasi-séparé~; 3\textsuperscript{o} pour
$\alpha\leq\lambda\leq\mu$, les morphismes
$u_{\lambda\mu}:S_\mu\to S_\lambda$ soient plats~; 4\textsuperscript{o}
$\overline{f_\alpha(X_\alpha)}$ soit constructible dans $Y_\alpha$. Alors,
pour que $f$ soit dominant, il faut et il suffit qu'il existe
$\lambda\geq\alpha$ tel que $f_\lambda$ soit dominant.
\end{proposition}

\begin{proof}
Les hypothèses entraînent que $Y_\alpha$ est quasi-compact et que le morphisme
$f_\alpha$ est quasi-compact (I, 1.2.4)~; par suite
$\overline{f_\alpha(X_\alpha)}=Z_\alpha$ est pro-constructible (I, 1.9.5,
(vii)) dans $Y_\alpha$. Si l'on pose
$Z_\lambda=w_{\alpha\lambda}^{-1}(Z_\alpha)$ pour $\lambda\geq\alpha$ et
$Z=w_\alpha^{-1}(Z_\alpha)$, on a donc
$Z_\lambda=\overline{f_\lambda(X_\lambda)}$ et
$Z=\overline{f(X)}$ (I, 3.4.8), et $Z_\lambda$ est pro-constructible dans
$Y_\lambda$ (I, 1.9.5, (vii)). Il suffit alors d'appliquer (8.3.13) en
remplaçant $S_\lambda$, $Z'_\lambda$ et $Z''_\lambda$ par $Y_\lambda$,
$Y_\lambda$ et $Z_\lambda$ respectivement.
\end{proof}

\begin{proposition}[8.10.4]
\label{IV.8.10.4-fr}
Supposons qu'il existe $\alpha$ tel que $Y_\alpha$ soit quasi-compact et
$f_\alpha$ de type fini et quasi-séparé. Pour que le morphisme $f$ soit
séparé, il faut et il suffit qu'il existe $\lambda\geq\alpha$ tel que
$f_\lambda$ soit séparé.
\end{proposition}

\begin{proof}
La question étant locale sur $Y_\alpha$ (puisque $Y_\alpha$ est quasi-compact
et $L$ filtrant), on peut se borner au cas où $Y_\alpha$ est affine, donc
quasi-séparé, et l'hypothèse entraîne que $X_\alpha$ (donc les $X_\lambda$ et
$X$) sont quasi-compacts et quasi-séparés. Posons
$X'_\lambda=X_\lambda\times_{Y_\lambda}X_\lambda$ pour
$\lambda\geq\alpha$ et $X'=X\times_YX$~; on a
$X'_\lambda=X'_\alpha\times_{Y_\alpha}Y_\lambda$ et
$X'=X'_\alpha\times_{Y_\alpha}Y$~; le morphisme première projection
$X'_\alpha\to X_\alpha$ est quasi-compact et quasi-séparé par hypothèse
(I, 1.2.3, (iii)), donc $X'_\lambda$ est quasi-compact et quasi-séparé.
Notons maintenant que si l'on désigne par $\Delta_\lambda$ (resp. $\Delta$)
la diagonale de $X_\lambda\times_{Y_\lambda}X_\lambda$ (resp. de
$X\times_YX$), il résulte de (I, 5.3.4 et 3.4.8) que $\Delta_\mu$ (resp.
$\Delta$) est l'image réciproque de $\Delta_\lambda$ par l'application
$v'_{\lambda\mu}:X'_\mu\to X'_\lambda$ (resp.
$v'_\lambda:X'\to X'_\lambda$). D'autre part $\Delta_\lambda$ est
constructible dans $X'_\lambda$~: en effet, puisque $f_\lambda$ est
quasi-séparé, l'immersion diagonale $X_\lambda\to X'_\lambda$ est
quasi-compacte, et localement de présentation finie puisque $f_\lambda$ est
de type fini (I, 4.3 et I, 5.4, (iii))~; il résulte alors de (1.8.4.1) que
$\Delta_\lambda$ est localement constructible, donc constructible puisque
$X'_\lambda$ est quasi-compact. On peut maintenant appliquer (8.3.12) en
remplaçant $S_\lambda$ et $Z_\lambda$ par $X'_\lambda$ et $\Delta_\lambda$
respectivement.
\end{proof}

\begin{theorem}[8.10.5]
\label{IV.8.10.5-fr}
Supposons $S_0$ quasi-compact, $X_\alpha$ et $Y_\alpha$ de présentation finie
sur $S_\alpha$, et soit $f_\alpha:X_\alpha\to Y_\alpha$ un
$S_\alpha$-morphisme. Considérons, pour un morphisme, la propriété d'être~:
\begin{enumerate}
\item[(i)] un isomorphisme~;
\item[(i \emph{bis})] un monomorphisme~;
\item[(ii)] une immersion~;
\item[(iii)] une immersion ouverte~;
\item[(iv)] une immersion fermée~;
\item[(v)] séparé~;
\item[(vi)] surjectif~;
\oldpage[IV-3]{38}
\item[(vii)] radiciel~;
\item[(viii)] affine~;
\item[(ix)] quasi-affine~;
\item[(x)] fini~;
\item[(xi)] quasi-fini~;
\item[(xii)] propre.
\end{enumerate}
Alors, si $\mathbf P$ désigne une des propriétés précédentes, pour que $f$
possède la propriété $\mathbf P$, il faut et il suffit qu'il existe
$\lambda\geq\alpha$ tel que $f_\lambda$ possède la propriété $\mathbf P$
(auquel cas $f_\mu$ la possède aussi pour $\mu\geq\lambda$).

Si $S_0$ est en outre supposé quasi-séparé, la même conclusion est valable
lorsque $\mathbf P$ est la propriété d'être~:
\begin{enumerate}
\item[(xiii)] projectif~;
\item[(xiv)] quasi-projectif.
\end{enumerate}
\end{theorem}

\begin{proof}
Le cas où $\mathbf P$ est l'une des propriétés (i) ou (v) n'a été inséré dans
l'énoncé que pour mémoire~; dans ces cas, le théorème découle de ce qui a été
démontré respectivement dans (8.8.2.4) et (8.10.4). D'ailleurs, compte tenu de
(I, 5.4.1 et 5.3.4), (v) résulte aussi de (iv). Le cas (i \emph{bis}) se
déduit aussitôt de (i), en utilisant (I, 5.3.8) et en notant (comme on l'a
déjà utilisé dans (8.10.4)) que la diagonale $\Delta_f$ se déduit de
$\Delta_{f_\lambda}$ par le changement de base $S\to S_\lambda$.

On notera d'autre part que (vi), (vii) et (xi) sont en fait des conditions
sur les fibres $f^{-1}(y)$ des morphismes envisagés, compte tenu de la
transitivité des fibres par changement de base (I, 3.6.4)~: la condition (vi)
signifie en effet que toutes les fibres doivent être non vides, la condition
(vii) qu'elles doivent être radicielles (I, 3.5.8) et la condition (xi)
qu'elles doivent être finies (III, 6.2.2 et 6.2.3 et I, 6.4.4, compte tenu de
ce que $f$ et les $f_\lambda$ sont des morphismes de type fini par (I, 5.4,
(v)). Le théorème dans ces trois cas sera donc encore conséquence d'un
résultat général sur ce type de propriétés ne concernant que les fibres, qui
sera établi dans (9.3.3)~; nous repoussons donc jusqu'à ce moment-là la
démonstration du théorème dans le cas (xi) (bien entendu, le lecteur pourra
vérifier que, sauf aux n\textsuperscript{os}~8.11 et 8.12, nous ne ferons pas
usage du théorème dans ce cas jusqu'à (9.3.3) et que (8.11) et (8.12) ne
seront pas utilisés avant (9.3.3)).

Pour les cas qui restent à démontrer, on peut se borner à prouver que la
condition de l'énoncé est nécessaire, toutes les propriétés $\mathbf P$
envisagées étant invariantes par changement de base (voir chap. I\textsuperscript{er}
et II dans les numéros concernant chacune de ces propriétés). On peut en
outre supposer que $S_\alpha=S_0$ et que $Y_\alpha=S_\alpha$, donc
$Y_\lambda=S_\lambda$ pour tout $\lambda\geq\alpha$. Enfin, les propriétés
(i) à (xii) sont locales sur $S_0$, donc, comme $S_0$ est réunion finie
d'ouverts affines et $L$ filtrant, on peut se borner pour les démontrer au cas
où $S_0=\operatorname{Spec}(A_0)$ est affine (donc quasi-séparé). On désignera
par $A_\lambda$ (resp. $A$) l'anneau de $S_\lambda$ (resp. $S$).

\emph{Cas (ii), (iii), (iv)}~: Supposons que $f$ soit une immersion (resp.
une immersion ouverte, resp. une immersion fermée), et soit $X'$ le
sous-préschéma (resp. induit sur un ouvert, resp. fermé) de $S$ associé à $f$,
qui est donc un $S$-préschéma de présentation finie.
\oldpage[IV-3]{39}
En vertu de (8.6.3), il existe donc un $\lambda\geq\alpha$ et un
sous-préschéma (resp. induit sur un ouvert, resp. fermé) $X'_\lambda$ de
$S_\lambda$, de présentation finie sur $S_\lambda$, tels que $X'$ soit
isomorphe à $X'_\lambda\times_{S_\lambda}S$. Pour tout $\mu\geq\lambda$,
$X'_\mu=X'_\lambda\times_{S_\lambda}S_\mu$ est donc un sous-préschéma (resp.
induit sur un ouvert, resp. fermé) de $S_\mu$, de présentation finie sur
$S_\mu$, et il résulte donc de (8.8.2.4) et (8.8.2.5) qu'il existe un
$\mu\geq\lambda$ et un isomorphisme $g_\mu:X_\mu\to X'_\mu$ tel que
$g=g_\mu\times1_S$ soit l'isomorphisme $X\to X'$ associé à $f$~; d'où la
conclusion.

\emph{Cas (vi) et (vii)}~: On sait (1.8.4.1) que
$Z_\alpha=f_\alpha(X_\alpha)$ est constructible dans $S_\alpha$~; si l'on
pose $Z_\lambda=u_{\alpha\lambda}^{-1}(Z_\alpha)$ pour $\lambda\geq\alpha$
et $Z=u_\alpha^{-1}(Z_\alpha)$, on a $Z_\lambda=f_\lambda(X_\lambda)$ et
$Z=f(X)$ (I, 3.4.8). Comme en vertu de (8.3.11) l'application canonique
$\varinjlim\mathfrak C(S_\lambda)\to\mathfrak C(S)$ est injective, la
relation $f(X)=S$ implique l'existence d'un $\lambda\geq\alpha$ tel que
$f_\lambda(X_\lambda)=S_\lambda$, ce qui prouve le théorème dans le cas (vi).
Pour le prouver dans le cas (vii), il suffit de noter que le morphisme
structural $X_\alpha\times_{S_\alpha}X_\alpha\to S_\alpha$ est de présentation
finie puisqu'il en est ainsi de la première projection
$X_\alpha\times_{S_\alpha}X_\alpha\to X_\alpha$ (1.6.2)~; il suffit donc en
vertu de (1.8.7.1) d'appliquer le cas (vi) du théorème au morphisme diagonal
$\Delta_{f_\alpha}:X_\alpha\to X_\alpha\times_{S_\alpha}X_\alpha$ en notant
que l'on a $\Delta_{f_\lambda}=\Delta_{f_\alpha}\times1_{S_\lambda}$ et
$\Delta_f=\Delta_{f_\alpha}\times1_S$ (I, 5.3.4 et 3.3.11).

\emph{Cas (viii) et (ix)}~: Comme $S=\operatorname{Spec}(A)$ est affine, dire
que $f:X\to S$ est affine (resp. quasi-affine) signifie qu'il existe un entier
$r$ et une immersion fermée (resp. une immersion)
$j:X\to\mathbf V_S^r=\operatorname{Spec}(A[T_1,\ldots,T_r])$ de
$S$-préschémas, puisque $f$ est de type fini et $S$ quasi-compact
(II, 5.1.9). Comme
$\mathbf V_S^r=\mathbf V_{S_0}^r\times_{S_0}S$, et que
$\mathbf V_{S_0}^r$ est un $S_0$-préschéma de présentation finie, il résulte
de (8.8.2, (i)) appliqué au $S$-morphisme $j$, qu'il existe un $\lambda$ et un
$S_\lambda$-morphisme $j_\lambda:X_\lambda\to\mathbf V_{S_\lambda}^r$ tel que
$j=j_\lambda\times1_S$~; appliquant alors (iv) (resp. (ii)) à $j$, on en
déduit qu'il existe $\mu\geq\lambda$ tel que $j_\mu$ soit une immersion
fermée (resp. une immersion)~; par suite $f_\mu$ est affine (resp.
quasi-affine).

\emph{Cas (x)}~: Par hypothèse, on a $X=\operatorname{Spec}(B)$, où $B$ est
une $A$-algèbre qui est un $A$-module de présentation finie (1.4.7)~; il
résulte donc de (8.5.2, (ii)) qu'il y a un $\lambda$ et un $A_\lambda$-module
de présentation finie $B_\lambda$ tels que le $A$-module $B$ soit isomorphe à
$B_\lambda\otimes_{A_\lambda}A$. La structure de $A$-algèbre de $B$ est
définie par un $A$-homomorphisme $m:B\otimes_AB\to B$~; comme on a
$B\otimes_AB=(B_\lambda\otimes_{A_\lambda}B_\lambda)\otimes_AA$, il existe
d'après (8.5.2, (i)) un $\mu\geq\lambda$ et un $A_\mu$-homomorphisme
$m_\mu:B_\mu\otimes_{A_\mu}B_\mu\to B_\mu$ tel que $m=m_\mu\otimes1$. En
considérant les diagrammes usuels exprimant l'associativité et la
commutativité de $m$, on voit en appliquant encore (8.5.2, (i)) qu'il existe
$\nu\geq\mu$ tel que $m_\nu$ définisse sur $B_\nu$ une multiplication
associative et commutative~; de la même manière on voit qu'on peut supposer
$\nu$ pris assez grand pour que l'anneau $B_\nu$ ainsi défini admette un
élément unité. Si $X'_\nu=\operatorname{Spec}(B_\nu)$, il est clair alors que
$X$ est $S$-isomorphe à $X'_\nu\times_{S_\nu}S$, donc, en vertu de (i), il
existe $\rho\geq\nu$ tel que $X_\rho$ et $X'_\rho$ soient
$S_\rho$-isomorphes, ce qui achève la démonstration dans ce cas.

Pour démontrer le théorème dans le cas (xii), nous démontrerons d'abord la
proposition suivante~:
\end{proof}

\begin{proposition}[8.10.5.1]
\label{IV.8.10.5.1-fr}
(lemme de Chow pour les morphismes de présentation finie). --- Soient $A$ un
anneau, $X$, $Y$ deux $A$-préschémas de présentation finie,
$f:X\to Y$
\oldpage[IV-3]{40}
un $A$-morphisme séparé. Alors il existe deux $A$-préschémas $X'$, $P$ de
présentation finie, et des $A$-morphismes $p:P\to Y$, $j:X'\to P$,
$g:X'\to X$, tels que le diagramme
\[
\begin{tikzcd}[column sep=large,row sep=large]
X' \arrow[r,"j"] \arrow[d,"g"'] & P \arrow[d,"p"] \\
X \arrow[r,"f"'] & Y
\end{tikzcd}
\]
soit commutatif, et que~: 1\textsuperscript{o} $p$ soit projectif~;
2\textsuperscript{o} $g$ soit projectif et surjectif~;
3\textsuperscript{o} $j$ soit une immersion ouverte.
\end{proposition}

\begin{proof}
En effet, soient $A_0\subset A$, $X_0$, $Y_0$ et $f_0$ déterminés comme dans
(8.9.1) de sorte que $Y_0$ est noethérien et $f_0$ de type fini~; on peut en
outre supposer $f_0$ séparé d'après (8.10.4). Le lemme de Chow (II, 5.6.1)
montre alors l'existence de trois morphismes $p_0:P_0\to Y_0$,
$g_0:X'_0\to X_0$ et $j_0:X'_0\to P_0$, de type fini, tels que le diagramme
\[
\begin{tikzcd}[column sep=large,row sep=large]
X'_0 \arrow[r,"j_0"] \arrow[d,"g_0"'] & P_0 \arrow[d,"p_0"] \\
X_0 \arrow[r,"f_0"'] & Y_0
\end{tikzcd}
\]
soit commutatif, et que $p_0$ soit projectif, $g_0$ projectif et surjectif et
$j_0$ une immersion ouverte. Les propriétés de l'énoncé résultent alors de
l'invariance des propriétés précédentes par changement de base (II, 5.5.5,
(iii) et I, 3.5.2 et 4.3.2).
\end{proof}

\begin{proof}[Suite de la démonstration de (8.10.5)]
\emph{Cas (xii)}~: Appliquons au morphisme $f_0:X_0\to S_0$ la prop.
(8.10.5.1)~: on a donc un diagramme commutatif
\[
\begin{tikzcd}[column sep=large,row sep=large]
X'_0 \arrow[r,"j_0"] \arrow[d,"g_0"'] & P_0 \arrow[d,"p_0"] \\
X_0 \arrow[r,"f_0"'] & S_0
\end{tikzcd}
\]
où $p_0$ est projectif, $g_0$ projectif et surjectif, et $j_0$ une immersion
ouverte~; on en déduit pour chaque $\lambda$ un diagramme analogue où les
morphismes $p_\lambda=p_0\times1_{S_\lambda}$,
$g_\lambda=g_0\times1_{S_\lambda}$ et $j_\lambda=j_0\times1_{S_\lambda}$ ont
respectivement les mêmes propriétés, et de même pour les morphismes limites
projectives $p=p_0\times1_S$, $g=g_0\times1_S$, $j=j_0\times1_S$. Comme $g$
est propre (II, 5.5.3), il en est de même de $p\circ j=f\circ g$ (II, 5.4.2)
et comme $p$ est séparé, $j$ est propre, donc une immersion fermée~;
appliquant le cas (iv) au morphisme $j$ (en notant que $X'_0$ et $P_0$ sont
des $S_0$-préschémas de présentation finie (8.10.5.1 et 1.6.2)), on voit
qu'il existe $\lambda$ tel que $j_\lambda$ soit une immersion fermée, donc
soit propre (II, 5.4.2). Mais alors
$f_\lambda\circ g_\lambda=j_\lambda\circ p_\lambda$ est propre (II, 5.5.3 et
5.4.2) et comme $g_\lambda$ est surjectif, et qu'on peut supposer $f_\lambda$
séparé en vertu de l'hypothèse sur $f$ et de (v), il résulte de (II, 5.4.3)
que $f_\lambda$ est propre.

\emph{Cas (xiii) et (xiv)}~: En vertu de (xii) et de (II, 5.5.3) (qui est
applicable puisque les $S_\lambda$ sont quasi-compacts et quasi-séparés,
compte tenu de (1.7.19)), il suffit de
\oldpage[IV-3]{41}
considérer le cas où $f$ est quasi-projectif. Supposons donc qu'il existe un
$\mathcal O_X$-Module inversible $\mathcal L$ qui soit $f$-ample~; comme $S_0$
est quasi-compact et quasi-séparé il en est de même de $X_0$ (1.2.3), et il y
a donc un $\lambda$ et un $\mathcal O_{X_\lambda}$-Module quasi-cohérent
$\mathcal L_\lambda$ de présentation finie tel que
$\mathcal L=v_\lambda^*(\mathcal L_\lambda)$ (8.5.2, (ii))~; en outre, en
vertu de (8.5.5) on peut supposer $\mathcal L_\lambda$ inversible. Le théorème
dans ce cas est alors conséquence du lemme plus précis suivant.

\begin{lemma}[8.10.5.2]
\label{IV.8.10.5.2-fr}
Supposons $S_0$ quasi-compact, et soit $\mathcal L_\lambda$ un
$\mathcal O_{X_\lambda}$-Module inversible. Pour que $\mathcal L$ soit un
$\mathcal O_X$-Module ample pour $f$ (resp. très ample pour $f$), il faut et
il suffit qu'il existe $\mu\geq\lambda$ tel que $\mathcal L_\mu$ soit ample
pour $f_\mu$ (resp. très ample pour $f_\mu$).
\end{lemma}

La condition étant évidemment suffisante (II, 4.4.10 et 4.6.13), montrons
qu'elle est nécessaire~; les $S_\lambda$ étant quasi-compacts et les
$f_\lambda$ de type fini, on peut, en remplaçant $\mathcal L$ par une puissance
tensorielle convenable, se borner à considérer le cas où $\mathcal L$ est très
ample (II, 4.6.11). En outre, la question étant ici locale sur $S_0$ (vu
(II, 4.4.5) et le fait que $L$ est filtrant), on peut se borner au cas où
$S_0$ (et par suite $S$) est affine. Alors, en vertu de (II, 4.4.1, (ii) et
II, 4.1.2), il existe une $S$-immersion
$j:X\to\mathbf P_S^r=P$ telle que $\mathcal L$ soit isomorphe à
$j^*(\mathcal O_P(1))$. Compte tenu de (8.8.2, (i)), de (ii) et de
(II, 4.1.3), il existe donc un $\mu\geq\lambda$ et une immersion
$j_\mu:X_\mu\to\mathbf P_{S_\mu}^r=P_\mu$ tels que
$j=j_\mu\times1_S$~; utilisant ensuite (II, 4.1.3.2) et (8.5.2.5), on voit
qu'il existe $\nu\geq\mu$ tel que $\mathcal L_\nu$ soit isomorphe à
$j_\nu^*(\mathcal O_{P_\nu}(1))$, ce qui montre que $\mathcal L_\nu$ est très
ample pour $f_\nu$ (II, 4.4.2).
\end{proof}

\subsection{Application aux morphismes quasi-finis}
\label{subsection:IV.8.11-fr}
\label{IV.8.11-fr}

Nous nous proposons dans cette section de démontrer les deux théorèmes
suivants~:

\begin{theorem}[8.11.1]
\label{IV.8.11.1-fr}
Soit $f:X\to Y$ un morphisme propre, localement de présentation finie, et
quasi-fini. Alors le morphisme $f$ est fini.
\end{theorem}

\begin{theorem}[8.11.2]
\label{IV.8.11.2-fr}
Soit $f:X\to Y$ un morphisme localement de présentation finie, quasi-fini et
séparé. Alors le morphisme $f$ est quasi-affine, et \emph{a fortiori}
quasi-projectif.
\end{theorem}

\begin{remark}[8.11.3]
\label{IV.8.11.3-fr}
On va voir ci-dessous que l'on peut se ramener, pour la démonstration de
(8.11.1) et (8.11.2), au cas où $Y$ est localement noethérien~; on notera que
dans ce cas on obtiendra ainsi une autre démonstration du théorème de
Chevalley (III, 4.4.2).
\end{remark}

\begin{env}[8.11.4]
\label{IV.8.11.4-fr}
Les hypothèses et les conclusions de (8.11.1) et (8.11.2) sont toutes locales
sur $Y$ (I, 6.1, I, 2.6, (II, 5.1.1), (II, 5.4.1) et (II, 6.2.2)), donc on
peut supposer $Y=\operatorname{Spec}(A)$ affine. On sait qu'il existe alors un
sous-anneau $A_0$ de $A$, qui est une $\mathbf Z$-algèbre de type fini, et un
morphisme de type fini $f_0:X_0\to\operatorname{Spec}(A_0)$ tel que $X$
s'identifie à $X_0\otimes_{A_0}A$ et $f$ à $f_0\times1$ (8.9.1). En outre,
$A$ peut être considéré comme limite inductive de ses sous-anneaux contenant
$A_0$ et qui sont des $\mathbf Z$-algèbres de type fini~; utilisant la méthode
de (8.1.2, c)) ainsi que (8.10.5, (v), (xi) et (xii)), on voit qu'il suffit de
démontrer les théorèmes (8.11.1) et (8.11.2) pour $f_0$. Supposons donc
désormais $Y$ noethérien~; utilisant maintenant la méthode de (8.1.2, a))
ainsi que (8.10.5, (v), (ix), (x), (xi) et (xii)) on peut remplacer $Y$ par
$\operatorname{Spec}(\mathcal O_y)$, où $y$ est un
\oldpage[IV-3]{42}
point quelconque de $Y$, donc on voit finalement que l'on peut supposer que
$Y=\operatorname{Spec}(A)$, où $A$ est un \emph{anneau local noethérien}.
Soient $\mathfrak m$ l'idéal maximal de $A$, $\widehat A$ le complété de $A$
pour la topologie $\mathfrak m$-préadique~; on sait que $\widehat A$ est un
anneau local noethérien et que le morphisme
$\operatorname{Spec}(\widehat A)\to\operatorname{Spec}(A)$ est fidèlement
plat et quasi-compact (0\textsubscript{I}, 7.3.5)~; quant à (2.7.1, (i),
(vii), (xiv), (xv) et (xvi)), on voit en outre qu'on peut se borner au cas où
$A$ est \emph{complet}. Il résulte alors de (II, 6.2.6) que
$X=X'\amalg X''$, où $X'$ est un $Y$-schéma fini et $X''$ un $Y$-schéma
quasi-fini tel que $X''\cap f^{-1}(y)=\varnothing$.
\end{env}

Plaçons-nous d'abord dans les hypothèses de (8.11.1)~; comme $f$ est propre,
$X''$, qui est fermé dans $X$, est propre sur $Y$ (II, 5.4.10), donc
$f(X'')$ est fermé dans $Y$~; mais $y$ n'est pas contenu dans $f(X'')$, et
par ailleurs est adhérent à tout point de $Y$, donc $f(X'')=\varnothing$, et
par suite $X''$ est vide et $f$ est \emph{fini}.

Plaçons-nous maintenant dans les hypothèses de (8.11.2) et, en nous ramenant
encore (comme on peut le faire d'après ce qui précède) au cas où
$Y=\operatorname{Spec}(A)$ est affine et noethérien de dimension finie,
raisonnons en outre par récurrence sur la dimension de $Y$. En se ramenant
comme ci-dessus au cas où $A$ est en outre local et complet, on a
$\dim(\mathcal O_y)=\dim(A)=\dim(Y)$ et pour tout $z\neq y$,
$\dim(\mathcal O_z)<\dim(\mathcal O_y)$, donc
$\dim(Y-\{y\})<\dim(Y)$. Or, par hypothèse on a
$f(X'')\subset Y-\{y\}$ et la restriction de $f$ à $X''$ est évidemment un
morphisme quasi-fini et séparé~; appliquant à $Y-\{y\}$ et à $X''$
l'hypothèse de récurrence, on voit que $X''$ est quasi-affine sur
$Y-\{y\}$~; mais $Y-\{y\}$ étant quasi-affine sur $Y$ puisque $Y$ est
noethérien, $X''$ est aussi quasi-affine sur $Y$ (II, 5.1.10, (ii))~; comme
par ailleurs $X'$ est fini (et \emph{a fortiori} affine) sur $Y$, $X$ est
quasi-affine sur $Y$ (II, 4.6.17 et 5.1.2, c')).

\begin{proposition}[8.11.5]
\label{IV.8.11.5-fr}
Soit $f:X\to Y$ un morphisme de présentation finie. Les conditions suivantes
sont équivalentes~:
\begin{enumerate}
\item[\emph{a)}] $f$ est une immersion fermée.
\item[\emph{b)}] $f$ est un monomorphisme propre.
\item[\emph{c)}] $f$ est propre et pour tout $y\in Y$, $f^{-1}(y)$ est
radiciel et géométriquement réduit sur $k(y)$ (c'est-à-dire vide ou
$k(y)$-isomorphe à $\operatorname{Spec}(k(y))$).
\end{enumerate}
\end{proposition}

Il est clair que \emph{a)} entraîne \emph{b)}. Pour voir que \emph{b)}
entraîne \emph{c)}, notons (I, 5.3.8) que pour tout $y\in Y$, le morphisme
$f^{-1}(y)\to\operatorname{Spec}(k(y))$ déduit de $f$ par extension de la base
est un monomorphisme, donc est injectif, et par suite $f^{-1}(y)$ est vide ou
réduit à un point, et en tout cas affine. De plus, si $A$ est l'anneau de
$f^{-1}(y)$, l'homomorphisme canonique $A\otimes_kA\to A$ est bijectif
(I, 5.3.8). Cela entraîne évidemment que $A=k$, sans quoi il y aurait un
élément $a\in A$ non dans $k$ et les images de $a\otimes1$ et $1\otimes a$
dans $A$ seraient toutes deux égales à $a$, alors que $a\otimes1\neq1\otimes a$
puisque $1$ et $a$ forment un système libre sur $k$.

Reste à prouver que \emph{c)} entraîne \emph{a)}. Il résulte tout d'abord de
(8.11.1) que $f$ est un morphisme \emph{fini}, donc
$X=\operatorname{Spec}(\mathcal A)$, où $\mathcal A$ est une
$\mathcal O_Y$-Algèbre \emph{finie}. Il suffit donc de prouver que
l'homomorphisme canonique $\mathcal O_Y\to\mathcal A$ est \emph{surjectif}
(II, 1.4.10), ou encore que pour tout $y\in Y$, l'homomorphisme
$\mathcal O_y\to\mathcal A_y$ est surjectif. Mais par hypothèse
\oldpage[IV-3]{43}
$f^{-1}(y)=\operatorname{Spec}(\mathcal A_y/\mathfrak m_y\mathcal A_y)$
(II, 1.5.5) est tel que l'homomorphisme correspondant
$k(y)=\mathcal O_y/\mathfrak m_y\to
\mathcal A_y/\mathfrak m_y\mathcal A_y$ soit bijectif~; comme $\mathcal A_y$
est un $\mathcal O_y$-module de type fini, le lemme de Nakayama montre que
$\mathcal O_y\to\mathcal A_y$ est surjectif, ce qui achève la démonstration.

\begin{remark}[8.11.5.1]
\label{IV.8.11.5.1-fr}
On notera que le raisonnement précédent prouve que si $f$ est un
\emph{monomorphisme}, alors, pour tout $y\in Y$, $f^{-1}(y)$ est vide ou
$k(y)$-isomorphe à $\operatorname{Spec}(k(y))$.
\end{remark}

\begin{proposition}[8.11.6]
\label{IV.8.11.6-fr}
Si un morphisme $f:X\to Y$ de présentation finie est un homéomorphisme
universel, il est fini, surjectif et radiciel (la réciproque étant vraie
d'après (2.4.5, (i))).
\end{proposition}

\begin{proof}
En effet, $f$ étant de type fini, universellement fermé, et séparé en vertu de
(2.4.4), est propre par définition (II, 5.4.1). Comme il est évidemment
quasi-fini (II, 6.2.3), il est fini d'après (8.11.1). On sait par ailleurs
qu'il est radiciel (2.4.4), et évidemment surjectif.
\end{proof}

\subsection{Nouvelle démonstration et généralisation du «~Main Theorem~» de Zariski}
\label{subsection:IV.8.12-fr}
\label{IV.8.12-fr}

\begin{lemma}[8.12.1]
\label{IV.8.12.1-fr}
Soient $f:X\to Y$ un morphisme quasi-compact et quasi-séparé, $\mathcal C$ une
$\mathcal O_Y$-Algèbre quasi-cohérente, $Z=\operatorname{Spec}(\mathcal C)$,
qui est un $Y$-préschéma affine au-dessus de $Y$. Soient $g:X\to Z$ un
$Y$-morphisme, $\varphi=\mathcal A(g):\mathcal C\to f_*(\mathcal O_X)$ le
$\mathcal O_Y$-homomorphisme correspondant de $\mathcal O_Y$-Algèbres
(II, 1.2.7). Supposons que $g$ soit une immersion. Alors, pour que l'image
fermée de $X$ par $g$ (I, 9.5.3) soit égale à $Z$, il faut et il suffit que
$\varphi$ soit injectif~; $g$ est alors une immersion ouverte.
\end{lemma}

\begin{proof}
L'hypothèse entraîne que $f_*(\mathcal O_X)$ est une $\mathcal O_Y$-Algèbre
quasi-cohérente (1.7.5)~; en outre, comme le morphisme canonique $h:Z\to Y$
est affine, donc quasi-compact et quasi-séparé (1.2.2 et 1.1.2), $g$ est un
morphisme quasi-compact et quasi-séparé, donc $g_*(\mathcal O_X)$ est une
$\mathcal O_Z$-Algèbre quasi-cohérente (1.7.5). Cela étant, dire que l'image
fermée de $X$ par $g$ est égale à $Z$ signifie (I, 9.5.2) que
l'homomorphisme canonique $\theta:\mathcal O_Z\to g_*(\mathcal O_X)$ est
injectif. Mais on a $h_*(\mathcal O_Z)=\mathcal C$ par définition de $Z$
(II, 1.3.1), et $h_*(g_*(\mathcal O_X))=f_*(\mathcal O_X)$. Comme $Z$ est
affine au-dessus de $Y$, il revient au même de dire que l'homomorphisme
$\theta:\mathcal O_Z\to g_*(\mathcal O_X)$ est injectif ou que
l'homomorphisme correspondant
$\varphi=h_*(\theta):\mathcal C\to f_*(\mathcal O_X)$ est injectif
(I, 1.3.9). Le fait que $g$ est alors une immersion ouverte résulte de
(I, 9.5.10) et de l'hypothèse que $g$ est une immersion.
\end{proof}

\begin{lemma}[8.12.2]
\label{IV.8.12.2-fr}
Soient $Y$ un préschéma quasi-compact et quasi-séparé, $f:X\to Y$ un
morphisme quasi-séparé de type fini, $\mathcal C$ une $\mathcal O_Y$-Algèbre
quasi-cohérente, $Z=\operatorname{Spec}(\mathcal C)$. Soient $g:X\to Z$ un
$Y$-morphisme, $\varphi:\mathcal C\to f_*(\mathcal O_X)$ le
$\mathcal O_Y$-homomorphisme correspondant de $\mathcal O_Y$-Algèbres. Soit
$(\mathcal C_\lambda)$ la famille filtrante croissante des sous-$\mathcal O_Y$-Algèbres
quasi-cohérentes de type fini de $\mathcal C$ (dont $\mathcal C$ est la
réunion ((I, 9.6.6) et (1.7.9)))~; posons
$Z_\lambda=\operatorname{Spec}(\mathcal C_\lambda)$ et soit $g_\lambda$ le
morphisme composé $X\xrightarrow{g}Z\to Z_\lambda$. Alors les conditions
suivantes sont équivalentes~:
\begin{enumerate}
\item[\emph{a)}] $g$ est une immersion.
\item[\emph{b)}] Il existe $\lambda$ tel que $g_\lambda$ soit une immersion.
\end{enumerate}
\oldpage[IV-3]{44}
De plus, lorsque $g_\lambda$ est une immersion, il en est de même de $g_\mu$
pour $\mu\geq\lambda$.

Il suffit d'appliquer (II, 3.8.4) en remplaçant $\mathcal L$ par
$\mathcal O_X$ et $\mathcal S$ par $\mathcal C[\mathrm T]$, et en tenant
compte de (II, 3.1.7).
\end{lemma}

\begin{proposition}[8.12.3]
\label{IV.8.12.3-fr}
Soient $Y$ un préschéma quasi-compact et quasi-séparé, $f:X\to Y$ un
morphisme séparé de type fini. Soit $\mathcal B=f_*(\mathcal O_X)$, qui est
une $\mathcal O_Y$-Algèbre quasi-cohérente (I, 9.2.2)~; soit $\mathcal C$ la
fermeture intégrale de $\mathcal O_Y$ dans $\mathcal B$, qui est une
$\mathcal O_Y$-Algèbre quasi-cohérente (II, 6.3.4)~; posons
$Z=\operatorname{Spec}(\mathcal C)$, et soit $g:X\to Z$ le $Y$-morphisme
correspondant à l'injection
$\varphi:\mathcal C\to\mathcal B=f_*(\mathcal O_X)$ (II, 1.2.7). Soit
$(\mathcal C_\lambda)$ la famille filtrante croissante des
sous-$\mathcal O_Y$-Algèbres quasi-cohérentes de type fini de $\mathcal C$
(dont $\mathcal C$ est la réunion ((I, 9.6.6) et (1.7.9))), et, pour tout
$\lambda$, posons $Z_\lambda=\operatorname{Spec}(\mathcal C_\lambda)$, et
soit $g_\lambda:X\to Z_\lambda$ le $Y$-morphisme correspondant à l'injection
$\mathcal C_\lambda\to\mathcal B=f_*(\mathcal O_X)$~; les conditions
suivantes sont équivalentes~:
\begin{enumerate}
\item[\emph{a)}] Il existe une factorisation de $f$ en
\[
X\xrightarrow{f'}Y'\xrightarrow{u}Y,
\]
où $f'$ est une immersion et $u$ un morphisme fini.
\item[\emph{a$'$)}] Il existe une factorisation
$X\xrightarrow{f'}Y'\xrightarrow{u}Y$ de $f$, où $f'$ est une immersion
ouverte et $u$ un morphisme fini.
\item[\emph{b)}] Le morphisme $g:X\to Z$ est une immersion.
\item[\emph{c)}] Il existe $\lambda$ tel que
$g_\lambda:X\to Z_\lambda$ soit une immersion.
\end{enumerate}
De plus, lorsqu'il en est ainsi, $g$ est une immersion ouverte, $g(X)$ est
dense dans $Z$, et il existe $\lambda$ tel que, pour $\mu\geq\lambda$,
$g_\mu$ soit une immersion ouverte.
\end{proposition}

Comme l'homomorphisme $\varphi:\mathcal C\to f_*(\mathcal O_X)$ est
injectif, il résulte de (8.12.1) que si $g$ est une immersion, c'est une
immersion ouverte et $g(X)$ est dense dans $Z$, et il en est de même pour
$g_\lambda$. Le fait que \emph{a)} entraîne \emph{a$'$)} résulte aussi de
(8.12.1), appliqué en remplaçant $Z$ par $Y'$ et $g$ par $f'$ ($Y'$ étant
fini, donc affine au-dessus de $Y$)~: en effet si $Y''$ est l'image fermée de
$X$ par $f'$, $Y''$ est fini au-dessus de $Y$ et $f'$ se factorise en
$X\xrightarrow{f''}Y''\xrightarrow{j}Y'$, où $j$ est l'injection canonique,
et $f''$ est une immersion (I, 4.1.10)~; il résulte donc de (8.12.1) que
$f''$ est une immersion ouverte.

L'équivalence de \emph{b)} et \emph{c)} résulte de (8.12.2) ainsi que le fait
que $g_\lambda$ est alors une immersion pour $\lambda$ assez grand. Il est
clair que \emph{c)} implique \emph{a)}, puisque $Z_\lambda$ est fini
au-dessus de $Y$ (II, 6.3.4 et 6.1.2). Montrons enfin que \emph{a)} implique
\emph{c)}. On a vu ci-dessus qu'on peut supposer que $Y'$ est l'image fermée
de $X$ par $f'$, et alors il résulte de (8.12.1) que si l'on pose
$\mathcal B'=u_*(\mathcal O_{Y'})$, de sorte que $Y'$ s'identifie à
$\operatorname{Spec}(\mathcal B')$, l'homomorphisme
$\varphi':\mathcal B'\to\mathcal B=f_*(\mathcal O_X)$ est \emph{injectif}.
Mais comme par hypothèse $\mathcal B'$ est une $\mathcal O_Y$-Algèbre finie,
elle s'identifie par définition de $\mathcal B$ à une des
sous-$\mathcal O_Y$-Algèbres $\mathcal C_\lambda$, ce qui démontre
\emph{c)}.

Nous dirons qu'un morphisme $f:X\to Y$, où $Y$ est quasi-compact et
quasi-séparé, est \emph{pseudo-fini}, s'il est de type fini et vérifie la
condition \emph{a)} de (8.12.3) (auquel cas il est nécessairement séparé).

\begin{corollary}[8.12.4]
\label{IV.8.12.4-fr}
Soient $Y$ un préschéma quasi-compact et quasi-séparé, $f:X\to Y$ un
morphisme.
\oldpage[IV-3]{45}
\begin{enumerate}
\item[(i)] Supposons $f$ pseudo-fini. Alors, pour tout morphisme $Y'\to Y$,
où $Y'$ est quasi-compact et quasi-séparé,
$f_{(Y')}:X'=X_{(Y')}\to Y'$ est pseudo-fini.
\item[(ii)] Soit $(U_\lambda)$ un recouvrement de $Y$ formé d'ouverts
quasi-compacts. Pour que $f$ soit pseudo-fini, il faut et il suffit que pour
tout $\lambda$, la restriction
$f_\lambda:f^{-1}(U_\lambda)\to U_\lambda$ de $f$ soit un morphisme
pseudo-fini.
\item[(iii)] Supposons en outre $Y$ noethérien, et $f$ de type fini. Alors,
pour que $f$ soit pseudo-fini, il faut et il suffit que, pour tout $y\in Y$,
le morphisme
$f_y=f\times 1:X_y=X\times_Y\operatorname{Spec}(\mathcal O_y)
\to\operatorname{Spec}(\mathcal O_y)$ le soit.
\end{enumerate}
\end{corollary}

\begin{proof}
(i) Il est clair que $f_{(Y')}$ est de type fini (1.5.4)~; en outre, une
factorisation
$X\xrightarrow{g}Z\xrightarrow{u}Y$, où $g$ est une immersion et $u$ est
fini, donne une factorisation
$X'\xrightarrow{g'}Z'\xrightarrow{u'}Y'$ de $f_{(Y')}$, où
$Z'=Z_{(Y')}$, $g'=g_{(Y')}$ et $u'=u_{(Y')}$~; $g'$ est une immersion
(I, 4.3.2) et $u'$ est fini (II, 6.1.5)~; donc $f_{(Y')}$ est pseudo-fini.

(ii) La condition est nécessaire en vertu de (i), les $U_\lambda$ étant
quasi-séparés puisque $Y$ l'est. Pour voir qu'elle est suffisante, observons
(avec les notations de (8.12.3)) que si l'on pose
$X_\lambda=f^{-1}(U_\lambda)$, on a
$\mathcal B|U_\lambda=(f_\lambda)_*(\mathcal O_{X_\lambda})$,
$\mathcal C|U_\lambda$ est la fermeture intégrale de $\mathcal O_{U_\lambda}$
dans $\mathcal B|U_\lambda$ et par suite, si $h:Z\to Y$ est le morphisme
canonique, $Z_\lambda=\operatorname{Spec}(\mathcal C|U_\lambda)$ s'identifie à
$h^{-1}(U_\lambda)$. Or, pour que $g:X\to Z$ soit une immersion, il faut et
il suffit que pour tout $\lambda$, la restriction
$g_\lambda:f^{-1}(U_\lambda)\to h^{-1}(U_\lambda)$ de $g$ le soit
(I, 4.2.4). Cela entraîne la conclusion en vertu de (8.12.3).

(iii) Il suffit, en vertu de (ii), de prouver que $y$ admet un voisinage $U$
tel que la restriction $f^{-1}(U)\to U$ de $f$ soit un morphisme pseudo-fini.
Désignons par $(U_\lambda)$ le système projectif filtrant décroissant des
voisinages ouverts affines de $y$, et appliquons la méthode de (8.1.2,
\emph{a)}). Puisque $Y$ est noethérien, les restrictions
$f_\lambda:f^{-1}(U_\lambda)\to U_\lambda$ de $f$ sont de présentation
finie et il en est de même de $f_y$. Par hypothèse $f_y$ se factorise en
$X_y\xrightarrow{g_y}Z_y\xrightarrow{u_y}\operatorname{Spec}(\mathcal O_y)$,
où $u_y$ est fini et $g_y$ est une immersion. Comme $\mathcal O_y$ est
noethérien, il en est de même de $Z_y$, et comme $Z_y$ est de présentation
finie sur $\operatorname{Spec}(\mathcal O_y)$, il existe un $\lambda$ et un
morphisme de présentation finie $u_\lambda:Z_\lambda\to U_\lambda$ tels que
$Z_y$ s'identifie à
$Z_\lambda\times_{U_\lambda}\operatorname{Spec}(\mathcal O_y)$ et $u_y$ à
$u_\lambda\times 1$ (8.8.2, (ii))~; en outre, il existe un morphisme
$g_\lambda:X_\lambda\to Z_\lambda$ tel que $g_y=g_\lambda\times 1$ et que
$f_\lambda=u_\lambda\circ g_\lambda$ (8.8.2, (i)). De plus, on peut supposer
$\lambda$ pris de sorte que $g_\lambda$ soit une immersion et $u_\lambda$ un
morphisme fini (8.10.5, (ii) et (x)), ce qui prouve que $f_\lambda$ est
pseudo-fini.
\end{proof}

\begin{env}[8.12.5]
\label{IV.8.12.5-fr}
Nous pouvons maintenant donner du «~Main theorem~» de Zariski (III, 4.4.3),
une démonstration n'utilisant pas les résultats cohomologiques de nature
«~globale~» du chap.~III, mais faisant appel par contre aux propriétés plus
fines des anneaux locaux noethériens~; nous généraliserons en outre l'énoncé
du théorème en le débarrassant des hypothèses noethériennes~:
\end{env}

\begin{theorem}[8.12.6]
\label{IV.8.12.6-fr}
(«~Main Theorem~» de Zariski). --- Soit $Y$ un préschéma quasi-compact et
quasi-séparé. Si un morphisme $f:X\to Y$ est quasi-fini, séparé et de
présentation finie, il existe une factorisation de $f$
\begin{equation}
\label{IV.8.12.6.1-fr}
X\xrightarrow{f'}Y'\xrightarrow{u}Y. \tag{8.12.6.1}
\end{equation}
où $f'$ est une immersion ouverte et $u$ un morphisme fini.
\end{theorem}

\oldpage[IV-3]{46}
\begin{proof}
En vertu de (8.12.4, (ii)) et du caractère local (sur $Y$) des notions de
morphisme quasi-fini, séparé et de présentation finie, on peut se borner au
cas où $Y=\operatorname{Spec}(A)$ est affine. Appliquant (8.9.1), on peut
supposer qu'il y a un sous-anneau $A_0$ de $A$, qui est une
$\mathbf Z$-algèbre de type fini, et un $A$-isomorphisme
$X_0\otimes_{A_0}A\xrightarrow{\sim}X$, $f$ s'identifiant par cet
isomorphisme à $f_0\times 1$, où
$f_0:X_0\to\operatorname{Spec}(A_0)$ est un morphisme \emph{de type fini}~;
en outre (8.10.5, (v) et (xii)) on peut supposer que $f_0$ est séparé et
quasi-fini~; si l'on prouve que $f_0$ est pseudo-fini, il en sera de même de
$f$ d'après (8.12.4, (i)). Comme $A_0$ est alors noethérien et que les
notions de morphisme de type fini, séparé et quasi-fini se conservent par
changement de base, il résulte de (8.12.4, (iii)) que l'on peut même supposer
que $A$ est un anneau local, \emph{essentiellement de type fini sur}
$\mathbf Z$ (1.3.8). Posons $n=\dim(A)$, et procédons par récurrence sur $n$~;
pour $n=0$, le théorème est évident, $A$ étant un corps et le morphisme $f$
étant déjà fini (II, 6.2.2). Posons $B=\Gamma(X,\mathcal O_X)$~; désignons
par $C$ la fermeture intégrale de $A$ dans $B$, posons
$Z=\operatorname{Spec}(C)$ et soit $g:X\to Z$ le $Y$-morphisme correspondant
au $A$-homomorphisme canonique injectif $C\to B$~; en vertu de (8.12.3), il
s'agit de montrer que $g$ est une immersion ouverte. Soit $a$ le point fermé
de $Y$, et soit $U=Y-\{a\}$~; $U$ est noethérien et tous ses anneaux locaux
sont essentiellement de type fini sur $\mathbf Z$ et de dimension $<n$~;
compte tenu de l'hypothèse de récurrence, et de (8.12.4, (iii)), on voit que
la restriction $f^{-1}(U)\to U$ de $f$ est un morphisme pseudo-fini. On en
conclut (8.12.3) que, si $h:Z\to Y$ est le morphisme structural, la
restriction $f^{-1}(U)\to h^{-1}(U)$ de $g$ est une immersion ouverte.
Posons $A'=\widehat A$, $Y'=\operatorname{Spec}(A')$,
$X'=X_{(Y')}$, $f'=f_{(Y')}:X'\to Y'$. Comme le morphisme canonique
$u:Y'\to Y$ est plat, il résulte de (2.3.1) que
$B'=\Gamma(X',\mathcal O_{X'})$ s'identifie à la $A'$-algèbre
$B\otimes_A A'$.

D'autre part, comme $A$ est un anneau local \emph{excellent} (7.8.3), le
morphisme $u:Y'\to Y$ est régulier, et \emph{a fortiori} normal, et par suite
(6.14.4) \emph{la fermeture intégrale de $A'$ dans $B'$ est égale à}
$C'=C\otimes_A A'$. On voit donc que $Z'=\operatorname{Spec}(C')$ est égal à
$Z_{(Y')}$ et le morphisme $g':X'\to Z'$ provenant de l'injection $C'\to B'$
est égal à $g_{(Y')}$. Comme $u:Y'\to Y$ est fidèlement plat et
quasi-compact, pour prouver que $g$ est une immersion ouverte, il suffit de
prouver que $g'$ est une immersion ouverte (2.7.1, (x)). Notons maintenant
que $u^{-1}(a)$ est réduit au point fermé $a'$ de $Y'$ et par suite
$U'=Y'-\{a'\}=u^{-1}(U)$. Si $h':Z'\to Y'$ est le morphisme canonique, le
fait que la restriction $f^{-1}(U)\to h^{-1}(U)$ de $g$ soit une immersion
ouverte entraîne qu'il en est de même de la restriction
$f'^{-1}(U')\to h'^{-1}(U')$ de $g'$. Notons maintenant que $f'$ est un
morphisme séparé et quasi-fini (II, 6.2.4)~; comme $A'$ est \emph{complet},
on déduit de (II, 6.2.6) que $X'$ est $Y'$-isomorphe à une somme
$X'_1\amalg X'_2$, où la restriction
$f'|X'_1=f'_1:X'_1\to Y'$ est un morphisme \emph{fini}, et
$X'_2\subset f'^{-1}(U')$. Il en résulte que $B'$ est composée directe des
deux $A'$-algèbres
$\Gamma(X'_1,\mathcal O_{X'_1})=B'_1$ et
$\Gamma(X'_2,\mathcal O_{X'_2})=B'_2$~; on en conclut aussitôt que la
fermeture intégrale $C'$ de $A'$ dans $B'$ est composée directe des
fermetures intégrales $C'_1$, $C'_2$ de $A'$ dans $B'_1$, $B'_2$
respectivement, d'où $Z'=Z'_1\amalg Z'_2$, où
$Z'_i=\operatorname{Spec}(C'_i)$ ($i=1,2$)~; et le morphisme canonique
$g':X'\to Z'$ est tel que $g'|X'_i$ soit le morphisme canonique
$g'_i:X'_i\to Z'_i$ ($i=1,2$). Mais comme $B'_1$ est déjà une
$A'$-algèbre finie, on a $C'_1=B'_1$, et $g'_1$ est donc un isomorphisme.
D'autre part, comme $X'_2\subset f'^{-1}(U')$ et est ouvert dans
$f'^{-1}(U')$, on sait
\oldpage[IV-3]{47}
déjà que $g'_2$ est une immersion ouverte. On conclut bien que $g'$ est une
immersion ouverte.
\end{proof}

\begin{remark}[8.12.7]
\label{IV.8.12.7-fr}
Lorsque, dans (8.12.6), on suppose que $Y$ est un \emph{schéma affine}, la
démonstration par réduction au cas noethérien montre que, dans la
factorisation (8.12.6.1), les morphismes $f'$ et $u$ sont aussi des
morphismes \emph{de présentation finie} (1.6.2).
\end{remark}

\begin{corollary}[8.12.8]
\label{IV.8.12.8-fr}
Soient $Y$ un schéma quasi-compact tel qu'il existe un $\mathcal O_Y$-Module
ample (II, 4.5.3), $f:X\to Y$ un morphisme quasi-fini et quasi-projectif.
Alors il existe une factorisation de $f$ en
\[
X\xrightarrow{f'}Y'\xrightarrow{u}Y
\]
où $f'$ est une immersion ouverte et $u$ un morphisme fini.
\end{corollary}

\begin{proof}
L'hypothèse entraîne que $X$ s'identifie à un sous-$Y$-schéma quasi-compact
d'un $Y$-schéma de la forme $Z=\mathbf P_Y^r$ (II, 5.3.3). Il y a par suite
un voisinage ouvert quasi-compact $U$ de $X$ dans $Z$ tel que $X$ soit fermé
dans $U$~; comme $Z$ est un schéma, l'injection canonique $U\to Z$ est un
morphisme de présentation finie ((1.2.7) et (1.6.2)), donc le morphisme
composé $g:U\to Z\to Y$ est aussi un morphisme de présentation finie (le
fait que $\mathbf P_Y^r$ est de présentation finie sur $Y$ résultant
aussitôt de la définition (II, 4.1.1)). Soit $\mathcal J$ l'Idéal
quasi-cohérent de $\mathcal O_U$ définissant le sous-préschéma fermé $X$~;
comme $U$ est un schéma quasi-compact, $\mathcal J$ est limite inductive
filtrante de ses sous-Idéaux quasi-cohérents de type fini $\mathcal J_\lambda$
(I, 9.4.9). Si $X_\lambda$ est le sous-préschéma fermé de $U$ défini par
$\mathcal J_\lambda$, on a par suite $X=\bigcap_\lambda X_\lambda$. Pour tout
$y\in Y$, on a donc
$f^{-1}(y)=\bigcap_\lambda(X_\lambda\cap g^{-1}(y))$, et comme les ensembles
$X_\lambda\cap g^{-1}(y)$ sont fermés dans l'espace noethérien
$g^{-1}(y)$, il existe pour tout $y$ un indice $\lambda(y)$ tel que
$f^{-1}(y)=X_{\lambda(y)}\cap g^{-1}(y)$. Désignons par $E_\lambda$
l'ensemble des $y\in Y$ tels que la fibre $X_\lambda\cap g^{-1}(y)$ de la
restriction de $g$ à $X_\lambda$ soit un $k(y)$-préschéma fini. L'hypothèse
que $f$ est quasi-fini entraîne, en vertu de ce qui précède, que
$Y=\bigcup_\lambda E_\lambda$. Or, chacun des $X_\lambda$ est, par
définition, \emph{de présentation finie} sur $Y$~; il résulte donc de
(9.2.3) et (9.2.6)\footnote{Le lecteur vérifiera que les corollaires (8.12.8)
à (8.12.11) ne sont pas utilisés dans le \S~9.} que les $E_\lambda$ sont des
ensembles \emph{constructibles} dans le schéma $Y$~; comme ils forment un
ensemble filtrant croissant, il existe un indice $\lambda$ tel que
$E_\lambda=Y$ (1.9.9), et pour cet indice $\lambda$, le morphisme
$f_\lambda:X_\lambda\to Y$, restriction de $g$ à $X_\lambda$, est donc
quasi-fini. Comme il est de présentation finie et séparé, on peut lui
appliquer (8.12.6), et $f_\lambda$ se factorise donc en
\[
X_\lambda\xrightarrow{j_\lambda}Y'_\lambda
\xrightarrow{u_\lambda}Y
\]
où $j_\lambda$ est une immersion et $u_\lambda$ un morphisme fini. Comme
$X$ est un sous-préschéma fermé de $X_\lambda$, on a ainsi prouvé que $f$
possède la propriété (8.12.3, \emph{a)}), d'où le corollaire en vertu de
l'équivalence de (8.12.3, \emph{a)}) et (8.12.3, \emph{a$'$)}).
\end{proof}

\oldpage[IV-3]{48}
\begin{corollary}[8.12.9]
\label{IV.8.12.9-fr}
Soit $f:X\to Y$ un morphisme localement quasi-fini
$(\mathbf{Err}_{\mathrm{III}}, 20)$. Pour tout $x\in X$ il existe un
voisinage ouvert $U$ de $x$ dans $X$, un voisinage ouvert $V$ de
$y=f(x)$ dans $Y$, tels que $f(U)\subset V$ et une factorisation
\[
U\xrightarrow{f'}V'\xrightarrow{u}V
\]
de la restriction de $f$ à $U$, où $f'$ est une immersion ouverte et $u$ un
morphisme fini.
\end{corollary}

Il suffit évidemment de prendre pour $V$ un voisinage affine de $y$ dans
$Y$, pour $U$ un voisinage affine de $x$ dans $X$ contenu dans $f^{-1}(U)$
et tel que $f|U$ soit quasi-fini. Le morphisme $U\to V$ restriction de $f$
étant alors affine (donc quasi-projectif), on peut lui appliquer (8.12.8).

\begin{corollary}[8.12.10]
\label{IV.8.12.10-fr}
Soient $Y$ un préschéma intègre et normal, $X$ un préschéma intègre,
$f:X\to Y$ un morphisme birationnel et localement quasi-fini
$(\mathbf{Err}_{\mathrm{III}}, 20)$. Alors $f$ est un isomorphisme local~;
pour que $f$ soit une immersion ouverte, il faut et il suffit que $f$ soit
en outre séparé.
\end{corollary}

La seconde assertion résulte aussitôt de la première et de (I, 8.2.8). Pour
prouver la première assertion, on peut supposer $X$ et $Y$ affines et $f$
quasi-fini~; considérons la factorisation $f=u\circ f'$ de (8.12.8), qui
permet d'identifier $X$ par $f'$ à un sous-préschéma induit sur un ouvert de
$Y'$. Comme $X$ est intègre, on peut, en vertu de (I, 5.2.3), remplacer
$Y'$ par le sous-préschéma réduit de $Y'$ ayant pour espace sous-jacent
$\overline X$, donc on peut supposer aussi que $Y'$ est \emph{intègre}. En
outre, puisque $f$ est birationnel, $u$ l'est aussi. La conclusion résulte
alors du lemme suivant~:

\begin{lemma}[8.12.10.1]
\label{IV.8.12.10.1-fr}
Soient $Y'$ un préschéma intègre, $Y$ un préschéma intègre et normal~; alors
un morphisme $u:Y'\to Y$ fini et birationnel est un isomorphisme.
\end{lemma}

Posons en effet $\mathcal A=u_*(\mathcal O_{Y'})$, de sorte que $\mathcal A$
est une $\mathcal O_Y$-Algèbre finie, $Y'$ s'identifiant à
$\operatorname{Spec}(\mathcal A)$ (II, 1.3.6). Si $R(Y)$ est le corps des
fonctions rationnelles de $Y$, on a donc, pour tout $y\in Y$,
$\mathcal O_{Y,y}\subset\mathcal A_y\subset R(Y)$~; mais comme l'anneau
$\mathcal O_{Y,y}$ est par hypothèse intégralement clos et a $R(Y)$ pour
corps des fractions, on a nécessairement
$\mathcal A_y=\mathcal O_{Y,y}$, d'où le lemme.

\begin{corollary}[8.12.11]
\label{IV.8.12.11-fr}
Soient $Y$ un préschéma intègre, $X$ un préschéma intègre et normal,
$f:X\to Y$ un morphisme dominant et localement quasi-fini. Soient $K$ et
$L$ (extension de $K$) les corps de fonctions rationnelles de $Y$ et $X$
respectivement~; et soit $Y'$ la fermeture intégrale de $Y$ relativement à
$L$ (II, 6.3.4)~; alors $f$ se factorise d'une seule manière en
$f:X\xrightarrow{f'}Y'\xrightarrow{u}Y$, où $f'$ est birationnel, et
correspond à l'automorphisme identique de $L$~; $f'$ est alors un
isomorphisme local, et pour que $f'$ soit une immersion ouverte, il faut et
il suffit que $f$ soit séparé.
\end{corollary}

L'existence et l'unicité de la factorisation de $f$ résultent de (II, 6.3.9).
Il résulte de (II, 6.2.4, (v)), en se ramenant au cas affine, que $f'$ est
localement quasi-fini~; en outre, il résulte de (I, 5.5.1) que, pour que
$f'$ soit séparé, il faut et il suffit que $f$ le soit, puisque $u$ est
affine, donc séparé~; les deux dernières assertions sont donc des
conséquences de (8.12.10) appliqué à $f'$.

\oldpage[IV-3]{49}
\subsection{Traduction en termes de pro-objets}
\label{subsection:IV.8.13-fr}

La proposition suivante est essentiellement équivalente à (8.8.2, (i))~:

\begin{proposition}[8.13.1]
\label{IV.8.13.1-fr}
Soient $S$ un préschéma, $(X_\lambda,v_{\lambda\mu})$ un système projectif
filtrant de $S$-préschémas~; on suppose qu'il existe $\alpha$ tel que
$v_{\alpha\lambda}$ soit un morphisme affine pour tout $\lambda\geqslant
\alpha$ (ce qui entraîne (II, 1.6.2) que $v_{\lambda\mu}$ soit affine pour
$\alpha\leqslant\lambda\leqslant\mu$), de sorte que la limite projective
$X=\varprojlim_\lambda X_\lambda$ existe dans la catégorie des $S$-préschémas
(8.2.3). Soit $Y$ un $S$-préschéma, et, pour tout $\lambda\geqslant\alpha$,
soit
\[
e_\lambda:\operatorname{Hom}_S(X_\lambda,Y)
\longrightarrow\operatorname{Hom}_S(X,Y)
\]
l'application qui, à tout $S$-morphisme $f_\lambda:X_\lambda\to Y$, fait
correspondre $f=f_\lambda\circ v_\lambda$, où $v_\lambda:X\to X_\lambda$
est le morphisme canonique. La famille $(e_\lambda)$ est un système inductif
d'applications, qui définit donc une application canonique
\begin{equation}
\label{IV.8.13.1.1-fr}
\varinjlim_\lambda\operatorname{Hom}_S(X_\lambda,Y)
\longrightarrow\operatorname{Hom}_S(X,Y). \tag{8.13.1.1}
\end{equation}
Supposons $X_\alpha$ quasi-compact (resp. quasi-compact et quasi-séparé), et le
morphisme structural $Y\to S$ localement de type fini (resp. localement de
présentation finie). Alors l'application (8.13.1.1) est injective (resp.
bijective).
\end{proposition}

Posons en effet, pour $\lambda\geqslant\alpha$,
$Z_\lambda=Y\times_S X_\lambda$, de sorte que l'on a
$Z_\lambda=Z_\alpha\times_{X_\alpha}X_\lambda$. Posons de même
$Z=Y\times_S X=Z_\alpha\times_{X_\alpha}X$~; on sait alors (8.2.5) que, si
l'on pose aussi
\[
w_{\lambda\mu}=1\times v_{\lambda\mu}:Z_\mu\to Z_\lambda
\quad(\alpha\leqslant\lambda\leqslant\mu),
\qquad
w_\lambda=1\times v_\lambda:Z\to Z_\lambda
\quad(\alpha\leqslant\lambda),
\]
$Z$ est limite projective du système projectif $(Z_\lambda,w_{\lambda\mu})$
et les $w_\lambda$ sont les morphismes canoniques correspondants. Notons
d'autre part que le morphisme $Z_\alpha\to X_\alpha$ est localement de type
fini (resp. localement de présentation finie) (1.3.4 et 1.4.3). Enfin, on
sait que l'on a
\[
\operatorname{Hom}_S(X_\lambda,Y)
=\operatorname{Hom}_{X_\lambda}(X_\lambda,Z_\lambda)
\quad\hbox{et}\quad
\operatorname{Hom}_S(X,Y)=\operatorname{Hom}_X(X,Z)
\]
(I, 3.3.14). Il suffit maintenant d'appliquer (8.8.2, (i)) en y faisant
$X_\lambda=S_\lambda$ et en y remplaçant $Y_\lambda$ par $Z_\lambda$.

\begin{corollary}[8.13.2]
\label{IV.8.13.2-fr}
Avec les notations de (8.13.1), on suppose $X_\alpha$ quasi-compact et
quasi-séparé, et les $v_{\alpha\lambda}$ affines pour
$\alpha\leqslant\lambda$~; on suppose en outre que
$Y=\varprojlim_\rho Y_\rho$, où $(Y_\rho,t_{\rho\sigma})$ est un système
projectif filtrant de $S$-préschémas tel que, pour chaque $\rho$, le morphisme
structural $Y_\rho\to S$ soit localement de présentation finie. On a alors une
bijection canonique
\begin{equation}
\label{IV.8.13.2.1-fr}
\operatorname{Hom}_S(X,Y)
\xrightarrow{\sim}
\varprojlim_\rho\left(\varinjlim_\lambda
\operatorname{Hom}_S(X_\lambda,Y_\rho)\right). \tag{8.13.2.1}
\end{equation}
\end{corollary}

En effet, le fait que $Y$ soit limite projective des $Y_\rho$ entraîne en
particulier que l'application canonique
\[
\operatorname{Hom}_S(X,Y)\longrightarrow
\varprojlim_\rho\operatorname{Hom}_S(X,Y_\rho)
\]
est bijective~; et d'autre part, les hypothèses entraînent, pour chaque
$\rho$, l'existence d'une bijection canonique
\[
\operatorname{Hom}_S(X,Y_\rho)\xrightarrow{\sim}
\varinjlim_\lambda\operatorname{Hom}_S(S_\lambda,Y_\rho)
\]
en vertu de (8.13.1)~; d'où la conclusion.

\begin{env}[8.13.3]
\label{IV.8.13.3-fr}
Les résultats précédents permettent d'interpréter dans la théorie des
préschémas les notions de «~pro-variété~» ou de «~pro-schéma~» qui
s'introduisent dans certaines applications (par exemple dans la théorie du
corps de classes local suivant les idées de Serre [39] ou dans la théorie de
Néron sur la réduction des variétés
\oldpage[IV-3]{50}
abéliennes [32]). Rappelons rapidement ici la notion de \emph{pro-objet} d'une
catégorie, renvoyant au chap.~V pour de plus amples développements (nous
n'utiliserons d'ailleurs pas avant le chap.~V l'interprétation qui va suivre,
et le lecteur peut donc omettre jusque-là la lecture de la fin de ce numéro).
Étant donnée une catégorie $\mathcal C$, la catégorie
$\mathbf{Pro}(\mathcal C)$ des \emph{pro-objets} de $\mathcal C$ a pour objets
les \emph{systèmes projectifs} (dans l'univers où l'on se place)
$\mathbf X=(X_\mu)_{\mu\in M}$ d'objets de $\mathcal C$ dont les ensembles
d'indices (dépendant du système projectif considéré) sont supposés préordonnés
filtrants. Étant donnés deux tels pro-objets
$\mathbf X=(X_\mu)_{\mu\in M}$,
$\mathbf X'=(X'_{\mu'})_{\mu'\in M'}$, les \emph{morphismes} de $\mathbf X$
dans $\mathbf X'$ sont par définition les éléments de l'ensemble
\[
\varprojlim_{\mu'}\left(\varinjlim_\mu
\operatorname{Hom}(X_\mu,X'_{\mu'})\right)~;
\]
la vérification du fait que l'on peut prendre ces ensembles pour ensembles de
morphismes est immédiate, la composition de systèmes de morphismes
$u_{\mu'}^\mu:X_\mu\to X'_{\mu'}$,
$u_{\mu''}^{\prime\mu'}:X'_{\mu'}\to X''_{\mu''}$ qui sont inductifs en
l'indice supérieur et projectifs en l'indice inférieur, se faisant
«~argument par argument~», autrement dit en considérant le système des
$u_{\mu''}^{\prime\mu'}\circ u_{\mu'}^\mu$.
\end{env}

\begin{env}[8.13.4]
\label{IV.8.13.4-fr}
Considérons alors un préschéma quasi-compact et quasi-séparé $S$, et désignons
par $\mathcal C$ la sous-catégorie pleine de la catégorie
$(\mathbf{Sch})_{/S}$ des $S$-préschémas, formée des $S$-préschémas $X$ ayant
la propriété suivante~: le morphisme structural $X\to S$ se factorise en
\[
X\xrightarrow{g}X_0\xrightarrow{f}S,
\]
où $g:X\to X_0$ est affine et $f:X_0\to S$ de présentation finie~; nous
dirons pour abréger que les préschémas de $\mathcal C$ sont
\emph{essentiellement affines au-dessus de $S$}.

Considérons d'autre part la sous-catégorie pleine $\mathcal C'_0$ de
$(\mathbf{Sch})_{/S}$ formée des $S$-préschémas \emph{de présentation finie},
et la catégorie $\mathbf{Pro}(\mathcal C'_0)$ des pro-objets de
$\mathcal C'_0$. Nous dirons qu'un objet
$\mathbf X=(X_\mu)_{\mu\in M}$ de $\mathbf{Pro}(\mathcal C'_0)$ est
\emph{essentiellement affine} s'il existe un $\gamma\in M$ tel que pour tout
$\mu\geqslant\gamma$, le morphisme de transition
$v_{\gamma\mu}:X_\mu\to X_\gamma$ soit affine (ce qui entraîne que pour
$\gamma\leqslant\mu\leqslant\nu$, $v_{\mu\nu}:X_\nu\to X_\mu$ est affine).
On notera qu'un objet de $\mathbf{Pro}(\mathcal C'_0)$ \emph{isomorphe} à un
objet essentiellement affine \emph{n'est pas} nécessairement essentiellement
affine lui-même.

Nous désignerons par $\mathcal C'$ la sous-catégorie pleine de
$\mathbf{Pro}(\mathcal C'_0)$ formée des pro-objets de $\mathcal C'_0$
essentiellement affines.

Cela étant, il résulte de (8.2.2) et (8.2.3) que pour tout objet
$\mathbf X=(X_\mu)_{\mu\in M}$ de $\mathcal C'$, le $S$-préschéma
$X=\varprojlim_\mu X_\mu$ existe~; en outre, comme, pour $\mu$ assez grand,
le morphisme $X\to X_\mu$ est affine (8.2.2), $X$ est
\emph{essentiellement affine au-dessus de $S$} par définition. Posons
$X=L(\mathbf X)$~; montrons qu'on a ainsi défini un \emph{foncteur canonique}
\begin{equation}
\label{IV.8.13.4.1-fr}
L:\mathcal C'\longrightarrow\mathcal C. \tag{8.13.4.1}
\end{equation}

On a en effet, pour deux objets $\mathbf X=(X_\mu)$,
$\mathbf X'=(X'_{\mu'})$ de $\mathcal C'$, une application canonique pour
chaque $\mu'$
\[
\varinjlim_\mu\operatorname{Hom}_S(X_\mu,X'_{\mu'})
\longrightarrow
\operatorname{Hom}_S(\varprojlim_\mu X_\mu,X'_{\mu'})
\]
définie dans (8.13.1.1), et d'autre part, par définition de la limite
projective, une bijection canonique
\[
\varprojlim_{\mu'}
\operatorname{Hom}_S(\varprojlim_\mu X_\mu,X'_{\mu'})
\xrightarrow{\sim}
\operatorname{Hom}_S(\varprojlim_\mu X_\mu,
\varprojlim_{\mu'}X'_{\mu'}),
\]
\oldpage[IV-3]{51}
d'où une application canonique
\begin{equation}
\label{IV.8.13.4.2-fr}
\varprojlim_{\mu'}\left(\varinjlim_\mu
\operatorname{Hom}_S(X_\mu,X'_{\mu'})\right)
\longrightarrow
\operatorname{Hom}_S(\varprojlim_\mu X_\mu,
\varprojlim_{\mu'}X'_{\mu'}), \tag{8.13.4.2}
\end{equation}
évidemment fonctorielle en $\mathbf X$ et $\mathbf X'$, et qui complète la
définition du foncteur $L$.
\end{env}

\begin{proposition}[8.13.5]
\label{IV.8.13.5-fr}
Les hypothèses et notations étant celles de (8.13.4), le foncteur $L$ est
pleinement fidèle. Si de plus $S$ est un préschéma noethérien (ce qui implique
déjà que $S$ est quasi-compact et quasi-séparé (1.2.8)), $L$ est une
\emph{équivalence de catégories}.
\end{proposition}

Dire que $L$ est pleinement fidèle signifie que l'application (8.13.4.2) est
bijective quels que soient $\mathbf X$, $\mathbf X'$ dans $\mathcal C'$, ce
qui est un cas particulier de (8.13.2)~: en effet, les morphismes structuraux
$X_\mu\to S$ étant de présentation finie, sont en particulier quasi-compacts
et quasi-séparés, donc les $X_\mu$ sont quasi-compacts et quasi-séparés.

Pour montrer que lorsque $S$ est noethérien $L$ est une équivalence de
catégories, il suffit, puisque l'on sait déjà que $L$ est pleinement fidèle,
de prouver que tout $S$-préschéma essentiellement affine $X$ est
$S$-\emph{isomorphe} à un objet de la forme $L(\mathbf X)$ où
$\mathbf X\in\mathcal C'$ $(\mathbf 0_{\mathrm{III}},8.1.5)$.

Or, il y a par hypothèse une factorisation
$X\xrightarrow{g}X_0\xrightarrow{f}S$ du morphisme structural, $f$ étant de
présentation finie et $g$ affine. On peut donc écrire
$X=\operatorname{Spec}(\mathcal A)$, où $\mathcal A$ est une
$\mathcal O_{X_0}$-Algèbre quasi-cohérente (II, 1.3.1). Or, comme $X_0$ est
noethérien (puisqu'il en est ainsi de $S$ et que $f$ est de type fini),
$\mathcal A$ est limite inductive filtrante de la famille
$(\mathcal A_\lambda)$ de ses sous-$\mathcal O_{X_0}$-Algèbres quasi-cohérentes
de type fini (I, 9.6.6). Posons
$X_\lambda=\operatorname{Spec}(\mathcal A_\lambda)$~; les morphismes
$X_\lambda\to X_0$ sont de type fini, donc de présentation finie puisque
$X_0$ est noethérien, et par suite il en est de même des morphismes composés
$X_\lambda\to X_0\to S$ (1.6.2)~; autrement dit, les $X_\lambda$
appartiennent à $\mathcal C'_0$ et comme les morphismes $X_\lambda\to X_0$
sont affines, $\mathbf X=(X_\lambda)$ est un objet de $\mathcal C'$ dont la
limite projective existe et est $S$-isomorphe à $X$ en vertu de (8.2.2). Ceci
achève la démonstration.

\begin{remark}[8.13.6]
\label{IV.8.13.6-fr}
Il résulte de (1.6.2) et de (II, 1.6.2) que si $X$ et $Y$ sont
essentiellement affines au-dessus de $S$, alors il en est de même de
$X\times_S Y$. On en conclut par exemple $(\mathbf 0_{\mathrm{III}},8.2.5)$
qu'un $\mathcal C$-groupe n'est autre qu'un $(\mathbf{Sch})_{/S}$-groupe qui
est un préschéma essentiellement affine au-dessus de $S$. D'autre part, les
produits finis existent dans la catégorie $\mathcal C'$~: en effet, si
$\mathbf X=(X_\mu)_{\mu\in M}$, $\mathbf Y=(Y_\rho)_{\rho\in R}$ sont deux
objets de $\mathcal C'$, les produits $X_\mu\times_S Y_\rho$ sont des
$S$-préschémas de présentation finie, et en prenant pour morphismes de
transition
\[
X_\nu\times_S Y_\sigma\longrightarrow X_\mu\times_S Y_\rho
\]
les produits des morphismes de transition $X_\nu\to X_\mu$ et
$Y_\sigma\to Y_\rho$, on voit aussitôt que
$(X_\mu\times_S Y_\rho)$ est produit de $\mathbf X$ et de $\mathbf Y$ dans
$\mathbf{Pro}(\mathcal C'_0)$~; en outre (II, 1.6.2) les morphismes de
transition ainsi définis sont affines pour $\mu$ et $\rho$ assez grands, donc
le produit $\mathbf X\times\mathbf Y$ ainsi défini appartient bien à
$\mathcal C'$. On conclut alors comme plus haut qu'un $\mathcal C'$-groupe est
un $\mathbf{Pro}(\mathcal C'_0)$-groupe qui est essentiellement affine. On
déduit donc de (8.13.5) que les catégories des $\mathcal C'$-groupes et des
$\mathcal C$-groupes sont \emph{équivalentes} lorsque $S$ est
\emph{noethérien}. Il semble plausible que lorsque $S$ est le spectre d'un
corps $k$, la catégorie des $\mathcal C$-groupes soit \emph{équivalente} à
celle des $\boldsymbol K$-groupes, où $\boldsymbol K$ est la catégorie des
$S$-préschémas \emph{quasi-compacts}~; autrement dit, tout préschéma en
groupes sur $k$ qui est quasi-compact serait essentiellement affine. D'autre
part, si l'on désigne
\oldpage[IV-3]{52}
par $\mathcal C'_0$-$\boldsymbol{gr}$ la catégorie des
$\mathcal C'_0$-groupes, il est vraisemblable que la catégorie des
$\mathcal C'$-groupes est équivalente à la sous-catégorie pleine de
$\mathbf{Pro}(\mathcal C'_0\text{-}\boldsymbol{gr})$ formée des «~groupes
proalgébriques essentiellement affines~», c'est-à-dire des systèmes
projectifs $\mathbf G=(G_\mu)_{\mu\in M}$, où les $G_\mu$ sont des groupes
algébriques sur $k$ et les morphismes de transition $G_\nu\to G_\mu$ sont
affines pour $\mu$ assez grand (ce que l'on peut exprimer aussi en disant que
$\mathbf G$ est extension d'un groupe algébrique par un pro-groupe algébrique
affine). La conjonction de ces deux conjectures équivaut d'ailleurs à la
suivante~: tout préschéma en groupes quasi-compacts sur $k$ est «~extension~»
d'un «~groupe algébrique~» (i.e. un préschéma en groupes de type fini sur $k$)
par un préschéma en groupes \emph{affine} sur $k$.

Les seuls pro-groupes algébriques rencontrés en pratique jusqu'à présent étant
en fait essentiellement affines, il y aura donc sans doute avantage à
substituer à l'étude des groupes pro-algébriques généraux (introduits et
étudiés par Serre [40]) celle des schémas en groupes quasi-compacts sur $k$,
dont la définition est conceptuellement plus simple.
\end{remark}

\subsection{Caractérisation d'un préschéma localement de présentation finie sur un autre, en termes du foncteur qu'il représente}
\label{subsection:IV.8.14-fr}

\begin{env}[8.14.1]
\label{IV.8.14.1-fr}
Étant donné un préschéma $S$, nous dirons encore, comme en (8.13.4), qu'un
système projectif filtrant $(X_\lambda,v_{\lambda\mu})$ de $S$-préschémas est
\emph{essentiellement affine} s'il existe $\alpha$ tel que
$v_{\alpha\lambda}$ soit un morphisme affine pour $\lambda\geqslant\alpha$.

L'énoncé suivant, qui sera surtout utile au chap.~V, précise (8.8.2, (i)) en
lui fournissant une réciproque~:
\end{env}

\begin{proposition}[8.14.2]
\label{IV.8.14.2-fr}
Soient $S$ un préschéma, $f:X\to S$ un morphisme. Pour tout $S$-préschéma
$T$, posons
\[
h_X(T)=\operatorname{Hom}_S(T,X)
\]
de sorte que $h_X$ est un foncteur contravariant de la catégorie
$(\mathbf{Sch})_{/S}$ des $S$-préschémas dans la catégorie
$\mathbf{Ens}$ des ensembles $(\mathbf 0_{\mathrm{III}},8.1.1)$, et $X$ un
objet représentant ce foncteur $(\mathbf 0_{\mathrm{III}},8.1.8)$. Les
conditions suivantes sont équivalentes~:
\begin{enumerate}[label=\emph{\alph*)}]
  \item $f$ est localement de présentation finie.
  \item Pour tout système projectif filtrant $(Z_\lambda)$ de
  $S$-préschémas, essentiellement affine (8.13.4) et formé de préschémas
  quasi-compacts et quasi-séparés, l'application canonique (8.13.1.1)
\begin{equation}
\label{IV.8.14.2.1-fr}
\varinjlim_\lambda h_X(Z_\lambda)
\longrightarrow h_X(\varprojlim_\lambda Z_\lambda) \tag{8.14.2.1}
\end{equation}
  est bijective.
  \item Pour tout système projectif filtrant $(Z_\lambda)$ de
  $S$-préschémas, tel que les $Z_\lambda$ soient des schémas affines,
  l'application (8.14.2.1) est bijective.
  \item[\rm c$'$)] Pour tout ouvert affine $U$ de $S$ et tout système
  projectif filtrant $(Z_\lambda)$ de $U$-préschémas, tel que les $Z_\lambda$
  soient des schémas affines, l'application (8.14.2.1) est bijective.
\end{enumerate}
\end{proposition}

Le fait que \emph{a)} implique \emph{b)} n'est autre que (8.13.1)~; il est
trivial que \emph{b)} implique \emph{c)} et que \emph{c)} implique
\emph{c$'$)}. Reste à voir que \emph{c$'$)} entraîne \emph{a)}, et comme la
propriété \emph{a)} est locale sur $S$, on peut se borner au cas où $S$ est
\emph{affine}.

\oldpage[IV-3]{53}
Supposons d'abord que $X$ soit aussi affine~; l'assertion à prouver est alors
équivalente au

\begin{corollary}[8.14.2.2]
\label{IV.8.14.2.2-fr}
Soient $A$ un anneau, $B$ une $A$-algèbre. Afin que, pour tout système
inductif filtrant $(C_\lambda)$ de $A$-algèbres, l'application canonique
\begin{equation}
\label{IV.8.14.2.3-fr}
\varinjlim_\lambda\operatorname{Hom}_{A\text{-}\mathrm{alg.}}(B,C_\lambda)
\longrightarrow
\operatorname{Hom}_{A\text{-}\mathrm{alg.}}
(B,\varinjlim_\lambda C_\lambda) \tag{8.14.2.3}
\end{equation}
soit bijective, il faut et il suffit que $B$ soit une $A$-algèbre de
présentation finie.
\end{corollary}

Il reste seulement à montrer que la condition est nécessaire. Prenons d'abord
pour $(C_\lambda)$ le système inductif filtrant des sous-$A$-algèbres de type
fini de $B$, de sorte que $\varinjlim_\lambda C_\lambda=B$. Le fait que (8.14.2.3)
soit bijective entraîne en particulier que l'application identique
$1_B$ se factorise en $B\to C_\lambda\to B$ pour un $\lambda$ convenable,
ce qui entraîne $C_\lambda=B$, donc $B$ est une $A$-algèbre \emph{de type
fini}. Posons alors $B=C/\mathfrak J$, où
$C=A[T_1,\ldots,T_n]$ et $\mathfrak J$ est un idéal de $C$. Alors
$\mathfrak J$ est limite inductive filtrante des idéaux de type fini
$\mathfrak J_\lambda\subset\mathfrak J$ de $C$~; posant
$C_\lambda=C/\mathfrak J_\lambda$, et utilisant l'exactitude du foncteur
$\varinjlim$, on voit que $B$ est encore isomorphe à la limite inductive du
système inductif filtrant $(C_\lambda)$. Il existe donc un $\lambda$ et un
$A$-homomorphisme $u:B\to C_\lambda$ tels que le composé
\[
B\xrightarrow{u}C_\lambda\xrightarrow{p_\lambda}B
\]
(où $p_\lambda$ est l'homomorphisme canonique) soit l'identité. Soit
$q_\lambda:C\to C_\lambda$ l'homomorphisme canonique, et posons
$t_i=p_\lambda(q_\lambda(T_i))$~; on a donc
$p_\lambda(u(t_i))=p_\lambda(q_\lambda(T_i))$, autrement dit
\[
u(t_i)-q_\lambda(T_i)\in\mathfrak J/\mathfrak J_\lambda.
\]
Il existe donc un $\mu\geqslant\lambda$ tel que les $n$ éléments
$u(t_i)-q_\lambda(T_i)$ appartiennent à
$\mathfrak J_\mu/\mathfrak J_\lambda$ $(1\leqslant i\leqslant n)$~; si
$p_{\mu\lambda}:C_\lambda\to C_\mu$ est l'homomorphisme canonique, on a par
suite
\[
p_{\mu\lambda}(u(t_i))
=p_{\mu\lambda}(q_\lambda(T_i))=q_\mu(T_i).
\]
Remplaçant $\lambda$ par $\mu$ et $u$ par $p_{\mu\lambda}\circ u$, on peut
donc supposer que $u(t_i)=q_\lambda(T_i)$ pour tout $i$, et si
$r=p_\lambda\circ q_\lambda$ est l'homomorphisme canonique
$C\to C/\mathfrak J=B$, on peut donc écrire
$u(r(T_i))=q_\lambda(T_i)$ pour tout $i$, d'où
$q_\lambda=u\circ r$. Mais cela entraîne nécessairement que
$\mathfrak J=\mathfrak J_\lambda$, car si $z\in\mathfrak J$, on a
$r(z)=0$~; donc on a $B=C_\lambda$.

Passons maintenant au cas où $S$ est affine et $X$ quelconque~; tout revient à
prouver qu'un ouvert affine $V$ de $X$ est de présentation finie sur $S$, et
en vertu de ce qui vient d'être démontré, il suffit de prouver que pour tout
système projectif filtrant $(Z_\lambda)$ de $S$-préschémas affines,
l'application
\begin{equation}
\label{IV.8.14.2.4-fr}
\varinjlim_\lambda\operatorname{Hom}_S(Z_\lambda,V)
\longrightarrow
\operatorname{Hom}_S(\varprojlim_\lambda Z_\lambda,V) \tag{8.14.2.4}
\end{equation}
est bijective. Il est immédiat que cette application est injective, car si
$(v_\lambda)$, $(v'_\lambda)$ sont deux systèmes inductifs de
$S$-homomorphismes $v_\lambda:Z_\lambda\to V$,
$v'_\lambda:Z_\lambda\to V$ tels que les morphismes correspondants
\[
Z\xrightarrow{u_\lambda}Z_\lambda\xrightarrow{v_\lambda}V,
\qquad
Z\xrightarrow{u_\lambda}Z_\lambda\xrightarrow{v'_\lambda}V
\]
soient égaux ($u_\lambda$ étant le morphisme canonique), alors les morphismes
\[
Z\xrightarrow{u_\lambda}Z_\lambda\xrightarrow{v_\lambda}V
\xrightarrow{j}X,
\qquad
Z\xrightarrow{u_\lambda}Z_\lambda\xrightarrow{v'_\lambda}V
\xrightarrow{j}X
\]
(où $j$ est l'injection canonique) sont égaux, ce qui entraîne
$j\circ v_\lambda=j\circ v'_\lambda$ par hypothèse pour un $\lambda$
convenable, donc $v_\lambda=v'_\lambda$.

\oldpage[IV-3]{54}
Reste à prouver que (8.14.2.4) est surjective. Soit donc $v:Z\to V$ un
$S$-morphisme~; par hypothèse il existe un $\lambda$ et un $S$-morphisme
$w_\lambda:Z_\lambda\to X$ tels que $j\circ v$ se factorise en
\[
Z\xrightarrow{u_\lambda}Z_\lambda\xrightarrow{w_\lambda}X,
\]
et tout revient à prouver qu'il existe $\mu\geqslant\lambda$ tel que le
morphisme
\[
Z_\mu\xrightarrow{u_{\lambda\mu}}Z_\lambda
\xrightarrow{w_\lambda}X
\]
(où $u_{\lambda\mu}$ est le morphisme de transition) se factorise en
\[
Z_\mu\xrightarrow{v_\mu}V\xrightarrow{j}X.
\]
Posons, pour tout $\mu\geqslant\lambda$,
$U_\mu=u_{\lambda\mu}^{-1}(w_\lambda^{-1}(V))$. On a
\[
u_\mu^{-1}(U_\mu)=u_\lambda^{-1}(U_\lambda)
=u_\lambda^{-1}(w_\lambda^{-1}(V))=v^{-1}(V)=Z.
\]
Comme les $Z_\lambda$ sont quasi-compacts et que les $U_\mu$, étant ouverts,
sont ind-constructibles (1.9.6), on déduit de (8.3.4) qu'il existe
$\mu\geqslant\lambda$ tel que $U_\mu=Z_\mu$. C.Q.F.D.

\begin{remark}[8.14.3]
\label{IV.8.14.3-fr}
Le fait que l'application (8.14.2.1) soit injective lorsque $f$ est localement
de type fini (8.8.2, (i)) amène naturellement à se demander si ce résultat
admet aussi une réciproque. Il n'en est rien, même lorsque $S$ et $X$ sont
affines, car il existe des monomorphismes $X\to S$ qui ne sont pas de type fini
(I, 2.4.2), et qui mettent donc en défaut cette conjecture.
\end{remark}
